%========================================
%. CONFIGURATION ARTICLE  COURS  SEC 101
%========================================
% 				CONTRIBUTIONS
%========================================

%************************************************************
% Chargement des variables du modèle
%************************************************************

%-------------------------------------------------------------
%               FR CYBERDEF SECOPS COURSE
%                          CONFIG NOTES File
%          SET BEAMER or Article/book Mode
%                    PRZMODE OR DOCMODE
%                             2020 eduf@ction
%-------------------------------------------------------------

% insert InterTitlle for each file compose main "Notes"
%\newcommand{\INTERTITLE}{OK}

 %------------------------------------------
 % Use global project CYBERDEF 
 %            Configuration file
 %------------------------------------------
%-----------------------------------------------
%               FR CYBERDEF SECOPS COURSE
%                         GLOBAL CONFIG  File
%                             2020 eduf@ction
%-----------------------------------------------

% just used for quick access to samples on project editing tool
%%-------------------------------------------------------------
%               FR CYBERDEF SECOPS COURSE
%              SAMPLES for edx & u  commands
%                             2020 eduf@ction
%-------------------------------------------------------------


% 												I M A G E S   
%-------------------------------------------------------------------------
%\upicture{../Latex/Sources/EDU/Pictures/img-liovar-incidents}{Incidents}{01}{lb:mtbf}
%-------------------------------------------------------------------------


%\mode<all>{\picframe{../Latex/Sources/EDU/Pictures/img-cycle}{Cycle de vie de gouvernance Cyberdef}{0.3}{lbl:cycle}}


% Planche PRZ 
% begin frame \mode<all>{\texframe{Déroulement}{présentation Gestion de la menace}
%{

%}} % end frame



%\begin{table}[!hbtp]
%\resizebox{\textwidth}{!}{%
%\begin{tabular}{ l|l l}
%%\begin{tabular}{l*{3}{S[table-format=2.1]}}
%\toprule
%\textbf{Package\_name }                 & \textbf{C\&C }                                        & \textbf{sha256}                                                            \\ \hline
%caracal.raceinspace.astronaut & http://api.lulquid.xyz                       & f1d32c17a169574369088...\\ \hline
%
%\bottomrule
%\end{tabular}
%}
%\caption{Exemple de  données de TI (Threat Intelligence)}
%\label{tab:ExempleTI}
%\end{table}

%----------------------------
% Path
\def\upath{../Latex}

%-------------------------------------------------------------
%			 FR CYBERDEF SECOPS COURSE
% 					INIT Commands File
% 		Select BEAMER or Article/book Mode
% 			with sub extension .doc or .prz
%                           2020 eduf@ction
%-------------------------------------------------------------
\RequirePackage{xstring}
\RequirePackage[realmainfile]{currfile}

\newcommand{\udocname}{\jobname}

\IfSubStr*{\currfilebase}{.prz}{\newcommand{\PRZMODE}{YES}}{\newcommand{\DOCMODE}{YES}}
   
\IfSubStr*{\currfilebase}{.edt}{\newcommand{\EDTMODE}{YES}}{\newcommand{\NOMODE}{YES}}

\input{../Latex/Templates/Magic/commons.inc/edx.library.tex}

\def\utemplatepath{../Latex/Templates/Magic}

\def\edxlibspath{\utemplatepath/edxlibs}
%\input{\edxlibspath/load.tex} à mettre dans les class article, book ...

\def\uauthorphoto{\utemplatepath/commons.inc/CommonsPictures/missionauthorphoto.png}


\def\utemplatetheme{cnam}
\def\ubeamermodel{edx}
\def\uproject{CNAMSEC101}


\newcommand{\upictureext}{pdf}

 %----------------------------
 % Basic Var
 
 %\setcounter{cntx}{\year+1}
\newcommand{\INFODistrib}{Notes de cours SECOPS 2022-2023}
\newcommand{\uJournalInfo}{Orange Cyber et Cloud avec CNAM Bretagne, Cybersécurité SEC101, edufaction}

\newcommand{\udescription}{La sécurité opérationnelle au coeur de la cyberdéfense d'entreprise}

\newcommand{\uUniversity}{Conservatoire National des Arts et Métiers}
\newcommand{\uCompany}{Orange Campus Cyber}
 %\newcommand{\uJournalInfo}{CNAM Bretagne, Cybersécurité SEC101, eduf@ction}
\newcommand{\uinstituteshort}{CNAM}
\newcommand{\uinstitute}{\uUniversity}
\newcommand{\uchaire}{Chaire de Cybersécurité}
\newcommand{\uversion}{Cnam Bretagne}





%----------------------------
% All other variables are defined in
 % file : *.lesson.tex
\newcommand{\ushorttitle}{CYBERDEF SEC101}
\newcommand{\uCoursetittle}{Eléments de Sécurité Opérationnelle}
\newcommand{\uMinilogo}{Cyberdéfense d'entreprise}
\newcommand{\uCourseLongName} {Eléments de sécurité opérationnelle en cyberdéfense d'entreprise}
%\newcommand{\utitle}{Cyberdéfense d'entreprise}

%----------------------------
% Book versus Course Notes
%----------------------------
% Use "e document' "e cours"
\newcommand{\edoc}{e document\xspace}
\newcommand{\ecours}{e cours\xspace} 
\newcommand{\etitle}{éléments de sécurité opérationnelle en cyberdéfense d'entreprise} 
\newcommand{\fichetech}{une fiche TECHNO}


%%----------------------------
%% Spécifique BOOK
% %----------------------------
\newcommand{\printer}{Publication limitée pour le \uUniversity \xspace et Orange Campus Cyber}

%%----------------------------
% Attention : Use hard coded Pictures Path
% (PicframeS)
\newcommand{\picpath}{../Latex/Sources/EDU/SRC1/Pictures}
%-------------------------------------------------------------

%%----------------------------
%% Spécifique COURSE
% %----------------------------

\newcommand {\parttitlecourse}{TITRE PARTIE (config)}
\newcommand{\coursenumber}{1}

\newcommand{\uGlossaryAcronyms}{
\loadglsentries{../Latex/Sources/commons/acronymes.tex}
\loadglsentries{../Latex/Sources/commons/glossaire.tex}
}

\newcommand{\ucourseresume}{\normalsize Aborder la sécurité des systèmes d’information sous l’angle d'une sécurité dynamique est un axe qui depuis quelques années apporte de nouvelle manière d’aborder la protection, la défense, et la résilience des systèmes d’information. La transformation digitale de l'entreprise modifie et rend plus flous les périmètres des systèmes d’informations. Cela nécessite une approche élargie du risque numérique et des nouvelles architectures de systèmes d'information intégrant des technologies de cloud et de cyberdéfense. Malgré la mise en place de mesures et de technologies de protection de plus en plus élaborées, l'impact d'une attaque ayant franchi ces frontières poreuses à considérablement augmenté.  Cette compilation des notes de cours élaborée dans le cadre d'un cours d'introduction à la cybersécurité et la cyberdéfense d'entreprise aborde les grands éléments fondamentaux permettant d'appréhender le domaine. Il permet d'en comprendre les enjeux, les codes, et acteurs. Protéger l'ensemble de l'entreprise alors qu'il est complexe de définir ses frontières est illusoire. Identifier les actifs essentiels ou vitaux et mettre en place les moyens adaptés à leur protection et leur défense est une démarche tactique qui permet de graduellement réduire ses cyber-risques tout en assurant un pilotage globale de la gouvernance de la sécurité sur ses volets de conformité.}

 %------------------------------------------
 
\renewcommand{\uMinilogo}{\includegraphics[width=0.04\paperwidth]{\utemplatepath/commons.inc/CommonsPictures/\utemplatetheme.theme/shield-20.pdf}}
 
\newcommand {\Ucontribute}
{
%%-------------------------------------------------------------
%               FR CYBERDEF SECOPS COURSE
%                            Contributions Rules
%                              2020 eduf@ction
%-------------------------------------------------------------

% \section{Contributions}

% \subsection{Comment contribuer}

Les notes et les présentations sont réalisées sous \hologo{LaTeX}. 

Vous pouvez contribuer au projet du cours CYBERDEF101. Les contributions peuvent se faire sous deux formes :

\begin{itemize}
  \item Corriger, amender, améliorer les notes publiées. A chaque session des modifications et évolutions sont apportées pour tenir compte des corrections de fond et de formes. 
  \item Ajouter, compléter, modifier des parties de notes sur la base de votre lecture du cours et de vos expertises dans chacun des domaines évoqués.
\end{itemize}

 Les fichiers sources sont publiés sur GITHUB dans l'espace : \ulink{https://github.com/edufaction/CYBERDEF101}{(edufaction/CYBERDEF101)}.
 
%%-------------------------------------------------------------
%               FR CYBERDEF SECOPS COURSE
%                Contributions DATAS & NAMES
%                              2020 eduf@ction
%-------------------------------------------------------------

\subsection{Les contributeurs/auteurs du cours}

\subsubsection{co-auteurs}

%-------------------------------------------------------------
%               FR CYBERDEF SECOPS COURSE
%                Contributions DATAS & NAMES
%                              2020 eduf@ction
%-------------------------------------------------------------
%           CO AUTEURS - Rédacteur de chapitre
%-------------------------------------------------------------

\coauthor{David  BATANY}{2019-2020}{Architecture et fonctionnement des Botnets}{Cnam SEC101}

\subsubsection{contributeurs}

%-------------------------------------------------------------
%               FR CYBERDEF SECOPS COURSE
%                Contributions DATAS & NAMES
%                              2020 eduf@ction
%-------------------------------------------------------------
%                 CONTRIBUTIONS AU COURS
%-------------------------------------------------------------


\coauthor{Céline JUBY}{2020}{Contributions d'amélioration et relectures}{Orange Cyberdefense}


} 

%------------------------------------------
 %     BEAMER OU ARTICLE MODEL
 %------------------------------------------
  
\ifdefined\PRZMODE 
	\providecommand{\umainload}{%-------------------------------------------------------------
%               FR CYBERDEF SECOPS COURSE
%                          BEAMER MAIN FILE
%                             2020 eduf@ction
%-------------------------------------------------------------

%\begin{frame}
%\frametitle<presentation>{TITRE DIAPO}
% frame text
%\end{frame}

%-------------------------------------------------------------
\documentclass[ignorenonframetext, 8pt, allowframebreaks]{beamer}

%-----------------------------------------
% Load Edxlibs in MAGIC path
%-----------------------------------------
% Need : \edxlibspath defined

\input{\edxlibspath/edxlib-placeholder.tex}
\input{\edxlibspath/edxlib-epigraphe.tex}
\input{\edxlibspath/edxlib-quizzmoodle.tex}


\usepackage[utf8]{inputenc}
\usepackage[T1]{fontenc}
\usepackage{csquotes}
\usepackage[french]{babel}
\usepackage{lmodern}
\usepackage{multicol}

\input{\utemplatepath/commons.inc/CommonsPictures/\utemplatetheme.theme/colors.inc} % set uMainColor

%-------------------------------------------------------------
%               FR CYBERDEF SECOPS COURSE
%                                  EDX PACKAGES
%.    Commons packages for BOOK, NOTES (.doc) 
%                            and BEAMER (.prz)
%                              2020 eduf@ction
%-------------------------------------------------------------

%\usepackage[hyphens,spaces,obeyspaces]{url}
%\usepackage[hyphens,spaces,obeyspaces]{url}



\PassOptionsToPackage{hyphens,spaces,obeyspaces}{url}

\usepackage[utf8]{inputenc}
\usepackage[french]{babel}
\usepackage[T1]{fontenc}
%\usepackage{markdown}



\usepackage{ccicons}

\usepackage{tikz} % Required for drawing custom shapes
\usetikzlibrary{tikzmark} 
\usetikzlibrary{mindmap,trees, backgrounds}


\usepackage{verbatim}
\usepackage{wrapfig}
\usepackage{xpatch}

\usepackage{hologo}

\usepackage[figure]{totalcount}

\usepackage{fontawesome}

\usepackage[skins]{tcolorbox} % don't forget SKINS option

\usepackage{listings}
\usepackage{upquote}

\lstset{
upquote=true,columns=flexible, basicstyle=\ttfamily,
language=HTML, 
frameround=tttt,
commentstyle=\color{gray},
identifierstyle=\color{blue},
keywordstyle=\color{uMainColor}\bfseries,
xleftmargin=2em,
xrightmargin=2em,
aboveskip=\topsep,
belowskip=\topsep, 
frame=single,
rulecolor=\color{uMainColor},
backgroundcolor=\color{uMainColor!5},
breaklines,
breakindent=1.5em,
showspaces=false,
showstringspaces=false,
showtabs=false,
xleftmargin=2em,
xrightmargin=2em,
aboveskip=1em,
belowskip=1em,
}



\usepackage{graphicx} % Required for including pictures
%\graphicspath{{Pictures/}} % Specifies the directory where pictures are stored

\usepackage{booktabs} % Required for nicer horizontal rules in tables

\usepackage{setspace}
\setstretch{1,2}

\usepackage{csquotes}

\usepackage{fancyhdr}

\usepackage{makeidx} 

\usepackage{xspace} 

\usepackage{floatflt} 

\usepackage{hyperref}
\usepackage{cleveref}

%pdfdisplaydoctitle={\utitle}, Attention le titre ne doit pas contenir de saut de ligne 

\hypersetup{
pdftitle={\uCoursetittle},
pdfauthor={\uauthor},
hidelinks,
backref=true,
pagebackref=true,
hyperindex=true,
pdfcenterwindow = true,
colorlinks=true,
breaklinks=true,
bookmarks=true,
allcolors=uMainColor,
pdfsubject = {\udescription},
pdfcreator = {MACTEX on TEXPAD},
bookmarksopen=true}


% urlcolor=cnam,
% runcolor=uMainColor,
%\pdfstringdefDisableCommands{%
%  \def\\{}%
%  \def\texttt#1{<#1>}%
%}

\usepackage[style=numeric,citestyle=numeric,sorting=nyt,sortcites=true,autopunct=true,hyperref=true,abbreviate=false,backref=true,backend=biber]{biblatex}

%\usepackage{cleveref} old position

\usepackage{multicol}

%--------------------------------------------------------
% pour la génération du QRCODE en page de garde d'article
\usepackage{qrcode} 

%--------------------------------------------------------
% POur la date d'impression en premiere page d'article
\usepackage[useregional]{datetime2}

%--------------------------------------------------------
% Gestion des glossaires et Sources/EDU/acronymes
% Un fichier acroymes.tex centralisé
\usepackage[acronym]{glossaries}

%\makeglossaries
\makenoidxglossaries 

\loadglsentries{Latex/Sources/commons/acronymes.tex}
\loadglsentries{Latex/Sources/commons/glossaire.tex}


%\renewcommand{\bibname}{R\'ef\'erences}

%--------------------------------------------------------
% Pour exclure des "environnement", ici techworkbox par exemple, travaux pour étudiants, non affiché dans le BOOK.
\usepackage{comment}


%--------------------------------------------------------
% pour filigramme sur les pages
\usepackage{eso-pic}

%--------------------------------------------------------
% for \gappto{\UrlBreaks}{\UrlOrds}
% rather than 
%\makeatletter
% \g@addto@macro{\UrlBreaks}{\UrlOrds}
% \makeatother

\usepackage{etoolbox}

\usepackage{breakurl}


\newcommand{\ucite}[1]{\cite{#1}}
%\newcommand{\ucite}[1]{REF BIB #1}
\newcommand{\ugls}[1]{\gls{#1}}
%\newcommand{\ugls}[1]{GLOSSAIRE (#1)}
%\usepackage[Rejne]{fncychap}

%\usepackage{pythontex} % need shell escape

%
%\makeatletter
%\DeclareRobustCommand{\afilename}[1]{%
% \begingroup
%  % \lstname seems to change hyphens into \textendash
%  \def\textendash{-}%
%  \afilename@parse{#1}%
%  \edef\afilename@base{\detokenize\expandafter{\afilename@base}}%
%  \texttt{\afilename@base.\afilename@ext}%
% \endgroup
%}
%\makeatother

%-------------------------------------------------------------
%               FR CYBERDEF SECOPS COURSE
%                                  EDX PACKAGES
%.    Commons COMMANDS for BOOK, NOTES (.doc) 
%                            and BEAMER (.prz)
%                              2020 eduf@ction
%-------------------------------------------------------------

\newcommand{\ucite}[1]{\cite{#1}}
%\newcommand{\ucite}[1]{REF BIB #1}
\newcommand{\ugls}[1]{\gls{#1}}
%\newcommand{\ugls}[1]{GLOSSAIRE (#1)}
\newcommand{\uac}[1]{\acrshort{#1}}



% \newcommand{\ExtractBeforeLastSlash}[1]{\directlua{extract_before_last_slash("#1")}}
% \newcommand{\ExtractAfterLastSlash}[1]{\directlua{extract_after_last_slash("#1")}}
% \newcommand{\RemoveExtension}[1]{\directlua{remove_extension("#1")}}
% \newcommand{\ExtractExtension}[1]{\directlua{extract_extension("#1")}}


\newcommand{\includepicturesubpath}[1]
{#1}

%{\ExtractBeforeLastSlash{#1}/\upicturethemepath\RemoveExtension{\ExtractAfterLastSlash{#1}}.\ExtractExtension{#1}}
    

%---------------------------------------------------
%										 	AUDIO NOTE 
%---------------------------------------------------

\newcommand\audionote[2]{%
  \note{#2}
  \immediate\write\audioexport{==AudioSlide=============}
  \immediate\write\audioexport{#1}
  \immediate\write\audioexport{\insertframenumber}
  \immediate\write\audioexport{---------------------------------}
  \immediate\write\audioexport{#2}
  \immediate\write\audioexport{---------------------------------}
}%

%---------------------------------------------------
%										 	TEXT STYLE SHORTCUT 
%---------------------------------------------------

\newcommand{\tcb}[1]{\textbf{\textcolor{uMainColor}{#1}}}

\newcommand{\voir}[1]{(Schéma~\ref{#1},~page~\pageref{#1}}

\newcommand{\graphicxextension}{pdf}

\newcommand\myemptypage{
   
    \thispagestyle{empty}
    \addtocounter{page}{-1}
    \newpage
  \null
    %CECI EST UNE PAGE VIDE
       \vfill
    }


\newcommand{\uNA}[2]
{
\newacronym{a#1}{#1}{#2}
}

\newcommand{\coauthor}[4]
{
\noindent (#2) \textbf{#1} - #4 : \textit{#3}
}


\newcommand\UKword[1]{\emph{#1}\xspace}
\newcommand\textUK[1]{\textit{\textbf{#1}}\xspace}
\newcommand\Ucenter[1]{\begin{center}{#1}\end{center}}

\newcommand\head[1]{\textbf{#1}}
\newcommand\tb[1]{\textbf{#1}}
\newcommand{\g}[1]{\og #1 \fg{}}
\newcommand{\uload}[1] {\input{../Latex/#1}}
\newcommand{\ucurcolor}{black}

\newcommand{\mak}[1]{\faicon{\aArrowCircleORight} #1}


%-----------------------------------------------------
%										 	\INFOGitHub 
%---------------------------------------------------


\newcommand{\GITfilename}{https://github.com/edufaction/CYBERDEF101/raw/main/Release/Builder/\jobname.pdf}

%\newcommand{\safeqrcode}[2][]{%
%\qrcode[#1]{\detokenize{#2}} }

\newcommand{\safeqrcode}[2][]{\qrcode[#1]{#2}} 

\newcommand{\INFOGitHub}
{
\begin{center} 
\setstretch{1.5}
 {Vérifiez la disponibilité d'une version plus récente de} \\
 {\ulink{\GITfilename}{\textbf{\jobname.pdf }sur GITHUB CYBERDEF} }   \\ 
{{\huge\ccbyncndeu}}  \\  
{Publication en \href{https://creativecommons.org/licenses/by-nc-nd/2.0/fr/}{Creative Common BY-NC-ND} by eduf@ction}    \\  % Copyright notice
{\safeqrcode[height=4 cm,padding]{\GITfilename}}  
\end{center} 
}

\newcommand{\INFOGitHubS}
{
\noindent{\safeqrcode[height=3 cm]{\GITfilename}} \\
\noindent {\tiny Téléchargez une version à jour} 
}



%-----------------------------------------------------
%										 	\upspicture 
%---------------------------------------------------
\newcommand{\upspicture}[3]
{
\renewcommand{\ucurcolor}{#3}
\resizebox{#2\textwidth}{!}{\input{../#1}}
}


%-----------------------------------------------------
%										 	\uindex 
%---------------------------------------------------


\newcommand{\uindex}[1]
{
#1\index{#1}\xspace
}

\newcommand{\ulindex}[1]{\index{#1}}

%-----------------------------------------------------
%										 	\includer 
%---------------------------------------------------

\newcommand{\includer}[1]
{
\input{../Latex/Sources/EDU/SRC1/Chapters/#1}
}

%-----------------------------------------------------
%										 	\edxdico 
%---------------------------------------------------

\newcommand{\edxdico}[2] {\textbf{#1}:\textit{#2} \index{#1}}

%-----------------------------------------------------
%  command DEF : InTextx
%-----------------------------------------------------
\newcommand\InTexta{no text}
\newcommand\InTextb{no text}
\newcommand\InTextc{no text}
\newcommand\rnc{\renewcommand}


%-----------------------------------------------------
%										 	ENV warningbox
%---------------------------------------------------

\newtcolorbox{warningbox}[2][]
{
  colframe = uMainColor!25,
  colback  = black!05,
  coltitle = uMainColor!20!black,
  fonttitle = \bfseries, 
  title    = #2,
  #1,
}


\newtcolorbox{techworkbox}[2][]{sharp corners,colframe = uMainColor!25,colback  = black!05,coltitle = uMainColor!20!black,fonttitle = \bfseries, title={\textcolor{uMainColor}{\faFile~\small {Sujet de mémoire SECOPS}}} : #2,#1,}


\newtcolorbox{notebox2}[2][]
{
sharp corners, 
colback = uMainColor!5!white, 
colframe = uMainColor!75!black,
fonttitle = \bfseries, 
colbacktitle= uMainColor!85!black,
title=#2,#1
}

\newtcolorbox{notebox}[2][]{colback=uMainColor!5!white,
colframe=uMainColor!75!black,fonttitle=\bfseries,
colbacktitle=uMainColor!85!black,enhanced,
attach boxed title to top right={yshift=+1mm},
title=#2,#1}


\newtcolorbox{toolsboxin}[2][]{sharp corners, colback=uMainColor!0!white,
colframe= black!25,fonttitle=\bfseries,
colbacktitle=black!60,enhanced,
attach boxed title to top left={yshift=+1mm},
title=#2,#1}


\newtcolorbox{techworkin}[2][]{sharp corners, colback=uMainColor!5!white,
colframe= uMainColor!100, fonttitle=\bfseries,
colbacktitle=black!5, enhanced,
 coltitle=black!50,
attach boxed title to top right ={yshift=+1mm},
title=#2,#1}



%-----------------------------------------------------
%										 	ENV ToolsBox
%---------------------------------------------------


\newcommand\tooleditor{No EDITOR in Datas.tex}
\newcommand\toolurl{No URL INSERTED in Datas.tex}
\newcommand\toolclass{No EDITOR in Datas.tex}
\newcommand\toolname{No NAME in Datas.tex}
\newcommand\toolanalyst{No ANALYST in Datas.tex}

\newcommand{\toolsbox}[2]
{
\input{../Latex/Sources/EDU/SRC1/SecComponents/#2.#1/datas.tex} 
\begin{toolsboxin}[sidebyside,righthand width=.7\textwidth]{#1}
\includegraphics[width=0.9\textwidth]{../Latex/Sources/EDU/SRC1/SecComponents/#2.#1/icon.png} \\
\tcblower
\input{../Latex/Sources/EDU/SRC1/SecComponents/#2.#1/description.tex}
\end{toolsboxin}
\noindent \small {\faCogs~Classe : \uppercase{\textbf{#2}}, Site de référence : \textbf{\ulink{\toolurl}{#1}}}
\newline
\small {\faGlobe~Editeur : \textbf{\tooleditor}}
\small {\faEye~Analyste : \textbf{\toolanalyst}}
}

%\renewcommand\toolsclass{SIEM}
%\renewcommand\toolsname{RSA Netwitness}
%\renewcommand\toolseditor{RSA}
%\renewcommand\toolsurl{https://www.rsa.com/fr-fr/products/threat-detection-response}

%-----------------------------------------------------
%										 	ENV  TechWork
%---------------------------------------------------


\newcommand{\techwork}[2]
{
\begin{techworkin} [righthand width=.8\textwidth]{\faCogs~\small {Fiche TECHNO} : #1}#2
\end{techworkin}
}

%-----------------------------------------------------
%										 	\uchap
%---------------------------------------------------
\ifdefined\INTERTITLE
\providecommand{\uchap}[1]
	{
	\begin{LARGE}
	\textbf{\Ucenter{#1 \\ -oOo-}}	
	\end{LARGE}
	}
\else
	\providecommand{\uchap}[1]{}
\fi

%-----------------------------------------------------
%											\uexpand
%---------------------------------------------------
\newcommand {\uexpand}[1]{
\mode<all>\input{../Latex/Sources/EDU/SRC1/Chapters/#1}
}


%-----------------------------------------------------
%											\spicture 
% direct include graphic with Scrivener
%-----------------------------------------------------

\newcommand{\spicture}[3]{
\begin{figure}[!hbtp] %h !hbtp
  \begin{center}
	 \includegraphics[width=0.8\textwidth]{#1}
  \end{center}
\caption{#3 - (slbl-#2)\label{slbl-#2}}
\end{figure}
}

%\StrBefore{#2}{.pdf}

%-----------------------------------------------------
%											\upicture 
%---------------------------------------------------

\ifdefined\PRZMODE

\providecommand{\upicture}[4]{
\framesubtitle{#2}
%\begin{figure}[!h] %h
  \begin{center}
	 \includegraphics[height=#3\textheight, width=#3\textwidth, keepaspectratio ]{\includepicturesubpath{#1.pdf}}
  \end{center}
%  \end{figure}
}

\else

\providecommand{\upicture}[4]{
\begin{figure}[!hbtp] %h !hbtp
  \begin{center}
	 \includegraphics[width=#3\textwidth]{\includepicturesubpath{#1.pdf}}
  \end{center}
\caption{#2 \label{#4}}
\end{figure}
}
	
\fi


\ifdefined\PRZMODE

\providecommand{\uwpicture}[4]{
\framesubtitle{#2}
%\begin{figure}[!h] %h
  \begin{center}
	 \includegraphics[height=0.8\textheight, width=0.8\textwidth, keepaspectratio ]{\includepicturesubpath{#1.pdf}}
  \end{center}
%  \end{figure}
}
\else
\providecommand{\uwpicture}[4]
{
\begin{wrapfigure}{O}{#3\textwidth} % inside / Outside
  \begin{center}
	 \includegraphics[width=#3\textwidth]{\includepicturesubpath{#1.pdf}}
\end{center}
\caption{#2 \label{#4}}
\end{wrapfigure}
}
\fi

%-----------------------------------------------------
%	\picframe for BEAMER (Image included in Beamer Slide)
%---------------------------------------------------

\ifdefined\PRZMODE

\providecommand{\picframe}[4]
{
\begin{frame}
\frametitle<presentation>{#2}
  \begin{center}

   \begin{tikzpicture}[remember picture, overlay]
    \node at (current page.center) {\includegraphics[width=#3\paperwidth, height=#3\paperheight,keepaspectratio]{\includepicturesubpath{#1.pdf}}};
    \end{tikzpicture}

  \end{center}
\end{frame}

}
\else
\providecommand{\picframe}[4]
{
\begin{figure}[!hbtp] % !hbtp ou !h 
\centering 
%{\includegraphics[width=#3\paperwidth, height=#3\paperheight,keepaspectratio]{\includepicturesubpath{#1.pdf}}}
\includegraphics[width=#3\textwidth]{\includepicturesubpath{#1.pdf}}
\caption{#2 \label{#4}}
\end{figure}
}
\fi

%-----------------------------------------------------
%.  for BEAMER (Image included in Beamer Slide)
%	\picframeS Simplified with Name and Label from One Param
%---------------------------------------------------


\ifdefined\PRZMODE

\providecommand{\picframeS}[2]
{
\begin{frame}
\frametitle<presentation>{#2}
  \begin{center}
	 \includegraphics[height=0.95\textheight, width=0.95\textwidth, keepaspectratio ]{\picpath/img-#1.\graphicxextension}
  \end{center}
\end{frame}

}
\else
\providecommand{\picframeS}[2]
{
\begin{figure}[!hbtp] % !hbtp ou !h 
\centering \includegraphics[width=0.9\textwidth]{\picpath/img-#1.\graphicxextension}
\caption{(s) #2 \label{lbl:#1}}
\end{figure}
}

\fi


%{figure}[!h]
%{wrapfigure}{O}{#3\textwidth}
%-----------------------------------------------------
%	FRAME			\wpicframe 
%---------------------------------------------------

\ifdefined\PRZMODE

\providecommand{\wpicframe}[4]
{
\begin{frame}
\frametitle<presentation>{#2}
  \begin{center}
	 \includegraphics[height=0.95\textheight, width=0.95\textwidth, keepaspectratio ]{#1.\graphicxextension}
  \end{center}
\end{frame}

}
\else
\providecommand{\wpicframe}[4]
{
\begin{wrapfigure}{O}{#3\textwidth}
\centering \includegraphics[width=#3\textwidth]{#1.\graphicxextension}
\caption{#2 \label{#4}}
\end{wrapfigure}
}
\fi

%-----------------------------------------------------
%	FRAME										\texframe 
%---------------------------------------------------
\ifdefined\PRZMODE

\providecommand{\texframe}[3]
{
\begin{frame}
\frametitle<presentation>{#1}
\framesubtitle<presentation>{#2}
#3
\end{frame}
}
\else
\providecommand{\texframe}[3]
{
#3
}

\fi

%-----------------------------------------------------
%											\uref 
%---------------------------------------------------
%\newcommand{\uref}[2]{(Voir~#1~\cref{#2} page~\pageref{#2})}

\newcommand{\uref}[2]{(Voir~#1~\ref{#2} page~\pageref{#2})}
%-----------------------------------------------------
%											\ubg
%---------------------------------------------------
\newcommand{\ubg}[1]{
\textbf{\g{\uppercase{#1}}}
}
%-----------------------------------------------------
%											\updfimage 
%---------------------------------------------------

\newcommand{\updfimage}[3]
{
\begin{wrapfigure}{R}{#3\textwidth}
  \begin{center}
	 \resizebox{#3\textwidth}{!}{\includegraphics{\includepicturesubpath{#1.pdf}}}
  \end{center}
\caption{#2}
\end{wrapfigure}
}


%************************PACKAGE***

%\usepackage{tikz}% Required for drawing custom shapes
\usetikzlibrary{shadows}


%-----------------------------------------------------
%											 		\rem 
% insertion d'une remarque avec R
%---------------------------------------------------
\newcommand\rem[1]{%
   \marginpar{%
   \tikzpicture[baseline={(title.base)}]
      \node[inner sep=5pt,text width=4cm,drop shadow={shadow yshift=-5pt,shadow xshift=5pt,uMainColor},fill=white] (box) {\vskip5pt \nointerlineskip #1};
      \node[right=10pt,inner sep=0pt,outer sep=10pt] (title) at (box.north west) {\bfseries\color{uMainColor}Remarque};
      \draw[draw=uMainColor,very thick](title.west)--(box.north west)--(box.south west)--(box.south east)--(box.north east)--(title.east);
      \fill[uMainColor]([yshift=-10pt]box.north west)--+(-5pt,-5pt)--+(0pt,-10pt);
   \endtikzpicture}%
}
%-----------------------------------------------------
%											 \utikzimage 
% insertion d'une image sur fichier tkz.tex
%---------------------------------------------------

\newcommand{\utikzimage}[3]
{
\begin{wrapfigure}{R}{#3\textwidth}
  \begin{center}
	 \resizebox{#3\textwidth}{!}{\input{../#1.tkz.tex}}
  \end{center}
\caption{#2}
\end{wrapfigure}
}

%-----------------------------------------------------
%														\link 
% Url et footnote du lien
%-----------------------------------------------------

% TT every character and hyphenate after it
\def\hyphenateAndTtWholeString #1{\xHyphenate#1$\wholeString\unskip}

\def\xHyphenate#1#2\wholeString {\if#1$%
    \else\transform{#1}%
    \takeTheRest#2\ofTheString\fi}

\def\takeTheRest#1\ofTheString\fi
{\fi \xHyphenate#1\wholeString}

\def\transform#1{\url{#1}\hskip 0pt plus 1pt}

% Define the \urlx command which works like \url, but with line brakes
\def\urlx #1{\href{#1}{\hyphenateAndTtWholeString{#1}}}
  

 %-----------------------------------------------------
%														\link 
%-----------------------------------------------------  
  
  \newcommand{\ulink}[2]
  {\href{#1}
  {#2~\raisebox{-0.2ex}{\faExternalLink}\footnote{\url{#1}}}\xspace}
  
  \newcommand{\wikipedia}[1]
  { \href{https://fr.wikipedia.org/wiki/#1}{\textcolor{grey}{#1}}\raisebox{0.5ex}{\tiny{\textcolor{uMainColor}{\faWikipediaW}}}}
  


 %-----------------------------------------------------
%														\utodo 
%-----------------------------------------------------

  \newcommand{\utodo}{\begin{center}{\noindent \textcolor{cyan}{\faGears~{En cours de rédaction, DRAFT non publiable} \faGears}}\end{center}}
  
    \newcommand{\utocomplete}{\begin{center}{\noindent \textcolor{blue}{\faClipboard~{à compléter, les éléments ne donnent qu'une vue trop réduite ou trop parcellaire du sujet}}}\end{center}}
  

 





%-------------------------------------------------------
\usepackage{\utemplatepath/LessonModel/Beamer/\ubeamermodel.model/edxbeamer}
%-------------------------------------------------------


% a charger après hyperref, pour avoir les guillemets sous beamer
\usepackage{ae}
\usepackage{aecompl}


\usepackage{environ}
\newwrite\audioexport



%\renewcommand{\familydefault}{pag}

%-------------------------------------------------------
% Optional: a subtitle to be displayed on the title slide

\title[\ushorttitle]{\lessonID~-~\utitle}
\subtitle{\tiny\uCourseLongName}

%-------------------------------------------------------
% The author(s) of the presentation:
%  - again first a short version to be displayed at the bottom;
%  - next the full list of authors, which may include contact information;

\author[\uauthorwriter]{ \LARGE{\uauthor} \\ \medskip {\small \url{\umaila} \and \url{\umailb}\\ {\small \url{\usitea}}}}

%-------------------------------------------------------
% The institute:
%  - to start the name of the university as displayed on the top of each slide
%    this can be adjusted such that you can also create a Dutch version
%  - next the institute information as displayed on the title slide

\institute[\uinstituteshort]{{\uinstitute \\ \uchaire}}


%-------------------------------------------------------
% Add a date and possibly the name of the event to the slides
%  - again first a short version to be shown at the bottom of each slide
%  - second the full date and event name for the title slide
\date[\uversion]{\tiny{Publication \INFODistrib ~du \\ \DTMnow}}


\setbeamertemplate{navigation symbols}{%
\insertslidenavigationsymbol
\insertframenavigationsymbol
\insertsubsectionnavigationsymbol
\insertsectionnavigationsymbol
\insertdocnavigationsymbol
\insertbackfindforwardnavigationsymbol
}

%\mode<all>{
%%-------------------------------------------------------------
%               FR CYBERDEF SECOPS COURSE
%
%                 DEFINITION DES Sources/EDU/acronymes

%                      $File :  Sources/EDU/acronymes.tex
%                             2020 eduf@ction
%-------------------------------------------------------------
% Require : glossaries package
%-------------------------------------------------------------
%\newacronym{XXX}{XXX}{XXX}
%-----------------------------------------------------------

%  Organismes

\newacronym{aCNIL}{CNIL}{Commission National Informatique et Liberté}
\newacronym{aOTAN}{OTAN}{Organisation du Traité de l'Atlantique Nord}
\newacronym{aANSSI}{ANSSI}{Agence Nationale de la Sécurité des systèmes d'information}
\newacronym{aCERT}{CERT}{Computer Emergency Response Team}
\newacronym{aAFNOR}{AFNOR}{Association Française de Normalisation}
\newacronym{aSIEM}{SIEM}{System Incident and Event Management}
\newacronym{aSIM}{SIM}{System Incident Management}
\newacronym{aSEM}{SEM}{System Event Management}
\newacronym{aIRP}{IRP}{Incident Response Plan, Plan de Réponse à Incident}
\newacronym{aISO}{ISO}{International Organization for Standardization}
\newacronym{aDICT}{DICT}{Disponibilité, Intégrité, Confidentialité, Traçabilité}
\newacronym{aPDCA}{PDCA}{Plan, Do, Check, Act (Roue de Deming)}
\newacronym{aGDPR}{GDPR}{General Data Protection Regulation ou RGPD Règlement général sur la protection des données}

%-----------------------------------------------------------

% ROLES ET METIERS

\newacronym{aOS}{OS}{Officier de sécurité}
\newacronym{aRSSI}{RSSI}{Responsable de la sécurité des systèmes d'information}
\newacronym{aDPO}{DPO}{Data Protection Officier}
\newacronym{aCIL}{CIL}{ Correspondant Informatique et Liberté}
\newacronym{aSSI}{SSI}{Sécurité des Systèmes d’information}
\newacronym{aGRC}{GRC}{Gouvernance Risques et Conformité}
\newacronym{aDTF}{DTF}{Defend The Flag}
\newacronym{aCTF}{CTF}{Capture The Flag}
\newacronym{aOIV}{OIV}{Opérateur d'Infrastructures Vitales}
\newacronym{aOSE}{OSE}{Opérateur de Services Essentiels}
\newacronym{aLPM}{LPM}{Loi de Programmation Militaire}
\newacronym{aEBIOS}{EBIOS}{Expression des Besoins et Identification des Objectifs de Sécurité}
\newacronym{aDLL}{DLL}{Dynamic Loads Library}
\newacronym{aACID}{ACID}{Authentification, Confidentialité, Intégrité, Disponibilité}

%-------------------------------------------------------
% Sources/EDU/acronymes DATABASE XLS Generation (cf. RESSOURCES DIR)
%-------------------------------------------------------
% DO NOT MODIFY DIRECTLY, USE XLS File (BEGIN)

\newacronym{a2FA}{2FA}{\textit{Two-Factor Authentication} : Authentification à deux facteurs}
\newacronym{aAAA}{AAA}{\textit{Authentication, Authorization, Accounting} : Authentification, autorisation, comptabilisation}
\newacronym{aAC}{AC}{\textit{Access Control} : Accès conditionnel}
\newacronym{aAES}{AES}{\textit{Advanced Encryption Standard} : Algorithme de chiffrement avancé}
\newacronym{aAPT}{APT}{\textit{Advanced Persistent Threat Group} : Groupe de cyber-espionnage avancé}
\newacronym{aASLR}{ASLR}{\textit{Address Space Layout Randomization} : Randomisation de l'espace d'adressage}
\newacronym{aAV}{AV}{\textit{Antivirus Software} : Logiciel antivirus}
\newacronym{aBOT}{BOT}{\textit{Robot (Malicious)} : Robot (malveillant)}
\newacronym{aCC}{C\&C}{\textit{Command And Control} : Commande et contrôle}
\newacronym{aCSRF}{CSRF}{\textit{Cross-Site Request Forgery Attack} : Attaque de confiance de site à site}
\newacronym{aDDOS}{DDOS}{\textit{Distributed Denial Of Service Attack} : Attaque de déni de service distribué}
\newacronym{aDES}{DES}{\textit{Data Encryption Standard} : Algorithme de chiffrement à clé de 64 bits}
\newacronym{aDLP}{DLP}{\textit{Data Loss Prevention} : Gestion de la protection des données sensibles}
\newacronym{aEDR}{EDR}{\textit{Endpoint Detection And Response} : Détection et réponse aux comportements anormaux d'un terminal}
\newacronym{aEK}{EK}{\textit{Exploit Kit} : Kit d'exploit}
\newacronym{aFTP}{FTP}{\textit{File Transfer Protocol} : Protocole de transfert de fichier}
\newacronym{aHTTPS}{HTTPS}{\textit{Protocol For Securing Communications On The Internet} : Protocole de sécurisation HTTP}
\newacronym{aIAM}{IAM}{\textit{Identity And Access Management} : Identité et Gestion des Accés}
\newacronym{aIDS}{IDS}{\textit{Intrusion Detection System} : Système de détection d'intrusion}
\newacronym{aIIS}{IIS}{\textit{Internet Information Services Server}}
\newacronym{aIPS}{IPS}{\textit{Intrusion Prevention System} : Système de prévention d'incident}
\newacronym{aISA}{ISA}{\textit{Internet Security And Acceleration Server} : Serveur de sécurité et d'accélération Internet (Microsoft)}
\newacronym{aLDAP}{LDAP}{\textit{Lightweight Directory Access Protocol} : Protocole de gestion d'annuaire léger}
\newacronym{aMD5}{MD5}{\textit{Message-Digest Algorithm 5}}
\newacronym{aMFA}{MFA}{\textit{Multi-Factor Authentication} : Authentification multi-facteurs}
\newacronym{aNAC}{NAC}{\textit{Network Access Control}}
\newacronym{aOAuth}{OAuth}{\textit{Open Authorization Protocol}}
\newacronym{aPAM}{PAM}{\textit{Privilege Access Management}}
\newacronym{aPKI}{PKI}{\textit{Public Key Infrastructure} : Infrastructure de clefs publiques}
\newacronym{aRSA}{RSA}{\textit{Rivest-Shamir-Adleman Encryption Algorithm}}
\newacronym{aSAML}{SAML}{\textit{Security Assertion Markup Language}}
\newacronym{aSMS}{SMS}{\textit{Short Message Service}}
\newacronym{aSQLI}{SQLI}{\textit{Sql Injection} : Injection de script de commande SQL}
\newacronym{aSSH}{SSH}{\textit{Secure Shell Protocol}}
\newacronym{aSSL}{SSL}{\textit{Protocol For Securing Communications On The Internet} : Protocole de sécurisation des communications sur Internet}
\newacronym{aTKIP}{TKIP}{\textit{Temporal Key Integrity Protocol}}
\newacronym{aTLS}{TLS}{\textit{Transport Layer Security}}
\newacronym{aURL}{URL}{\textit{Uniform Resource Locator}}
\newacronym{aVPN}{VPN}{\textit{Virtual Private Network} : Réseau privé virtuel}
\newacronym{aWAF}{WAF}{\textit{Web Application Firewall}}
\newacronym{aWPA}{WPA}{\textit{Wifi Protected Access}}
\newacronym{aXSS}{XSS}{\textit{Cross-Site Scripting Attack}}

% DO NOT MODIFY (END)
%-------------------------------------------------------




%%-------------------------------------------------------------
%          FR CYBERDEF SECOPS COURSE
%
%                        GLOSSAIRE
%
%                 $File :  Sources/EDU/glossaire.tex
%                    2020 eduf@ction
%-------------------------------------------------------------
% Require : glossaries package
%-------------------------------------------------------------

%\newglossaryentry{xxx}
%{
%name={xxx},
%description={xxx}
%}

\newglossaryentry{gLPM}
{
name={Loi de programmation militaire},
description={Une loi de programmation militaire (LPM), est une loi visant à établir une programmation pluriannuelle des dépenses et donc du budget que l'État français consacre à ses forces armées. Les lois de programmation militaire ont des durées d'application de quatre, cinq ou six ans. La cyberdéfense fait partie des axes importants des lois de programmation}
}

\newglossaryentry{gCTF}
{
name={Capture The Flag},
description={Challenge consistant à exploiter ses connaissance et les vulnérabilités affectant des logiciels  et systèmes numériques de manière à s’introduire sur des réseaux , des ordinateurs, pour récupérer (capturer) des preuves de l’intrusion réussie dénommées G\g{des drapeaux (\textit{Flags} en anglais)}.}
}

\newglossaryentry{gDTF}
{
name={Defend The Flag},
description={Challenge consistant à détecter, résister à des attaquants qui tentent de s'introduire ou sont introduits dans des systèmes numériques de manière à défendre le bien considéré comme le plus sensible convoité par l'attaquant dénommé G\g{drapeau (\textit{Flags} en anglais)}.}
}


\newglossaryentry{gROOTKIT}
{
name={Rootkit},
description={ou simplement « kit » (est un « outil de dissimulation d'activité » ou  « code furtif » qui est une sorte de trousse à outils pour un pirate informatique). Il est constitué  de techniques mises en œuvre par un ou plusieurs logiciels, dont le but est d'obtenir et de pérenniser un accès ceci avec en déployant des techniques de furtivité et de dissimulation.}
}

%-------------------------------------------------------
% GLS DATABASE XLS Generation (cf. RESSOURCES DIR)
%-------------------------------------------------------
% DO NOT MODIFY DIRECTLY, USE XLS File (BEGIN)

\newglossaryentry{gAffiliate}{name={\textit{Affiliates}},description={(Affiliés): Cybercriminels qui mènent des cyberattaques, parfois sophistiquées, par le biais de rançongiciels qui sont mis à disposition par des développeurs ou concepteurs. Les gains générés par les affiliés sont partagés avec les développeurs via un système de commissions.}}

% DO NOT MODIFY (END)
%-------------------------------------------------------


%}

%++++++++++++++++++++++++++++++++
\begin{document}



  %-------------------------------------------------------
  % AudioNote Open File
  %-------------------------------------------------------
\mode<all>{\immediate\openout\audioexport=\jobname.txt}

%================================
\begin{frame}[plain]
 					\begin{figure}
  					 \centering
  					 \begin{columns}[c]
  \begin{column}{0.4\textwidth}
  
    \begin{figure}
        \begin{flushleft}
  			\includegraphics[width=0.9\textwidth]{\utemplatepath/commons.inc/CommonsPictures/\utemplatetheme.theme/maincourselogo.pdf}
  \end{flushleft}
  \end{figure}
      \end{column}
  
  \begin{column}{0.6\textwidth}
    \begin{figure}
        \begin{flushright}
 		    \includegraphics[width=0.7\textwidth]{\utemplatepath/commons.inc/CommonsPictures/\utemplatetheme.theme/maininstitutelogo.pdf}
    \end{flushright}
  \end{figure}
  \end{column}
\end{columns}
   							  
   							  \end{figure}
							
				  				\titlepage
\end{frame}
%================================



%================================
\begin{frame}
	\frametitle{Abstract}
	\begin{techworkin} [righthand width=.8\textwidth]{\faCogs~\small {Hashtags} : \ukeywords}\uabstract
	\end{techworkin}
	\end{frame}
%================================

%================================
\begin{frame}
 \frametitle{Sommaire}
      \begin{multicols}{2}
         \tableofcontents[hideallsubsections]
     \end{multicols}
 \end{frame}
 %================================

%-------------------------------------------------------
%  BEAMER MAIN TEX configured with ubody
%-------------------------------------------------------
\mode<all>{\ubody}

%-------------------------------------------------------
% Last Frame
%-------------------------------------------------------
\begin{frame}
 \frametitle{des questions ?}
  \framesubtitle{contacter \umaila}
  					\begin{figure}
  					 \centering
   						 \includegraphics[width=0.2\textwidth]{\utemplatepath/commons.inc/CommonsPictures/\utemplatetheme.theme/mainlogo.pdf}\\
				  \end{figure}
				 	 \centering \textit{Tous les documents publiés dans le cadre de ce cours sont perfectibles, \\ne pas hésiter à m'envoyer vos remarques !}
 \end{frame}

%-------------------------------------------------------
% Contributions
%-------------------------------------------------------
\begin{frame}
 \frametitle{Contributions}
\begin{columns}
		\begin{column}{0.1\textwidth}

 					\begin{figure}
  						  \centering
   						 \includegraphics[width=0.9\textwidth]{../Latex/Sources/EDU/SRC1/Contribs/contribute/ContribPictures/github.pdf}\\
   						 	 \includegraphics[width=0.9\textwidth]{../Latex/Sources/EDU/SRC1/Contribs/contribute/ContribPictures/overleaf.pdf}
				  \end{figure}
		\end{column}
		\begin{column}{0.9\textwidth}
							%-------------------------------------------------------------
%               FR CYBERDEF SECOPS COURSE
%                            Contributions Rules
%                              2020 eduf@ction
%-------------------------------------------------------------

% \section{Contributions}

% \subsection{Comment contribuer}

Les notes et les présentations sont réalisées sous \hologo{LaTeX}. 

Vous pouvez contribuer au projet du cours CYBERDEF101. Les contributions peuvent se faire sous deux formes :

\begin{itemize}
  \item Corriger, amender, améliorer les notes publiées. A chaque session des modifications et évolutions sont apportées pour tenir compte des corrections de fond et de formes. 
  \item Ajouter, compléter, modifier des parties de notes sur la base de votre lecture du cours et de vos expertises dans chacun des domaines évoqués.
\end{itemize}

 Les fichiers sources sont publiés sur GITHUB dans l'espace : \ulink{https://github.com/edufaction/CYBERDEF101}{(edufaction/CYBERDEF101)}.
 
		\end{column}
\end{columns}
 \end{frame}
 
 %-------------------------------------------------------
% Link to download file and copyright
%-------------------------------------------------------
\begin{frame}
\frametitle{Mises à jour régulières}
\framesubtitle{Eduf@ction~\umaila}
%\INFOGitHub
{\begin{center} 
\setstretch{2}
 {Vérifiez la disponibilité d'une version plus récente de} \\
 {\ulink{\GITfilename}{\textbf{\udocname.pdf }sur GITHUB CYBERDEF} }   \\ 
{{\huge\ccbyncndeu}}  \\  
{\the\year~eduf@ction - Publication en Creative Common BY-NC-ND }    \\  % Copyright notice
{\safeqrcode[padding]{\GITfilename}}  
\end{center} }
\end{frame}


  %-------------------------------------------------------
  % AudioNote Open File
  %-------------------------------------------------------
  \mode<all>{\immediate\closeout\audioexport}


\end{document}




}
\else

\ifdefined\EDTMODE 
	\providecommand{\umainload}{%-------------------------------------------------------------
%               FR CYBERDEF SECOPS COURSE
%                          ARTICLE MAIN FILE
%                             2020 eduf@ction
%-------------------------------------------------------------

\documentclass[10pt,fleqn,twoside]{\utemplatepath/LessonModel/Article/edxstyle} 
%edxstyle
\usepackage{beamerarticle}

\usepackage{ae}
\usepackage[french]{babel}
\usepackage[utf8]{inputenc}
\usepackage[T1]{fontenc}
\usepackage{lmodern}




\usepackage{titletoc} % Required for manipulating the table of contents

\input{\utemplatepath/commons.inc/CommonsPictures/\utemplatetheme.theme/colors.inc}

%-------------------------------------------------------------
%               FR CYBERDEF SECOPS COURSE
%                                  EDX PACKAGES
%.    Commons packages for BOOK, NOTES (.doc) 
%                            and BEAMER (.prz)
%                              2020 eduf@ction
%-------------------------------------------------------------

%\usepackage[hyphens,spaces,obeyspaces]{url}
%\usepackage[hyphens,spaces,obeyspaces]{url}



\PassOptionsToPackage{hyphens,spaces,obeyspaces}{url}

\usepackage[utf8]{inputenc}
\usepackage[french]{babel}
\usepackage[T1]{fontenc}
%\usepackage{markdown}



\usepackage{ccicons}

\usepackage{tikz} % Required for drawing custom shapes
\usetikzlibrary{tikzmark} 
\usetikzlibrary{mindmap,trees, backgrounds}


\usepackage{verbatim}
\usepackage{wrapfig}
\usepackage{xpatch}

\usepackage{hologo}

\usepackage[figure]{totalcount}

\usepackage{fontawesome}

\usepackage[skins]{tcolorbox} % don't forget SKINS option

\usepackage{listings}
\usepackage{upquote}

\lstset{
upquote=true,columns=flexible, basicstyle=\ttfamily,
language=HTML, 
frameround=tttt,
commentstyle=\color{gray},
identifierstyle=\color{blue},
keywordstyle=\color{uMainColor}\bfseries,
xleftmargin=2em,
xrightmargin=2em,
aboveskip=\topsep,
belowskip=\topsep, 
frame=single,
rulecolor=\color{uMainColor},
backgroundcolor=\color{uMainColor!5},
breaklines,
breakindent=1.5em,
showspaces=false,
showstringspaces=false,
showtabs=false,
xleftmargin=2em,
xrightmargin=2em,
aboveskip=1em,
belowskip=1em,
}



\usepackage{graphicx} % Required for including pictures
%\graphicspath{{Pictures/}} % Specifies the directory where pictures are stored

\usepackage{booktabs} % Required for nicer horizontal rules in tables

\usepackage{setspace}
\setstretch{1,2}

\usepackage{csquotes}

\usepackage{fancyhdr}

\usepackage{makeidx} 

\usepackage{xspace} 

\usepackage{floatflt} 

\usepackage{hyperref}
\usepackage{cleveref}

%pdfdisplaydoctitle={\utitle}, Attention le titre ne doit pas contenir de saut de ligne 

\hypersetup{
pdftitle={\uCoursetittle},
pdfauthor={\uauthor},
hidelinks,
backref=true,
pagebackref=true,
hyperindex=true,
pdfcenterwindow = true,
colorlinks=true,
breaklinks=true,
bookmarks=true,
allcolors=uMainColor,
pdfsubject = {\udescription},
pdfcreator = {MACTEX on TEXPAD},
bookmarksopen=true}


% urlcolor=cnam,
% runcolor=uMainColor,
%\pdfstringdefDisableCommands{%
%  \def\\{}%
%  \def\texttt#1{<#1>}%
%}

\usepackage[style=numeric,citestyle=numeric,sorting=nyt,sortcites=true,autopunct=true,hyperref=true,abbreviate=false,backref=true,backend=biber]{biblatex}

%\usepackage{cleveref} old position

\usepackage{multicol}

%--------------------------------------------------------
% pour la génération du QRCODE en page de garde d'article
\usepackage{qrcode} 

%--------------------------------------------------------
% POur la date d'impression en premiere page d'article
\usepackage[useregional]{datetime2}

%--------------------------------------------------------
% Gestion des glossaires et Sources/EDU/acronymes
% Un fichier acroymes.tex centralisé
\usepackage[acronym]{glossaries}

%\makeglossaries
\makenoidxglossaries 

\loadglsentries{Latex/Sources/commons/acronymes.tex}
\loadglsentries{Latex/Sources/commons/glossaire.tex}


%\renewcommand{\bibname}{R\'ef\'erences}

%--------------------------------------------------------
% Pour exclure des "environnement", ici techworkbox par exemple, travaux pour étudiants, non affiché dans le BOOK.
\usepackage{comment}


%--------------------------------------------------------
% pour filigramme sur les pages
\usepackage{eso-pic}

%--------------------------------------------------------
% for \gappto{\UrlBreaks}{\UrlOrds}
% rather than 
%\makeatletter
% \g@addto@macro{\UrlBreaks}{\UrlOrds}
% \makeatother

\usepackage{etoolbox}

\usepackage{breakurl}


\newcommand{\ucite}[1]{\cite{#1}}
%\newcommand{\ucite}[1]{REF BIB #1}
\newcommand{\ugls}[1]{\gls{#1}}
%\newcommand{\ugls}[1]{GLOSSAIRE (#1)}
%\usepackage[Rejne]{fncychap}

%\usepackage{pythontex} % need shell escape

%
%\makeatletter
%\DeclareRobustCommand{\afilename}[1]{%
% \begingroup
%  % \lstname seems to change hyphens into \textendash
%  \def\textendash{-}%
%  \afilename@parse{#1}%
%  \edef\afilename@base{\detokenize\expandafter{\afilename@base}}%
%  \texttt{\afilename@base.\afilename@ext}%
% \endgroup
%}
%\makeatother

%-------------------------------------------------------------
%               FR CYBERDEF SECOPS COURSE
%                                  EDX PACKAGES
%.    Commons COMMANDS for BOOK, NOTES (.doc) 
%                            and BEAMER (.prz)
%                              2020 eduf@ction
%-------------------------------------------------------------

\newcommand{\ucite}[1]{\cite{#1}}
%\newcommand{\ucite}[1]{REF BIB #1}
\newcommand{\ugls}[1]{\gls{#1}}
%\newcommand{\ugls}[1]{GLOSSAIRE (#1)}
\newcommand{\uac}[1]{\acrshort{#1}}



% \newcommand{\ExtractBeforeLastSlash}[1]{\directlua{extract_before_last_slash("#1")}}
% \newcommand{\ExtractAfterLastSlash}[1]{\directlua{extract_after_last_slash("#1")}}
% \newcommand{\RemoveExtension}[1]{\directlua{remove_extension("#1")}}
% \newcommand{\ExtractExtension}[1]{\directlua{extract_extension("#1")}}


\newcommand{\includepicturesubpath}[1]
{#1}

%{\ExtractBeforeLastSlash{#1}/\upicturethemepath\RemoveExtension{\ExtractAfterLastSlash{#1}}.\ExtractExtension{#1}}
    

%---------------------------------------------------
%										 	AUDIO NOTE 
%---------------------------------------------------

\newcommand\audionote[2]{%
  \note{#2}
  \immediate\write\audioexport{==AudioSlide=============}
  \immediate\write\audioexport{#1}
  \immediate\write\audioexport{\insertframenumber}
  \immediate\write\audioexport{---------------------------------}
  \immediate\write\audioexport{#2}
  \immediate\write\audioexport{---------------------------------}
}%

%---------------------------------------------------
%										 	TEXT STYLE SHORTCUT 
%---------------------------------------------------

\newcommand{\tcb}[1]{\textbf{\textcolor{uMainColor}{#1}}}

\newcommand{\voir}[1]{(Schéma~\ref{#1},~page~\pageref{#1}}

\newcommand{\graphicxextension}{pdf}

\newcommand\myemptypage{
   
    \thispagestyle{empty}
    \addtocounter{page}{-1}
    \newpage
  \null
    %CECI EST UNE PAGE VIDE
       \vfill
    }


\newcommand{\uNA}[2]
{
\newacronym{a#1}{#1}{#2}
}

\newcommand{\coauthor}[4]
{
\noindent (#2) \textbf{#1} - #4 : \textit{#3}
}


\newcommand\UKword[1]{\emph{#1}\xspace}
\newcommand\textUK[1]{\textit{\textbf{#1}}\xspace}
\newcommand\Ucenter[1]{\begin{center}{#1}\end{center}}

\newcommand\head[1]{\textbf{#1}}
\newcommand\tb[1]{\textbf{#1}}
\newcommand{\g}[1]{\og #1 \fg{}}
\newcommand{\uload}[1] {\input{../Latex/#1}}
\newcommand{\ucurcolor}{black}

\newcommand{\mak}[1]{\faicon{\aArrowCircleORight} #1}


%-----------------------------------------------------
%										 	\INFOGitHub 
%---------------------------------------------------


\newcommand{\GITfilename}{https://github.com/edufaction/CYBERDEF101/raw/main/Release/Builder/\jobname.pdf}

%\newcommand{\safeqrcode}[2][]{%
%\qrcode[#1]{\detokenize{#2}} }

\newcommand{\safeqrcode}[2][]{\qrcode[#1]{#2}} 

\newcommand{\INFOGitHub}
{
\begin{center} 
\setstretch{1.5}
 {Vérifiez la disponibilité d'une version plus récente de} \\
 {\ulink{\GITfilename}{\textbf{\jobname.pdf }sur GITHUB CYBERDEF} }   \\ 
{{\huge\ccbyncndeu}}  \\  
{Publication en \href{https://creativecommons.org/licenses/by-nc-nd/2.0/fr/}{Creative Common BY-NC-ND} by eduf@ction}    \\  % Copyright notice
{\safeqrcode[height=4 cm,padding]{\GITfilename}}  
\end{center} 
}

\newcommand{\INFOGitHubS}
{
\noindent{\safeqrcode[height=3 cm]{\GITfilename}} \\
\noindent {\tiny Téléchargez une version à jour} 
}



%-----------------------------------------------------
%										 	\upspicture 
%---------------------------------------------------
\newcommand{\upspicture}[3]
{
\renewcommand{\ucurcolor}{#3}
\resizebox{#2\textwidth}{!}{\input{../#1}}
}


%-----------------------------------------------------
%										 	\uindex 
%---------------------------------------------------


\newcommand{\uindex}[1]
{
#1\index{#1}\xspace
}

\newcommand{\ulindex}[1]{\index{#1}}

%-----------------------------------------------------
%										 	\includer 
%---------------------------------------------------

\newcommand{\includer}[1]
{
\input{../Latex/Sources/EDU/SRC1/Chapters/#1}
}

%-----------------------------------------------------
%										 	\edxdico 
%---------------------------------------------------

\newcommand{\edxdico}[2] {\textbf{#1}:\textit{#2} \index{#1}}

%-----------------------------------------------------
%  command DEF : InTextx
%-----------------------------------------------------
\newcommand\InTexta{no text}
\newcommand\InTextb{no text}
\newcommand\InTextc{no text}
\newcommand\rnc{\renewcommand}


%-----------------------------------------------------
%										 	ENV warningbox
%---------------------------------------------------

\newtcolorbox{warningbox}[2][]
{
  colframe = uMainColor!25,
  colback  = black!05,
  coltitle = uMainColor!20!black,
  fonttitle = \bfseries, 
  title    = #2,
  #1,
}


\newtcolorbox{techworkbox}[2][]{sharp corners,colframe = uMainColor!25,colback  = black!05,coltitle = uMainColor!20!black,fonttitle = \bfseries, title={\textcolor{uMainColor}{\faFile~\small {Sujet de mémoire SECOPS}}} : #2,#1,}


\newtcolorbox{notebox2}[2][]
{
sharp corners, 
colback = uMainColor!5!white, 
colframe = uMainColor!75!black,
fonttitle = \bfseries, 
colbacktitle= uMainColor!85!black,
title=#2,#1
}

\newtcolorbox{notebox}[2][]{colback=uMainColor!5!white,
colframe=uMainColor!75!black,fonttitle=\bfseries,
colbacktitle=uMainColor!85!black,enhanced,
attach boxed title to top right={yshift=+1mm},
title=#2,#1}


\newtcolorbox{toolsboxin}[2][]{sharp corners, colback=uMainColor!0!white,
colframe= black!25,fonttitle=\bfseries,
colbacktitle=black!60,enhanced,
attach boxed title to top left={yshift=+1mm},
title=#2,#1}


\newtcolorbox{techworkin}[2][]{sharp corners, colback=uMainColor!5!white,
colframe= uMainColor!100, fonttitle=\bfseries,
colbacktitle=black!5, enhanced,
 coltitle=black!50,
attach boxed title to top right ={yshift=+1mm},
title=#2,#1}



%-----------------------------------------------------
%										 	ENV ToolsBox
%---------------------------------------------------


\newcommand\tooleditor{No EDITOR in Datas.tex}
\newcommand\toolurl{No URL INSERTED in Datas.tex}
\newcommand\toolclass{No EDITOR in Datas.tex}
\newcommand\toolname{No NAME in Datas.tex}
\newcommand\toolanalyst{No ANALYST in Datas.tex}

\newcommand{\toolsbox}[2]
{
\input{../Latex/Sources/EDU/SRC1/SecComponents/#2.#1/datas.tex} 
\begin{toolsboxin}[sidebyside,righthand width=.7\textwidth]{#1}
\includegraphics[width=0.9\textwidth]{../Latex/Sources/EDU/SRC1/SecComponents/#2.#1/icon.png} \\
\tcblower
\input{../Latex/Sources/EDU/SRC1/SecComponents/#2.#1/description.tex}
\end{toolsboxin}
\noindent \small {\faCogs~Classe : \uppercase{\textbf{#2}}, Site de référence : \textbf{\ulink{\toolurl}{#1}}}
\newline
\small {\faGlobe~Editeur : \textbf{\tooleditor}}
\small {\faEye~Analyste : \textbf{\toolanalyst}}
}

%\renewcommand\toolsclass{SIEM}
%\renewcommand\toolsname{RSA Netwitness}
%\renewcommand\toolseditor{RSA}
%\renewcommand\toolsurl{https://www.rsa.com/fr-fr/products/threat-detection-response}

%-----------------------------------------------------
%										 	ENV  TechWork
%---------------------------------------------------


\newcommand{\techwork}[2]
{
\begin{techworkin} [righthand width=.8\textwidth]{\faCogs~\small {Fiche TECHNO} : #1}#2
\end{techworkin}
}

%-----------------------------------------------------
%										 	\uchap
%---------------------------------------------------
\ifdefined\INTERTITLE
\providecommand{\uchap}[1]
	{
	\begin{LARGE}
	\textbf{\Ucenter{#1 \\ -oOo-}}	
	\end{LARGE}
	}
\else
	\providecommand{\uchap}[1]{}
\fi

%-----------------------------------------------------
%											\uexpand
%---------------------------------------------------
\newcommand {\uexpand}[1]{
\mode<all>\input{../Latex/Sources/EDU/SRC1/Chapters/#1}
}


%-----------------------------------------------------
%											\spicture 
% direct include graphic with Scrivener
%-----------------------------------------------------

\newcommand{\spicture}[3]{
\begin{figure}[!hbtp] %h !hbtp
  \begin{center}
	 \includegraphics[width=0.8\textwidth]{#1}
  \end{center}
\caption{#3 - (slbl-#2)\label{slbl-#2}}
\end{figure}
}

%\StrBefore{#2}{.pdf}

%-----------------------------------------------------
%											\upicture 
%---------------------------------------------------

\ifdefined\PRZMODE

\providecommand{\upicture}[4]{
\framesubtitle{#2}
%\begin{figure}[!h] %h
  \begin{center}
	 \includegraphics[height=#3\textheight, width=#3\textwidth, keepaspectratio ]{\includepicturesubpath{#1.pdf}}
  \end{center}
%  \end{figure}
}

\else

\providecommand{\upicture}[4]{
\begin{figure}[!hbtp] %h !hbtp
  \begin{center}
	 \includegraphics[width=#3\textwidth]{\includepicturesubpath{#1.pdf}}
  \end{center}
\caption{#2 \label{#4}}
\end{figure}
}
	
\fi


\ifdefined\PRZMODE

\providecommand{\uwpicture}[4]{
\framesubtitle{#2}
%\begin{figure}[!h] %h
  \begin{center}
	 \includegraphics[height=0.8\textheight, width=0.8\textwidth, keepaspectratio ]{\includepicturesubpath{#1.pdf}}
  \end{center}
%  \end{figure}
}
\else
\providecommand{\uwpicture}[4]
{
\begin{wrapfigure}{O}{#3\textwidth} % inside / Outside
  \begin{center}
	 \includegraphics[width=#3\textwidth]{\includepicturesubpath{#1.pdf}}
\end{center}
\caption{#2 \label{#4}}
\end{wrapfigure}
}
\fi

%-----------------------------------------------------
%	\picframe for BEAMER (Image included in Beamer Slide)
%---------------------------------------------------

\ifdefined\PRZMODE

\providecommand{\picframe}[4]
{
\begin{frame}
\frametitle<presentation>{#2}
  \begin{center}

   \begin{tikzpicture}[remember picture, overlay]
    \node at (current page.center) {\includegraphics[width=#3\paperwidth, height=#3\paperheight,keepaspectratio]{\includepicturesubpath{#1.pdf}}};
    \end{tikzpicture}

  \end{center}
\end{frame}

}
\else
\providecommand{\picframe}[4]
{
\begin{figure}[!hbtp] % !hbtp ou !h 
\centering 
%{\includegraphics[width=#3\paperwidth, height=#3\paperheight,keepaspectratio]{\includepicturesubpath{#1.pdf}}}
\includegraphics[width=#3\textwidth]{\includepicturesubpath{#1.pdf}}
\caption{#2 \label{#4}}
\end{figure}
}
\fi

%-----------------------------------------------------
%.  for BEAMER (Image included in Beamer Slide)
%	\picframeS Simplified with Name and Label from One Param
%---------------------------------------------------


\ifdefined\PRZMODE

\providecommand{\picframeS}[2]
{
\begin{frame}
\frametitle<presentation>{#2}
  \begin{center}
	 \includegraphics[height=0.95\textheight, width=0.95\textwidth, keepaspectratio ]{\picpath/img-#1.\graphicxextension}
  \end{center}
\end{frame}

}
\else
\providecommand{\picframeS}[2]
{
\begin{figure}[!hbtp] % !hbtp ou !h 
\centering \includegraphics[width=0.9\textwidth]{\picpath/img-#1.\graphicxextension}
\caption{(s) #2 \label{lbl:#1}}
\end{figure}
}

\fi


%{figure}[!h]
%{wrapfigure}{O}{#3\textwidth}
%-----------------------------------------------------
%	FRAME			\wpicframe 
%---------------------------------------------------

\ifdefined\PRZMODE

\providecommand{\wpicframe}[4]
{
\begin{frame}
\frametitle<presentation>{#2}
  \begin{center}
	 \includegraphics[height=0.95\textheight, width=0.95\textwidth, keepaspectratio ]{#1.\graphicxextension}
  \end{center}
\end{frame}

}
\else
\providecommand{\wpicframe}[4]
{
\begin{wrapfigure}{O}{#3\textwidth}
\centering \includegraphics[width=#3\textwidth]{#1.\graphicxextension}
\caption{#2 \label{#4}}
\end{wrapfigure}
}
\fi

%-----------------------------------------------------
%	FRAME										\texframe 
%---------------------------------------------------
\ifdefined\PRZMODE

\providecommand{\texframe}[3]
{
\begin{frame}
\frametitle<presentation>{#1}
\framesubtitle<presentation>{#2}
#3
\end{frame}
}
\else
\providecommand{\texframe}[3]
{
#3
}

\fi

%-----------------------------------------------------
%											\uref 
%---------------------------------------------------
%\newcommand{\uref}[2]{(Voir~#1~\cref{#2} page~\pageref{#2})}

\newcommand{\uref}[2]{(Voir~#1~\ref{#2} page~\pageref{#2})}
%-----------------------------------------------------
%											\ubg
%---------------------------------------------------
\newcommand{\ubg}[1]{
\textbf{\g{\uppercase{#1}}}
}
%-----------------------------------------------------
%											\updfimage 
%---------------------------------------------------

\newcommand{\updfimage}[3]
{
\begin{wrapfigure}{R}{#3\textwidth}
  \begin{center}
	 \resizebox{#3\textwidth}{!}{\includegraphics{\includepicturesubpath{#1.pdf}}}
  \end{center}
\caption{#2}
\end{wrapfigure}
}


%************************PACKAGE***

%\usepackage{tikz}% Required for drawing custom shapes
\usetikzlibrary{shadows}


%-----------------------------------------------------
%											 		\rem 
% insertion d'une remarque avec R
%---------------------------------------------------
\newcommand\rem[1]{%
   \marginpar{%
   \tikzpicture[baseline={(title.base)}]
      \node[inner sep=5pt,text width=4cm,drop shadow={shadow yshift=-5pt,shadow xshift=5pt,uMainColor},fill=white] (box) {\vskip5pt \nointerlineskip #1};
      \node[right=10pt,inner sep=0pt,outer sep=10pt] (title) at (box.north west) {\bfseries\color{uMainColor}Remarque};
      \draw[draw=uMainColor,very thick](title.west)--(box.north west)--(box.south west)--(box.south east)--(box.north east)--(title.east);
      \fill[uMainColor]([yshift=-10pt]box.north west)--+(-5pt,-5pt)--+(0pt,-10pt);
   \endtikzpicture}%
}
%-----------------------------------------------------
%											 \utikzimage 
% insertion d'une image sur fichier tkz.tex
%---------------------------------------------------

\newcommand{\utikzimage}[3]
{
\begin{wrapfigure}{R}{#3\textwidth}
  \begin{center}
	 \resizebox{#3\textwidth}{!}{\input{../#1.tkz.tex}}
  \end{center}
\caption{#2}
\end{wrapfigure}
}

%-----------------------------------------------------
%														\link 
% Url et footnote du lien
%-----------------------------------------------------

% TT every character and hyphenate after it
\def\hyphenateAndTtWholeString #1{\xHyphenate#1$\wholeString\unskip}

\def\xHyphenate#1#2\wholeString {\if#1$%
    \else\transform{#1}%
    \takeTheRest#2\ofTheString\fi}

\def\takeTheRest#1\ofTheString\fi
{\fi \xHyphenate#1\wholeString}

\def\transform#1{\url{#1}\hskip 0pt plus 1pt}

% Define the \urlx command which works like \url, but with line brakes
\def\urlx #1{\href{#1}{\hyphenateAndTtWholeString{#1}}}
  

 %-----------------------------------------------------
%														\link 
%-----------------------------------------------------  
  
  \newcommand{\ulink}[2]
  {\href{#1}
  {#2~\raisebox{-0.2ex}{\faExternalLink}\footnote{\url{#1}}}\xspace}
  
  \newcommand{\wikipedia}[1]
  { \href{https://fr.wikipedia.org/wiki/#1}{\textcolor{grey}{#1}}\raisebox{0.5ex}{\tiny{\textcolor{uMainColor}{\faWikipediaW}}}}
  


 %-----------------------------------------------------
%														\utodo 
%-----------------------------------------------------

  \newcommand{\utodo}{\begin{center}{\noindent \textcolor{cyan}{\faGears~{En cours de rédaction, DRAFT non publiable} \faGears}}\end{center}}
  
    \newcommand{\utocomplete}{\begin{center}{\noindent \textcolor{blue}{\faClipboard~{à compléter, les éléments ne donnent qu'une vue trop réduite ou trop parcellaire du sujet}}}\end{center}}
  

 




%-------------------------------------------------------------
%               FR CYBERDEF SECOPS COURSE
%
%                                    EDX STYLE
%
%                              2020 eduf@ction
%-------------------------------------------------------------
\usepackage{enumitem}

\usepackage{parskip}
\setlength{\parindent}{0em}
\setlength{\parskip}{5pt} % 1ex plus 0.5ex minus 0.2ex}

\def\uCoverColor{}

%---------------------------------------------------------------
%	MARGINS
%---------------------------------------------------------------

\usepackage{geometry}

\geometry{
	paper=a4paper, 
%	top=1cm, 
%	bottom=1cm, 
%	left=1cm, 
%	right=1cm, 
%	headheight=20pt, 
%	footskip=1.4cm, 
%	headsep=10pt, 
%	showframe, 
	showcrop 
}

%---------------------------------------------------------------
%	FOOTNOTE (format)
%---------------------------------------------------------------

\makeatletter
\long\def\@makefntextFB#1{%
    \ifx\thefootnote\ftnISsymbol
        \@makefntextORI{#1}%
    \else
        \rule\z@\footnotesep
        \setbox\@tempboxa\hbox{\@thefnmark}%
            \ifdim\wd\@tempboxa>\z@
                \kern2em\llap{\@thefnmark.\kern0.5em}%
            \fi
        \hangindent2em\hangafter\@ne#1
    \fi}
\makeatother

%-----------------------------------------------------
%	COLOR
%-----------------------------------------------------
%\usepackage{xcolor} % Required for specifying colors by name
%\definecolor{uMainColor}{RGB}{160,0,0} 
\colorlet{edxcolorcover}{uCoverColor}
%\definecolor{edxcolorcover}{RGB}{160,0,0} 

%\definecolor{cnam}{RGB}{128, 0, 32} 
                 

  \definecolor{comments}{rgb}{0.7,0,0}    % rouge foncé
  \definecolor{link}{rgb}{0,0.4,0.6}      % ~RoyalBlue de dvips
  \definecolor{url}{rgb}{0.6,0,0}         % rouge-brun
  \definecolor{citation}{rgb}{0,0.5,0}    % vert foncé
  \definecolor{ULlinkcolor}{rgb}{0,0,0.3} % de ulthese.cls
  \definecolor{rouge}{rgb}{0.85,0,0.07}   % rouge bandeau identitaire
  \definecolor{or}{rgb}{1,0.8,0}          % or bandeau identitaire

%-----------------------------------------------------
%	FONTS
%-----------------------------------------------------

%%\usepackage{avant} % Use the Avantgarde font for headings
%\usepackage{times} % Use the Times font for headings
%\usepackage{mathptmx} % Use the Adobe Times Rgman as the default text font together with math symbols from the Sym­bol, Chancery and Com­puter Modern fonts

%%\usepackage{microtype} % Slightly tweak font spacing for aesthetics
%%\usepackage[utf8]{inputenc} % Required for including letters with accents
%\usepackage[T1]{fontenc} % Use 8-bit encoding that has 256 glyphs

% Using AVANT GARDE Family

%\renewcommand{\familydefault}{\sfdefault}

%\fontfamily{pag}\selectfont
%\renewcommand{\familydefault}{pag} % sfdefault pag lmss

%bch         Charter
%lmr         Latin Modern Roman
%lmss        Latin Modern Sans Serif
%lmssq       Latin Modern Sans Serif extended
%lmtt        Latin Modern Typewriter
%lmvtt       Latin Modern Typewriter proportional
%pag         Avant Garde
%pbk         Bookman
%pcr         Courier
%phv         Helvetica
%pnc         New Century Schoolbook
%ppl         Palatino
%ptm         Times
%put         Utopia
 
%\renewcommand{\rmdefault}{pag} % text 
%\renewcommand{\sfdefault}{pag} % titre


%\renewcommand{\rmdefault}{pag} % text 
%\renewcommand{\sfdefault}{phv} % titre


%-----------------------------------------------------
%	BIBLIOGRAPHY AND INDEX
%-----------------------------------------------------


\usepackage{calc} % For simpler calculation - used for spacing the index letter headings correctly

\usepackage{imakeidx}
%\makeindex % Tells LaTeX to create the files required for indexing

\makeindex[columns=2,intoc=true,options={-s \upath/indexstyle.ist}]

%-----------------------------------------------------
%	HEADERS AND FOOTERS
%-----------------------------------------------------


\usepackage{fancyhdr} % Required for header and footer configuration
\pagestyle{fancy} % Enable the custom headers and footers

\ifcsname chapter \endcsname   % verify if command name exists in the CLASS (Book, report, article ...)
% BEGIN IF
\renewcommand{\chaptermark}[1]{\markboth{\sffamily\normalsize\bfseries\chaptername\ \thechapter.\ #1}{}} % Styling for the current chapter in the header
  \else
%nothing
\fi
% END IF

\renewcommand{\sectionmark}[1]{\markright{\sffamily\normalsize\thesection\hspace{5pt}#1}{}} % Styling for the current section in the header

\fancyhf{} % Clear default headers and footers
\fancyhead[LE,RO]{\sffamily\normalsize\thepage} % Styling for the page number in the header
\fancyhead[LO]{\sffamily\normalsize\rightmark} % Print the nearest section name on the left side of odd pages
\fancyhead[RE]{\sffamily\normalsize\leftmark} % Print the current chapter name on the right side of even pages
\fancyfoot[RE,LO]{\sffamily\normalsize\uCoursetittle} % Uncomment to include a footer
\fancyfoot[LE,RO]{\hyperlink{toc}{\sffamily\normalsize\uMinilogo}} % Uncomment to include a footer
% \hyperlink{toc}
\renewcommand{\headrulewidth}{1.0pt} % Thickness of the rule under the header

\fancypagestyle{plain}{% Style for when a plain pagestyle is specified
	\fancyhead{}\renewcommand{\headrulewidth}{0pt}} 

% Removes the header from odd empty pages at the end of chapters
\makeatletter
\renewcommand{\cleardoublepage}{
\clearpage\ifodd\c@page\else
\hbox{}
\vspace*{\fill}
\thispagestyle{empty}
\newpage
\fi}


%----------------------------------------------------------------------------------------
%	MAIN TABLE OF CONTENTS
%----------------------------------------------------------------------------------------
%\usepackage{titletoc} % Required for manipulating the table of contents

\contentsmargin{0cm} % Removes the default margin

% Part text styling (this is mostly taken care of in the PART HEADINGS section of this file)
\titlecontents{part}
	[0cm] % Left indentation
	{\addvspace{20pt}\bfseries} % Spacing and font options for parts
	{}
	{}
	{}

% Chapter text styling
\titlecontents{chapter}
	[1.25cm] % Left indentation
	{\addvspace{12pt}\large\sffamily\bfseries} % Spacing and font options for chapters
	{\color{uMainColor!60}\contentslabel[\Large\thecontentslabel]{1.25cm}\color{uMainColor}} % Formatting of numbered sections of this type
	{\color{uMainColor}} % Formatting of numberless sections of this type
	{\color{uMainColor!60}\normalsize\;\titlerule*[.5pc]{.}\;\thecontentspage} % Formatting of the filler to the right of the heading and the page number

% Section text styling
\titlecontents{section}
	[1.25cm] % Left indentation
	{\addvspace{3pt}\sffamily\bfseries} % Spacing and font options for sections
	{\contentslabel[\thecontentslabel]{1.25cm}} % Formatting of numbered sections of this type
	{} % Formatting of numberless sections of this type
	{\hfill\color{black}\thecontentspage} % Formatting of the filler to the right of the heading and the page number

% Subsection text styling
\titlecontents{subsection}
	[1.25cm] % Left indentation
	{\addvspace{1pt}\sffamily\small} % Spacing and font options for subsections
	{\contentslabel[\thecontentslabel]{1.25cm}} % Formatting of numbered sections of this type
	{} % Formatting of numberless sections of this type
	{\ \titlerule*[.5pc]{.}\;\thecontentspage} % Formatting of the filler to the right of the heading and the page number

%\titlecontents{subsubsection}
%	[1.25cm] % Left indentation
%	{\addvspace{1pt}\sffamily\tiny} % Spacing and font options for subsections
%	{\contentslabel[\thecontentslabel]{1.25cm}} % Formatting of numbered sections of this type
%	{} % Formatting of numberless sections of this type
%	{\ \titlerule*[.5pc]{.}\;\thecontentspage} % Formatting of the filler to the right of the heading and the page number

\titlecontents{subsubsection}
[1.25cm]
  {\addvspace{1pt}\sffamily\tiny}
 {\contentslabel[\thecontentslabel]{1.25cm}}
  {}
  {}

% Figure text styling
\titlecontents{figure}
	[1.25cm] % Left indentation
	{\addvspace{1pt}\sffamily\small} % Spacing and font options for figures
	{\thecontentslabel\hspace*{1em}} % Formatting of numbered sections of this type
	{} % Formatting of numberless sections of this type
	{\ \titlerule*[.5pc]{.}\;\thecontentspage} % Formatting of the filler to the right of the heading and the page number

% Table text styling
\titlecontents{table}
	[1.25cm] % Left indentation
	{\addvspace{1pt}\sffamily\small} % Spacing and font options for tables
	{\thecontentslabel\hspace*{1em}} % Formatting of numbered sections of this type
	{} % Formatting of numberless sections of this type
	{\ \titlerule*[.5pc]{.}\;\thecontentspage} % Formatting of the filler to the right of the heading and the page number

%----------------------------------------------------------------------------------------
%	MINI TABLE OF CONTENTS IN PART HEADS
%----------------------------------------------------------------------------------------

% Chapter text styling
\titlecontents{lchapter}
	[0em] % Left indentation
	{\addvspace{15pt}\large\sffamily\bfseries} % Spacing and font options for chapters
	{\color{uMainColor}\contentslabel[\Large\thecontentslabel]{1.25cm}\color{uMainColor}} % Chapter number
	{}  
	{\color{uMainColor}\normalsize\sffamily\bfseries\;\titlerule*[.5pc]{.}\;\thecontentspage} % Page number

% Section text styling
\titlecontents{lsection}
	[0em] % Left indentation
	{\sffamily\small} % Spacing and font options for sections
	{\contentslabel[\thecontentslabel]{1.25cm}} % Section number
	{}
	{}

% Subsection text styling (note these aren't shown by default, display them by searchings this file for tgcdepth and reading the commented text)
%\titlecontents{lsubsection}
%	[.5em] % Left indentation
%	{\sffamily\footnotesize} % Spacing and font options for subsections
%	{\contentslabel[\thecontentslabel]{1.25cm}}
%	{}
%	{}


\usepackage[]{ccicons}


%----------------------------------------------------------------------------------------
%	THEOREM STYLES
%----------------------------------------------------------------------------------------

\usepackage{amsmath,amsfonts,amssymb,amsthm} % For math equations, theorems, symbols, etc

\newcommand{\intoo}[2]{\mathopen{]}#1\,;#2\mathclose{[}}
\newcommand{\ud}{\mathop{\mathrm{{}d}}\mathopen{}}
\newcommand{\intff}[2]{\mathopen{[}#1\,;#2\mathclose{]}}
\renewcommand{\qedsymbol}{$\blacksquare$}

\ifcsname chapter \endcsname   % verify if command name exists in the CLASS (Book, report, article ...)
	\newtheorem{notation}{Notation}[chapter]
 \else
	\newtheorem{notation}{Notation}[section]
\fi


% Boxed/framed environments
\newtheoremstyle{ocrenumbox}% Theorem style name
{0pt}% Space above
{0pt}% Space below
{\normalfont}% Body font
{}% Indent amount
{\small\bf\sffamily\color{uMainColor}}% Theorem head font
{\;}% Punctuation after theorem head
{0.25em}% Space after theorem head
{\small\sffamily\color{uMainColor}\thmname{#1}\nobreakspace\thmnumber{\@ifnotempty{#1}{}\@upn{#2}}% Theorem text (e.g. Theorem 2.1)
\thmnote{\nobreakspace\the\thm@notefont\sffamily\bfseries\color{black}---\nobreakspace#3.}} % Optional theorem note

\newtheoremstyle{ocrenumbox}% Theorem style name
{0pt}% Space above
{0pt}% Space below
{\normalfont}% Body font
{}% Indent amount
{\small\bf\sffamily\color{uMainColor}}% Theorem head font
{\;}% Punctuation after theorem head
{0.25em}% Space after theorem head
{\small\sffamily\color{uMainColor}\thmname{#1}\nobreakspace\thmnumber{\@ifnotempty{#1}{}\@upn{#2}}% Theorem text (e.g. Theorem 2.1)
\thmnote{\nobreakspace\the\thm@notefont\sffamily\bfseries\color{black}---\nobreakspace#3.}} % Optional theorem note

\newtheoremstyle{blacknumex}% Theorem style name
{5pt}% Space above
{5pt}% Space below
{\normalfont}% Body font
{} % Indent amount
{\small\bf\sffamily}% Theorem head font
{\;}% Punctuation after theorem head
{0.25em}% Space after theorem head
{\small\sffamily{\tiny\ensuremath{\blacksquare}}\nobreakspace\thmname{#1}\nobreakspace\thmnumber{\@ifnotempty{#1}{}\@upn{#2}}% Theorem text (e.g. Theorem 2.1)
\thmnote{\nobreakspace\the\thm@notefont\sffamily\bfseries---\nobreakspace#3.}}% Optional theorem note


\newtheoremstyle{blacknumboxS}% name of the style to be used
{5pt}% measure of space to leave above the theorem. E.g.: 3pt
{5pt}% measure of space to leave below the theorem. E.g.: 3pt
{\normalfont}% name of font to use in the body of the theorem
{}% measure of space to indent
{\small\bf\sffamily}% name of head font
{}% punctuation between head and body
{ }% space after theorem head; " " = normal interword space
{\color{black}\faEye~  \color{uMainColor}\small\sffamily\thmnote{#3 : }}

% Non-boxed/non-framed environments
\newtheoremstyle{ocrenum}% Theorem style name
{5pt}% Space above
{5pt}% Space below
{\normalfont}% Body font
{}% Indent amount
{\small\bf\sffamily\color{uMainColor}}% Theorem head font
{\;}% Punctuation after theorem head
{0.25em}% Space after theorem head
{\small\sffamily\color{uMainColor}\thmname{#1}\nobreakspace\thmnumber{\@ifnotempty{#1}{}\@upn{#2}}% Theorem text (e.g. Theorem 2.1)
\thmnote{\nobreakspace\the\thm@notefont\sffamily\bfseries\color{black}---\nobreakspace#3.}} % Optional theorem note
\makeatother

% Defines the theorem text style for each type of theorem to one of the three styles above



\ifcsname chapter \endcsname

\newcounter{dummy} 

\numberwithin{dummy}{section}
	\theoremstyle{ocrenumbox}
		\newtheorem{theoremeT}[dummy]{Proposition}
		\newtheorem{problem}{Problem}[chapter]
		\newtheorem{exerciseT}{Outillage}[chapter]
	\theoremstyle{blacknumex}
		\newtheorem{exampleT}{Example}[chapter]
	\theoremstyle{blacknumbox}
		\newtheorem{vocabulary}{Vocabulary}[chapter]
		\newtheorem{definitionT}{Definition}[section]
		\newtheorem{corollaryT}[dummy]{Concept}
	\theoremstyle{blacknumboxS}
		\newtheorem{remarqueST}[]{}
		\newtheorem{remarqueT}[dummy]{Remarque}
	\theoremstyle{ocrenum}
	\newtheorem{proposition}[dummy]{Proposition}
  
 \else
 
\newcounter{dummy} 
\numberwithin{dummy}{section}

	\theoremstyle{ocrenumbox}
		\newtheorem{theoremeT}[dummy]{Proposition}
		%\newtheorem{problem}{Problem}[section]
		\newtheorem{exerciseT}{Outillage}[section]
	\theoremstyle{blacknumex}
		\newtheorem{exampleT}{Example}[section]
	\theoremstyle{blacknumbox}
		\newtheorem{vocabulary}{Vocabulary}[section]
		\newtheorem{definitionT}{Definition}[section]
		\newtheorem{corollaryT}[dummy]{concept}
	\theoremstyle{blacknumboxS}
		\newtheorem{remarqueST}[]{}
		\newtheorem{remarqueT}[dummy]{Remarque}
	\theoremstyle{ocrenum}
		\newtheorem{proposition}[dummy]{Proposition}
\fi



%----------------------------------------------------------------------------------------
%	DEFINITION OF COLORED BOXES
%----------------------------------------------------------------------------------------

\RequirePackage[framemethod=default]{mdframed} % Required for creating the theorem, definition, exercise and corollary boxes

% Theorem box
\newmdenv[skipabove=7pt,
skipbelow=7pt,
backgroundcolor=black!5,
linecolor=uMainColor,
innerleftmargin=5pt,
innerrightmargin=5pt,
innertopmargin=5pt,
leftmargin=0cm,
rightmargin=0cm,
innerbottommargin=5pt]{tBox}

% Exercise box	  
\newmdenv[skipabove=7pt,
skipbelow=7pt,
rightline=false,
leftline=true,
topline=false,
bottomline=false,
backgroundcolor=uMainColor!10,
linecolor=uMainColor,
innerleftmargin=5pt,
innerrightmargin=5pt,
innertopmargin=5pt,
innerbottommargin=5pt,
leftmargin=0cm,
rightmargin=0cm,
linewidth=4pt]{eBox}	

% Definition box
\newmdenv[skipabove=7pt,
skipbelow=7pt,
rightline=false,
leftline=true,
topline=false,
bottomline=false,
linecolor=uMainColor,
innerleftmargin=5pt,
innerrightmargin=5pt,
innertopmargin=0pt,
leftmargin=0cm,
rightmargin=0cm,
linewidth=4pt,
innerbottommargin=0pt]{dBox}	

% Corollary box
\newmdenv[skipabove=7pt,
skipbelow=7pt,
rightline=false,
leftline=true,
topline=false,
bottomline=false,
linecolor=gray,
backgroundcolor=black!5,
innerleftmargin=5pt,
innerrightmargin=5pt,
innertopmargin=5pt,
leftmargin=0cm,
rightmargin=0cm,
linewidth=4pt,
innerbottommargin=5pt]{cBox}

% Creates an environment for each type of theorem and assigns it a theorem text style from the "Theorem Styles" section above and a colored box from above
%\newenvironment{theorem}{\begin{tBox}\begin{theoremeT}}{\end{theoremeT}\end{tBox}}
\newenvironment{exercise}{\begin{eBox}\begin{exerciseT}}{\hfill{\color{uMainColor}\tiny\ensuremath{\blacksquare}}\end{exerciseT}\end{eBox}}				  
%\newenvironment{definition}{\begin{dBox}\begin{definitionT}}{\end{definitionT}\end{dBox}}	
%\newenvironment{example}{\begin{exampleT}}{\hfill{\tiny\ensuremath{\blacksquare}}\end{exampleT}}		
%\newenvironment{corollary}{\begin{cBox}\begin{corollaryT}}{\end{corollaryT}\end{cBox}}	
\newenvironment{nota}{\begin{cBox}\begin{remarqueT}}{\end{remarqueT}\end{cBox}}	

%----------------------------------------------------------------------------------------
%	REMARK ENVIRONMENT
%----------------------------------------------------------------------------------------

\newenvironment{remark}{\par\vspace{10pt}\small % Vertical white space above the remark and smaller font size
\begin{list}{}{
\leftmargin=35pt % Indentation on the left
\rightmargin=25pt}\item\ignorespaces % Indentation on the right
\makebox[-2.5pt]{\begin{tikzpicture}[overlay]
\node[draw=uMainColor!60,line width=1pt,circle,fill=uMainColor!25,font=\sffamily\bfseries,inner sep=2pt,outer sep=0pt] at (-15pt,0pt){\textcolor{uMainColor}{i}};\end{tikzpicture}} % Orange R in a circle
\advance\baselineskip -1pt}{\end{list}\vskip5pt} % Tighter line spacing and white space after remark

%----------------------------------------------------------------------------------------
%	SECTION NUMBERING IN THE MARGIN
%----------------------------------------------------------------------------------------

\makeatletter
% \renewcommand{\@seccntformat}[1]{\textcolor{uMainColor}{\csname the#1\endcsname}\quad}   

%\makeatother

%----------------------------------------------------------------------------------------
%	PART HEADINGS
%----------------------------------------------------------------------------------------

% Numbered part in the table of contents
\newcommand{\@mypartnumtocformat}[2]{%
	\setlength\fboxsep{0pt}%
	\noindent\colorbox{uMainColor!20}{\strut\parbox[c][.7cm]{\ecart}{\color{uMainColor!70}\Large\sffamily\bfseries\centering#1}}\hskip\esp\colorbox{uMainColor!40}{\strut\parbox[c][.7cm]{\linewidth-\ecart-\esp}{\Large\sffamily\centering#2}}%
}

% Unnumbered part in the table of contents
\newcommand{\@myparttocformat}[1]{%
	\setlength\fboxsep{0pt}%
	\noindent\colorbox{uMainColor!40}{\strut\parbox[c][.7cm]{\linewidth}{\Large\sffamily\centering#1}}%
}

\newlength\esp
\setlength\esp{4pt}
\newlength\ecart
\setlength\ecart{1.2cm-\esp}
\newcommand{\thepartimage}{}%
\newcommand{\partimage}[1]{\renewcommand{\thepartimage}{#1}}%
\def\@part[#1]#2{%
\ifnum \c@secnumdepth >-2\relax%
\refstepcounter{part}%
\addcontentsline{toc}{part}{\texorpdfstring{\protect\@mypartnumtocformat{\thepart}{#1}}{\partname~\thepart\ ---\ #1}}
\else%
\addcontentsline{toc}{part}{\texorpdfstring{\protect\@myparttocformat{#1}}{#1}}%
\fi%
\startcontents%
\markboth{}{}%
{\thispagestyle{empty}%
\begin{tikzpicture}[remember picture,overlay]%
\node at (current page.north west){\begin{tikzpicture}[remember picture,overlay]%	
\fill[uMainColor!20](0cm,0cm) rectangle (\paperwidth,-\paperheight);
\node[anchor=north] at (4cm,-3.25cm){\color{uMainColor!40}\fontsize{220}{100}\sffamily\bfseries\thepart}; 
\node[anchor=south east] at (\paperwidth-1cm,-\paperheight+1cm){\parbox[t][][t]{8.5cm}{
\printcontents{l}{0}{\setcounter{tocdepth}{1}}% The depth to which the Part mini table of contents displays headings; 0 for chapters only, 1 for chapters and sections and 2 for chapters, sections and subsections
}};
\node[anchor=north east] at (\paperwidth-1.5cm,-3.25cm){\parbox[t][][t]{15cm}{\strut\raggedleft\color{white}\fontsize{30}{30}\sffamily\bfseries#2}};
\end{tikzpicture}};
\end{tikzpicture}}%
\@endpart}
\def\@spart#1{%
\startcontents%
\phantomsection
{\thispagestyle{empty}%
\begin{tikzpicture}[remember picture,overlay]%
\node at (current page.north west){\begin{tikzpicture}[remember picture,overlay]%	
\fill[uMainColor!20](0cm,0cm) rectangle (\paperwidth,-\paperheight);
\node[anchor=north east] at (\paperwidth-1.5cm,-3.25cm){\parbox[t][][t]{15cm}{\strut\raggedleft\color{white}\fontsize{30}{30}\sffamily\bfseries#1}};
\end{tikzpicture}};
\end{tikzpicture}}
\addcontentsline{toc}{part}{\texorpdfstring{%
\setlength\fboxsep{0pt}%
\noindent\protect\colorbox{uMainColor!40}{\strut\protect\parbox[c][.7cm]{\linewidth}{\Large\sffamily\protect\centering #1\quad\mbox{}}}}{#1}}%
\@endpart}
\def\@endpart{\vfil\newpage
\if@twoside
\if@openright
\null
\thispagestyle{empty}%
\newpage
\fi
\fi
\if@tempswa
\twocolumn
\fi}

%----------------------------------------------------------------------------------------
%	CHAPTER HEADINGS
%-------------------------------------------------%--------------------------------------

%
%
%% A switch to conditionally include a picture, implemented by Christian Hupfer
\newif\ifusechapterimage
\usechapterimagetrue
\newcommand{\thechapterimage}{}%
\newcommand{\chapterimage}[1]{\ifusechapterimage\renewcommand{\thechapterimage}{#1}\fi}%
\newcommand{\autodot}{.}

\def\@makechapterhead#1{%
{\parindent \z@ \raggedright \normalfont
\ifnum \c@secnumdepth >\m@ne
\if@mainmatter
\begin{tikzpicture}[remember picture,overlay]
	
\node at (current page.north west)
			{\begin{tikzpicture}[remember picture,overlay]
			\node[anchor=north west,inner sep=0pt] at (0,0) {\ifusechapterimage\includegraphics[width=\paperwidth]{\thechapterimage}\fi};
			\draw[anchor=west] (\Gm@lmargin,-9cm) node [line width=2pt,rounded corners=15pt,draw=uMainColor,fill=white,fill opacity=1,inner sep=15pt]{\strut\makebox[22cm]{}};
			
			\draw[anchor=west] (\Gm@lmargin+.3cm,-9cm) node {\fontsize{20}{30}\selectfont\sffamily\bfseries\color{uMainColor}~#1\strut};
			%{\Huge\sffamily\bfseries\color{uMainColor}\thechapter\autodot~#1\strut};
			\node at (current page.north east)  [xshift=-3.95cm, yshift=-4.05cm, text opacity=1]  {{\color{white}\centering\fontsize{200}{30}\selectfont \bfseries\sffamily\thechapter\strut} };
			\node at (current page.north east)  [xshift=-4cm, yshift=-4cm, text opacity=1]  {{\color{uMainColor}\centering\fontsize{200}{30}\selectfont \bfseries\sffamily\thechapter\strut} };
			
			\end{tikzpicture}};
	\end{tikzpicture}

\else

\begin{tikzpicture}[remember picture,overlay]
\node at (current page.north west)
{\begin{tikzpicture}[remember picture,overlay]
\node[anchor=north west,inner sep=0pt] at (0,0) {\ifusechapterimage\includegraphics[width=\paperwidth]{\thechapterimage}\fi};
\draw[anchor=west] (\Gm@lmargin,-9cm) node [line width=2pt,rounded corners=15pt,draw=uMainColor,fill=white,fill opacity=1,inner sep=15pt]{\strut\makebox[22cm]{}};

\draw[anchor=east] (\Gm@lmargin,-9cm) node [line width=2pt,rounded corners=15pt,draw=uMainColor,fill=white,fill opacity=1,inner sep=15pt]{\strut\makebox[22cm]{}};
\draw[anchor=west] (\Gm@lmargin+.3cm,-9cm) node {\huge\sffamily\bfseries\color{uMainColor}#1\strut};

\end{tikzpicture}};
 
\end{tikzpicture}
\fi\fi\par\vspace*{270\p@}}}

\makeatother
%-------------------------------------------

%\def\@makeschapterhead#1{%
%\begin{tikzpicture}[remember picture,overlay]
%\node at (current page.north west)
%{\begin{tikzpicture}[remember picture,overlay]
%\node[anchor=north west,inner sep=0pt] at (0,0) {\ifusechapterimage\includegraphics[width=\paperwidth]{\thechapterimage}\fi};
%\draw[anchor=west] (\Gm@lmargin,-9cm) node [line width=2pt,rounded corners=15pt,draw=uMainColor,fill=white, fill opacity=0.8,inner sep=15pt]{\strut\makebox[22cm]{}};
%\draw[anchor=west] (\Gm@lmargin+.3cm,-9cm) node {\huge\sffamily\bfseries\color{uMainColor}#1\strut};
%\end{tikzpicture}};
%\end{tikzpicture}
%\par\vspace*{270\p@}}




%----------------------------------------------------------------------------------------
%	LINKS
%----------------------------------------------------------------------------------------


\usepackage{bookmark}
\bookmarksetup{
open,
numbered,
addtohook={%
\ifnum\bookmarkget{level}=0 % chapter
\bookmarksetup{bold}%
\fi
\ifnum\bookmarkget{level}=-1 % part
\bookmarksetup{color=uMainColor,bold}%
\fi
}
}

%----------------------------------------------------------------------------------------
%	FANCY CHAPTER Header
%----------------------------------------------------------------------------------------





%\definecolor{uMainColor}{RGB}{255,128,0}   

%-------------------------------------------------
%	BIBLIOGRAPHY
%-------------------------------------------------

\addbibresource{Latex/Sources/commons/bibliography.bib} % BibTeX bibliography file
\defbibheading{bibempty}{}

%-------------------------------------------------
%	COLORS & BORDERS
%-------------------------------------------------
\setlength{\fboxrule}{0.75pt} % Width of the border around the abstract
%\definecolor{uMainColor}{RGB}{255,128,0} % Color of the article title and section
\definecolor{color2}{RGB}{220,220,220} % Color of the boxes behind the abstract and headings

%-------------------------------------------------
% ITEMS DEFINITION
%-------------------------------------------------

\setlist [itemize,1]{label=\color{uMainColor}\faCaretRight }

\begin{document}


\addtocontents{toc}{\protect\hypertarget{toc}{}}

%-------------------------------------------------
%	ARTICLE INFORMATION
%-------------------------------------------------

\Abstract {\uabstract%===========================
% COURS "INTRO CYBERDEF"
% Abstract général des articles
%===========================


Il fait partie du cours introductif aux fondamentaux de la sécurité des systèmes d'information vue sous deux prismes quelques fois opposés dans la littérature : la gouvernance et la gestion opérationnelle de la sécurité.
Le cours est constitué d'un ensemble de notes de synthèse indépendantes compilées en un document unique, mais édité par chapitre dans le cadre de ce cours.\\
C\edoc ne constitue pas à lui seul le référentiel du cours CYBERDEF101 (SEC101 du Cnam). Il compile des notes de cours mises à disposition de l'auditeur comme support pédagogique partiel à ce cours introductif à la cyberdéfense d'entreprise. }
\JournalInfo{\uJournalInfo} % Journal information
\Archive{Notes de cours éditées le  \DTMnow} % Additional notes (e.g. copyright, DOI, review/research article)
\PaperTitle{\utitle} % Article title
\Authors{\uauthor\textsuperscript{1,}\textsuperscript{2}*} % Authors
\affiliation{\textsuperscript{1}\textit{\uproa}} % Author affiliation
\affiliation{\textsuperscript{2}\textit{\uprob}} % Author affiliation
\affiliation{*\textbf{email}: \umaila\ -- \umailb} % Corresponding author
\Keywords{\ukeywords} % Keywords 
\newcommand{\keywordname}{Mots clefs} % Defines the keywords heading name

%-------------------------------------------------
%	MAKE TITLE
%-------------------------------------------------

%\maketitle

%\begin{tikzpicture}[remember picture,overlay]\node at (current page.north west)[xshift=4cm, yshift=-4cm, text opacity=1.0]{\includegraphics[width=0.1\paperwidth]{\utemplatepath/commons.inc/CommonsPictures/\utemplatetheme.theme/shield-20.pdf}}; 
%\end{tikzpicture} 
%
%
%\AddToShipoutPicture{%
%\AtPageCenter{%
%\begin{tikzpicture}[remember picture,overlay]
%\node [rotate=60,scale=5,text opacity=0.1] at (current page.center) { \textit{\color{uAlertColor}\INFODistrib} };
%\end{tikzpicture} 
%%\begin{tikzpicture}[remember picture,overlay]
%%\node [rotate=60,scale=10,text opacity=0.3]
%%at (current page.center) {moi};
%%\end{tikzpicture}
%}%
%}


%-------------------------------------------------	
%\begin{center}
%	\textcolor{uMainColor}{\textbf{\INFODistrib}} % Distribution information
%\end{center}
%
%\INFOGitHub % Link to download file and copyright
%%-------------------------------------------------

%%-------------------------------------------------------------
%               FR CYBERDEF SECOPS COURSE
%
%                 DEFINITION DES Sources/EDU/acronymes

%                      $File :  Sources/EDU/acronymes.tex
%                             2020 eduf@ction
%-------------------------------------------------------------
% Require : glossaries package
%-------------------------------------------------------------
%\newacronym{XXX}{XXX}{XXX}
%-----------------------------------------------------------

%  Organismes

\newacronym{aCNIL}{CNIL}{Commission National Informatique et Liberté}
\newacronym{aOTAN}{OTAN}{Organisation du Traité de l'Atlantique Nord}
\newacronym{aANSSI}{ANSSI}{Agence Nationale de la Sécurité des systèmes d'information}
\newacronym{aCERT}{CERT}{Computer Emergency Response Team}
\newacronym{aAFNOR}{AFNOR}{Association Française de Normalisation}
\newacronym{aSIEM}{SIEM}{System Incident and Event Management}
\newacronym{aSIM}{SIM}{System Incident Management}
\newacronym{aSEM}{SEM}{System Event Management}
\newacronym{aIRP}{IRP}{Incident Response Plan, Plan de Réponse à Incident}
\newacronym{aISO}{ISO}{International Organization for Standardization}
\newacronym{aDICT}{DICT}{Disponibilité, Intégrité, Confidentialité, Traçabilité}
\newacronym{aPDCA}{PDCA}{Plan, Do, Check, Act (Roue de Deming)}
\newacronym{aGDPR}{GDPR}{General Data Protection Regulation ou RGPD Règlement général sur la protection des données}

%-----------------------------------------------------------

% ROLES ET METIERS

\newacronym{aOS}{OS}{Officier de sécurité}
\newacronym{aRSSI}{RSSI}{Responsable de la sécurité des systèmes d'information}
\newacronym{aDPO}{DPO}{Data Protection Officier}
\newacronym{aCIL}{CIL}{ Correspondant Informatique et Liberté}
\newacronym{aSSI}{SSI}{Sécurité des Systèmes d’information}
\newacronym{aGRC}{GRC}{Gouvernance Risques et Conformité}
\newacronym{aDTF}{DTF}{Defend The Flag}
\newacronym{aCTF}{CTF}{Capture The Flag}
\newacronym{aOIV}{OIV}{Opérateur d'Infrastructures Vitales}
\newacronym{aOSE}{OSE}{Opérateur de Services Essentiels}
\newacronym{aLPM}{LPM}{Loi de Programmation Militaire}
\newacronym{aEBIOS}{EBIOS}{Expression des Besoins et Identification des Objectifs de Sécurité}
\newacronym{aDLL}{DLL}{Dynamic Loads Library}
\newacronym{aACID}{ACID}{Authentification, Confidentialité, Intégrité, Disponibilité}

%-------------------------------------------------------
% Sources/EDU/acronymes DATABASE XLS Generation (cf. RESSOURCES DIR)
%-------------------------------------------------------
% DO NOT MODIFY DIRECTLY, USE XLS File (BEGIN)

\newacronym{a2FA}{2FA}{\textit{Two-Factor Authentication} : Authentification à deux facteurs}
\newacronym{aAAA}{AAA}{\textit{Authentication, Authorization, Accounting} : Authentification, autorisation, comptabilisation}
\newacronym{aAC}{AC}{\textit{Access Control} : Accès conditionnel}
\newacronym{aAES}{AES}{\textit{Advanced Encryption Standard} : Algorithme de chiffrement avancé}
\newacronym{aAPT}{APT}{\textit{Advanced Persistent Threat Group} : Groupe de cyber-espionnage avancé}
\newacronym{aASLR}{ASLR}{\textit{Address Space Layout Randomization} : Randomisation de l'espace d'adressage}
\newacronym{aAV}{AV}{\textit{Antivirus Software} : Logiciel antivirus}
\newacronym{aBOT}{BOT}{\textit{Robot (Malicious)} : Robot (malveillant)}
\newacronym{aCC}{C\&C}{\textit{Command And Control} : Commande et contrôle}
\newacronym{aCSRF}{CSRF}{\textit{Cross-Site Request Forgery Attack} : Attaque de confiance de site à site}
\newacronym{aDDOS}{DDOS}{\textit{Distributed Denial Of Service Attack} : Attaque de déni de service distribué}
\newacronym{aDES}{DES}{\textit{Data Encryption Standard} : Algorithme de chiffrement à clé de 64 bits}
\newacronym{aDLP}{DLP}{\textit{Data Loss Prevention} : Gestion de la protection des données sensibles}
\newacronym{aEDR}{EDR}{\textit{Endpoint Detection And Response} : Détection et réponse aux comportements anormaux d'un terminal}
\newacronym{aEK}{EK}{\textit{Exploit Kit} : Kit d'exploit}
\newacronym{aFTP}{FTP}{\textit{File Transfer Protocol} : Protocole de transfert de fichier}
\newacronym{aHTTPS}{HTTPS}{\textit{Protocol For Securing Communications On The Internet} : Protocole de sécurisation HTTP}
\newacronym{aIAM}{IAM}{\textit{Identity And Access Management} : Identité et Gestion des Accés}
\newacronym{aIDS}{IDS}{\textit{Intrusion Detection System} : Système de détection d'intrusion}
\newacronym{aIIS}{IIS}{\textit{Internet Information Services Server}}
\newacronym{aIPS}{IPS}{\textit{Intrusion Prevention System} : Système de prévention d'incident}
\newacronym{aISA}{ISA}{\textit{Internet Security And Acceleration Server} : Serveur de sécurité et d'accélération Internet (Microsoft)}
\newacronym{aLDAP}{LDAP}{\textit{Lightweight Directory Access Protocol} : Protocole de gestion d'annuaire léger}
\newacronym{aMD5}{MD5}{\textit{Message-Digest Algorithm 5}}
\newacronym{aMFA}{MFA}{\textit{Multi-Factor Authentication} : Authentification multi-facteurs}
\newacronym{aNAC}{NAC}{\textit{Network Access Control}}
\newacronym{aOAuth}{OAuth}{\textit{Open Authorization Protocol}}
\newacronym{aPAM}{PAM}{\textit{Privilege Access Management}}
\newacronym{aPKI}{PKI}{\textit{Public Key Infrastructure} : Infrastructure de clefs publiques}
\newacronym{aRSA}{RSA}{\textit{Rivest-Shamir-Adleman Encryption Algorithm}}
\newacronym{aSAML}{SAML}{\textit{Security Assertion Markup Language}}
\newacronym{aSMS}{SMS}{\textit{Short Message Service}}
\newacronym{aSQLI}{SQLI}{\textit{Sql Injection} : Injection de script de commande SQL}
\newacronym{aSSH}{SSH}{\textit{Secure Shell Protocol}}
\newacronym{aSSL}{SSL}{\textit{Protocol For Securing Communications On The Internet} : Protocole de sécurisation des communications sur Internet}
\newacronym{aTKIP}{TKIP}{\textit{Temporal Key Integrity Protocol}}
\newacronym{aTLS}{TLS}{\textit{Transport Layer Security}}
\newacronym{aURL}{URL}{\textit{Uniform Resource Locator}}
\newacronym{aVPN}{VPN}{\textit{Virtual Private Network} : Réseau privé virtuel}
\newacronym{aWAF}{WAF}{\textit{Web Application Firewall}}
\newacronym{aWPA}{WPA}{\textit{Wifi Protected Access}}
\newacronym{aXSS}{XSS}{\textit{Cross-Site Scripting Attack}}

% DO NOT MODIFY (END)
%-------------------------------------------------------





%-------------------------------------------------
\newpage
\ubody % MAIN BODY defined external
%-------------------------------------------------

	\printbibliography
 \printglossary[type=\acronymtype]
%	\printglossary
%-------------------------------------------------	
\newpage
\Ucontribute % Contribution TEXT
%-------------------------------------------------


%-------------------------------------------------
	\newpage
	\tableofcontents  % TOC
	\iftotalfigures
  			\listoffigures
	\fi
%-------------------------------------------------

\end{document}
}
	\else
	\providecommand{\umainload}{%-------------------------------------------------------------
%               FR CYBERDEF SECOPS COURSE
%                          ARTICLE MAIN FILE
%                             2020 eduf@ction
%-------------------------------------------------------------

\documentclass[10pt,fleqn,twoside]{\utemplatepath/LessonModel/Article/edxstyle} %edxstyle

\usepackage{beamerarticle}

\usepackage{ae}
\usepackage[french]{babel}
\usepackage[utf8]{inputenc}
\usepackage[T1]{fontenc}
\usepackage{lmodern}




\usepackage{titletoc} % Required for manipulating the table of contents

\input{\utemplatepath/commons.inc/CommonsPictures/\utemplatetheme.theme/colors.inc}

%-------------------------------------------------------------
%               FR CYBERDEF SECOPS COURSE
%                                  EDX PACKAGES
%.    Commons packages for BOOK, NOTES (.doc) 
%                            and BEAMER (.prz)
%                              2020 eduf@ction
%-------------------------------------------------------------

%\usepackage[hyphens,spaces,obeyspaces]{url}
%\usepackage[hyphens,spaces,obeyspaces]{url}



\PassOptionsToPackage{hyphens,spaces,obeyspaces}{url}

\usepackage[utf8]{inputenc}
\usepackage[french]{babel}
\usepackage[T1]{fontenc}
%\usepackage{markdown}



\usepackage{ccicons}

\usepackage{tikz} % Required for drawing custom shapes
\usetikzlibrary{tikzmark} 
\usetikzlibrary{mindmap,trees, backgrounds}


\usepackage{verbatim}
\usepackage{wrapfig}
\usepackage{xpatch}

\usepackage{hologo}

\usepackage[figure]{totalcount}

\usepackage{fontawesome}

\usepackage[skins]{tcolorbox} % don't forget SKINS option

\usepackage{listings}
\usepackage{upquote}

\lstset{
upquote=true,columns=flexible, basicstyle=\ttfamily,
language=HTML, 
frameround=tttt,
commentstyle=\color{gray},
identifierstyle=\color{blue},
keywordstyle=\color{uMainColor}\bfseries,
xleftmargin=2em,
xrightmargin=2em,
aboveskip=\topsep,
belowskip=\topsep, 
frame=single,
rulecolor=\color{uMainColor},
backgroundcolor=\color{uMainColor!5},
breaklines,
breakindent=1.5em,
showspaces=false,
showstringspaces=false,
showtabs=false,
xleftmargin=2em,
xrightmargin=2em,
aboveskip=1em,
belowskip=1em,
}



\usepackage{graphicx} % Required for including pictures
%\graphicspath{{Pictures/}} % Specifies the directory where pictures are stored

\usepackage{booktabs} % Required for nicer horizontal rules in tables

\usepackage{setspace}
\setstretch{1,2}

\usepackage{csquotes}

\usepackage{fancyhdr}

\usepackage{makeidx} 

\usepackage{xspace} 

\usepackage{floatflt} 

\usepackage{hyperref}
\usepackage{cleveref}

%pdfdisplaydoctitle={\utitle}, Attention le titre ne doit pas contenir de saut de ligne 

\hypersetup{
pdftitle={\uCoursetittle},
pdfauthor={\uauthor},
hidelinks,
backref=true,
pagebackref=true,
hyperindex=true,
pdfcenterwindow = true,
colorlinks=true,
breaklinks=true,
bookmarks=true,
allcolors=uMainColor,
pdfsubject = {\udescription},
pdfcreator = {MACTEX on TEXPAD},
bookmarksopen=true}


% urlcolor=cnam,
% runcolor=uMainColor,
%\pdfstringdefDisableCommands{%
%  \def\\{}%
%  \def\texttt#1{<#1>}%
%}

\usepackage[style=numeric,citestyle=numeric,sorting=nyt,sortcites=true,autopunct=true,hyperref=true,abbreviate=false,backref=true,backend=biber]{biblatex}

%\usepackage{cleveref} old position

\usepackage{multicol}

%--------------------------------------------------------
% pour la génération du QRCODE en page de garde d'article
\usepackage{qrcode} 

%--------------------------------------------------------
% POur la date d'impression en premiere page d'article
\usepackage[useregional]{datetime2}

%--------------------------------------------------------
% Gestion des glossaires et Sources/EDU/acronymes
% Un fichier acroymes.tex centralisé
\usepackage[acronym]{glossaries}

%\makeglossaries
\makenoidxglossaries 

\loadglsentries{Latex/Sources/commons/acronymes.tex}
\loadglsentries{Latex/Sources/commons/glossaire.tex}


%\renewcommand{\bibname}{R\'ef\'erences}

%--------------------------------------------------------
% Pour exclure des "environnement", ici techworkbox par exemple, travaux pour étudiants, non affiché dans le BOOK.
\usepackage{comment}


%--------------------------------------------------------
% pour filigramme sur les pages
\usepackage{eso-pic}

%--------------------------------------------------------
% for \gappto{\UrlBreaks}{\UrlOrds}
% rather than 
%\makeatletter
% \g@addto@macro{\UrlBreaks}{\UrlOrds}
% \makeatother

\usepackage{etoolbox}

\usepackage{breakurl}


\newcommand{\ucite}[1]{\cite{#1}}
%\newcommand{\ucite}[1]{REF BIB #1}
\newcommand{\ugls}[1]{\gls{#1}}
%\newcommand{\ugls}[1]{GLOSSAIRE (#1)}
%\usepackage[Rejne]{fncychap}

%\usepackage{pythontex} % need shell escape

%
%\makeatletter
%\DeclareRobustCommand{\afilename}[1]{%
% \begingroup
%  % \lstname seems to change hyphens into \textendash
%  \def\textendash{-}%
%  \afilename@parse{#1}%
%  \edef\afilename@base{\detokenize\expandafter{\afilename@base}}%
%  \texttt{\afilename@base.\afilename@ext}%
% \endgroup
%}
%\makeatother

%-------------------------------------------------------------
%               FR CYBERDEF SECOPS COURSE
%                                  EDX PACKAGES
%.    Commons COMMANDS for BOOK, NOTES (.doc) 
%                            and BEAMER (.prz)
%                              2020 eduf@ction
%-------------------------------------------------------------

\newcommand{\ucite}[1]{\cite{#1}}
%\newcommand{\ucite}[1]{REF BIB #1}
\newcommand{\ugls}[1]{\gls{#1}}
%\newcommand{\ugls}[1]{GLOSSAIRE (#1)}
\newcommand{\uac}[1]{\acrshort{#1}}



% \newcommand{\ExtractBeforeLastSlash}[1]{\directlua{extract_before_last_slash("#1")}}
% \newcommand{\ExtractAfterLastSlash}[1]{\directlua{extract_after_last_slash("#1")}}
% \newcommand{\RemoveExtension}[1]{\directlua{remove_extension("#1")}}
% \newcommand{\ExtractExtension}[1]{\directlua{extract_extension("#1")}}


\newcommand{\includepicturesubpath}[1]
{#1}

%{\ExtractBeforeLastSlash{#1}/\upicturethemepath\RemoveExtension{\ExtractAfterLastSlash{#1}}.\ExtractExtension{#1}}
    

%---------------------------------------------------
%										 	AUDIO NOTE 
%---------------------------------------------------

\newcommand\audionote[2]{%
  \note{#2}
  \immediate\write\audioexport{==AudioSlide=============}
  \immediate\write\audioexport{#1}
  \immediate\write\audioexport{\insertframenumber}
  \immediate\write\audioexport{---------------------------------}
  \immediate\write\audioexport{#2}
  \immediate\write\audioexport{---------------------------------}
}%

%---------------------------------------------------
%										 	TEXT STYLE SHORTCUT 
%---------------------------------------------------

\newcommand{\tcb}[1]{\textbf{\textcolor{uMainColor}{#1}}}

\newcommand{\voir}[1]{(Schéma~\ref{#1},~page~\pageref{#1}}

\newcommand{\graphicxextension}{pdf}

\newcommand\myemptypage{
   
    \thispagestyle{empty}
    \addtocounter{page}{-1}
    \newpage
  \null
    %CECI EST UNE PAGE VIDE
       \vfill
    }


\newcommand{\uNA}[2]
{
\newacronym{a#1}{#1}{#2}
}

\newcommand{\coauthor}[4]
{
\noindent (#2) \textbf{#1} - #4 : \textit{#3}
}


\newcommand\UKword[1]{\emph{#1}\xspace}
\newcommand\textUK[1]{\textit{\textbf{#1}}\xspace}
\newcommand\Ucenter[1]{\begin{center}{#1}\end{center}}

\newcommand\head[1]{\textbf{#1}}
\newcommand\tb[1]{\textbf{#1}}
\newcommand{\g}[1]{\og #1 \fg{}}
\newcommand{\uload}[1] {\input{../Latex/#1}}
\newcommand{\ucurcolor}{black}

\newcommand{\mak}[1]{\faicon{\aArrowCircleORight} #1}


%-----------------------------------------------------
%										 	\INFOGitHub 
%---------------------------------------------------


\newcommand{\GITfilename}{https://github.com/edufaction/CYBERDEF101/raw/main/Release/Builder/\jobname.pdf}

%\newcommand{\safeqrcode}[2][]{%
%\qrcode[#1]{\detokenize{#2}} }

\newcommand{\safeqrcode}[2][]{\qrcode[#1]{#2}} 

\newcommand{\INFOGitHub}
{
\begin{center} 
\setstretch{1.5}
 {Vérifiez la disponibilité d'une version plus récente de} \\
 {\ulink{\GITfilename}{\textbf{\jobname.pdf }sur GITHUB CYBERDEF} }   \\ 
{{\huge\ccbyncndeu}}  \\  
{Publication en \href{https://creativecommons.org/licenses/by-nc-nd/2.0/fr/}{Creative Common BY-NC-ND} by eduf@ction}    \\  % Copyright notice
{\safeqrcode[height=4 cm,padding]{\GITfilename}}  
\end{center} 
}

\newcommand{\INFOGitHubS}
{
\noindent{\safeqrcode[height=3 cm]{\GITfilename}} \\
\noindent {\tiny Téléchargez une version à jour} 
}



%-----------------------------------------------------
%										 	\upspicture 
%---------------------------------------------------
\newcommand{\upspicture}[3]
{
\renewcommand{\ucurcolor}{#3}
\resizebox{#2\textwidth}{!}{\input{../#1}}
}


%-----------------------------------------------------
%										 	\uindex 
%---------------------------------------------------


\newcommand{\uindex}[1]
{
#1\index{#1}\xspace
}

\newcommand{\ulindex}[1]{\index{#1}}

%-----------------------------------------------------
%										 	\includer 
%---------------------------------------------------

\newcommand{\includer}[1]
{
\input{../Latex/Sources/EDU/SRC1/Chapters/#1}
}

%-----------------------------------------------------
%										 	\edxdico 
%---------------------------------------------------

\newcommand{\edxdico}[2] {\textbf{#1}:\textit{#2} \index{#1}}

%-----------------------------------------------------
%  command DEF : InTextx
%-----------------------------------------------------
\newcommand\InTexta{no text}
\newcommand\InTextb{no text}
\newcommand\InTextc{no text}
\newcommand\rnc{\renewcommand}


%-----------------------------------------------------
%										 	ENV warningbox
%---------------------------------------------------

\newtcolorbox{warningbox}[2][]
{
  colframe = uMainColor!25,
  colback  = black!05,
  coltitle = uMainColor!20!black,
  fonttitle = \bfseries, 
  title    = #2,
  #1,
}


\newtcolorbox{techworkbox}[2][]{sharp corners,colframe = uMainColor!25,colback  = black!05,coltitle = uMainColor!20!black,fonttitle = \bfseries, title={\textcolor{uMainColor}{\faFile~\small {Sujet de mémoire SECOPS}}} : #2,#1,}


\newtcolorbox{notebox2}[2][]
{
sharp corners, 
colback = uMainColor!5!white, 
colframe = uMainColor!75!black,
fonttitle = \bfseries, 
colbacktitle= uMainColor!85!black,
title=#2,#1
}

\newtcolorbox{notebox}[2][]{colback=uMainColor!5!white,
colframe=uMainColor!75!black,fonttitle=\bfseries,
colbacktitle=uMainColor!85!black,enhanced,
attach boxed title to top right={yshift=+1mm},
title=#2,#1}


\newtcolorbox{toolsboxin}[2][]{sharp corners, colback=uMainColor!0!white,
colframe= black!25,fonttitle=\bfseries,
colbacktitle=black!60,enhanced,
attach boxed title to top left={yshift=+1mm},
title=#2,#1}


\newtcolorbox{techworkin}[2][]{sharp corners, colback=uMainColor!5!white,
colframe= uMainColor!100, fonttitle=\bfseries,
colbacktitle=black!5, enhanced,
 coltitle=black!50,
attach boxed title to top right ={yshift=+1mm},
title=#2,#1}



%-----------------------------------------------------
%										 	ENV ToolsBox
%---------------------------------------------------


\newcommand\tooleditor{No EDITOR in Datas.tex}
\newcommand\toolurl{No URL INSERTED in Datas.tex}
\newcommand\toolclass{No EDITOR in Datas.tex}
\newcommand\toolname{No NAME in Datas.tex}
\newcommand\toolanalyst{No ANALYST in Datas.tex}

\newcommand{\toolsbox}[2]
{
\input{../Latex/Sources/EDU/SRC1/SecComponents/#2.#1/datas.tex} 
\begin{toolsboxin}[sidebyside,righthand width=.7\textwidth]{#1}
\includegraphics[width=0.9\textwidth]{../Latex/Sources/EDU/SRC1/SecComponents/#2.#1/icon.png} \\
\tcblower
\input{../Latex/Sources/EDU/SRC1/SecComponents/#2.#1/description.tex}
\end{toolsboxin}
\noindent \small {\faCogs~Classe : \uppercase{\textbf{#2}}, Site de référence : \textbf{\ulink{\toolurl}{#1}}}
\newline
\small {\faGlobe~Editeur : \textbf{\tooleditor}}
\small {\faEye~Analyste : \textbf{\toolanalyst}}
}

%\renewcommand\toolsclass{SIEM}
%\renewcommand\toolsname{RSA Netwitness}
%\renewcommand\toolseditor{RSA}
%\renewcommand\toolsurl{https://www.rsa.com/fr-fr/products/threat-detection-response}

%-----------------------------------------------------
%										 	ENV  TechWork
%---------------------------------------------------


\newcommand{\techwork}[2]
{
\begin{techworkin} [righthand width=.8\textwidth]{\faCogs~\small {Fiche TECHNO} : #1}#2
\end{techworkin}
}

%-----------------------------------------------------
%										 	\uchap
%---------------------------------------------------
\ifdefined\INTERTITLE
\providecommand{\uchap}[1]
	{
	\begin{LARGE}
	\textbf{\Ucenter{#1 \\ -oOo-}}	
	\end{LARGE}
	}
\else
	\providecommand{\uchap}[1]{}
\fi

%-----------------------------------------------------
%											\uexpand
%---------------------------------------------------
\newcommand {\uexpand}[1]{
\mode<all>\input{../Latex/Sources/EDU/SRC1/Chapters/#1}
}


%-----------------------------------------------------
%											\spicture 
% direct include graphic with Scrivener
%-----------------------------------------------------

\newcommand{\spicture}[3]{
\begin{figure}[!hbtp] %h !hbtp
  \begin{center}
	 \includegraphics[width=0.8\textwidth]{#1}
  \end{center}
\caption{#3 - (slbl-#2)\label{slbl-#2}}
\end{figure}
}

%\StrBefore{#2}{.pdf}

%-----------------------------------------------------
%											\upicture 
%---------------------------------------------------

\ifdefined\PRZMODE

\providecommand{\upicture}[4]{
\framesubtitle{#2}
%\begin{figure}[!h] %h
  \begin{center}
	 \includegraphics[height=#3\textheight, width=#3\textwidth, keepaspectratio ]{\includepicturesubpath{#1.pdf}}
  \end{center}
%  \end{figure}
}

\else

\providecommand{\upicture}[4]{
\begin{figure}[!hbtp] %h !hbtp
  \begin{center}
	 \includegraphics[width=#3\textwidth]{\includepicturesubpath{#1.pdf}}
  \end{center}
\caption{#2 \label{#4}}
\end{figure}
}
	
\fi


\ifdefined\PRZMODE

\providecommand{\uwpicture}[4]{
\framesubtitle{#2}
%\begin{figure}[!h] %h
  \begin{center}
	 \includegraphics[height=0.8\textheight, width=0.8\textwidth, keepaspectratio ]{\includepicturesubpath{#1.pdf}}
  \end{center}
%  \end{figure}
}
\else
\providecommand{\uwpicture}[4]
{
\begin{wrapfigure}{O}{#3\textwidth} % inside / Outside
  \begin{center}
	 \includegraphics[width=#3\textwidth]{\includepicturesubpath{#1.pdf}}
\end{center}
\caption{#2 \label{#4}}
\end{wrapfigure}
}
\fi

%-----------------------------------------------------
%	\picframe for BEAMER (Image included in Beamer Slide)
%---------------------------------------------------

\ifdefined\PRZMODE

\providecommand{\picframe}[4]
{
\begin{frame}
\frametitle<presentation>{#2}
  \begin{center}

   \begin{tikzpicture}[remember picture, overlay]
    \node at (current page.center) {\includegraphics[width=#3\paperwidth, height=#3\paperheight,keepaspectratio]{\includepicturesubpath{#1.pdf}}};
    \end{tikzpicture}

  \end{center}
\end{frame}

}
\else
\providecommand{\picframe}[4]
{
\begin{figure}[!hbtp] % !hbtp ou !h 
\centering 
%{\includegraphics[width=#3\paperwidth, height=#3\paperheight,keepaspectratio]{\includepicturesubpath{#1.pdf}}}
\includegraphics[width=#3\textwidth]{\includepicturesubpath{#1.pdf}}
\caption{#2 \label{#4}}
\end{figure}
}
\fi

%-----------------------------------------------------
%.  for BEAMER (Image included in Beamer Slide)
%	\picframeS Simplified with Name and Label from One Param
%---------------------------------------------------


\ifdefined\PRZMODE

\providecommand{\picframeS}[2]
{
\begin{frame}
\frametitle<presentation>{#2}
  \begin{center}
	 \includegraphics[height=0.95\textheight, width=0.95\textwidth, keepaspectratio ]{\picpath/img-#1.\graphicxextension}
  \end{center}
\end{frame}

}
\else
\providecommand{\picframeS}[2]
{
\begin{figure}[!hbtp] % !hbtp ou !h 
\centering \includegraphics[width=0.9\textwidth]{\picpath/img-#1.\graphicxextension}
\caption{(s) #2 \label{lbl:#1}}
\end{figure}
}

\fi


%{figure}[!h]
%{wrapfigure}{O}{#3\textwidth}
%-----------------------------------------------------
%	FRAME			\wpicframe 
%---------------------------------------------------

\ifdefined\PRZMODE

\providecommand{\wpicframe}[4]
{
\begin{frame}
\frametitle<presentation>{#2}
  \begin{center}
	 \includegraphics[height=0.95\textheight, width=0.95\textwidth, keepaspectratio ]{#1.\graphicxextension}
  \end{center}
\end{frame}

}
\else
\providecommand{\wpicframe}[4]
{
\begin{wrapfigure}{O}{#3\textwidth}
\centering \includegraphics[width=#3\textwidth]{#1.\graphicxextension}
\caption{#2 \label{#4}}
\end{wrapfigure}
}
\fi

%-----------------------------------------------------
%	FRAME										\texframe 
%---------------------------------------------------
\ifdefined\PRZMODE

\providecommand{\texframe}[3]
{
\begin{frame}
\frametitle<presentation>{#1}
\framesubtitle<presentation>{#2}
#3
\end{frame}
}
\else
\providecommand{\texframe}[3]
{
#3
}

\fi

%-----------------------------------------------------
%											\uref 
%---------------------------------------------------
%\newcommand{\uref}[2]{(Voir~#1~\cref{#2} page~\pageref{#2})}

\newcommand{\uref}[2]{(Voir~#1~\ref{#2} page~\pageref{#2})}
%-----------------------------------------------------
%											\ubg
%---------------------------------------------------
\newcommand{\ubg}[1]{
\textbf{\g{\uppercase{#1}}}
}
%-----------------------------------------------------
%											\updfimage 
%---------------------------------------------------

\newcommand{\updfimage}[3]
{
\begin{wrapfigure}{R}{#3\textwidth}
  \begin{center}
	 \resizebox{#3\textwidth}{!}{\includegraphics{\includepicturesubpath{#1.pdf}}}
  \end{center}
\caption{#2}
\end{wrapfigure}
}


%************************PACKAGE***

%\usepackage{tikz}% Required for drawing custom shapes
\usetikzlibrary{shadows}


%-----------------------------------------------------
%											 		\rem 
% insertion d'une remarque avec R
%---------------------------------------------------
\newcommand\rem[1]{%
   \marginpar{%
   \tikzpicture[baseline={(title.base)}]
      \node[inner sep=5pt,text width=4cm,drop shadow={shadow yshift=-5pt,shadow xshift=5pt,uMainColor},fill=white] (box) {\vskip5pt \nointerlineskip #1};
      \node[right=10pt,inner sep=0pt,outer sep=10pt] (title) at (box.north west) {\bfseries\color{uMainColor}Remarque};
      \draw[draw=uMainColor,very thick](title.west)--(box.north west)--(box.south west)--(box.south east)--(box.north east)--(title.east);
      \fill[uMainColor]([yshift=-10pt]box.north west)--+(-5pt,-5pt)--+(0pt,-10pt);
   \endtikzpicture}%
}
%-----------------------------------------------------
%											 \utikzimage 
% insertion d'une image sur fichier tkz.tex
%---------------------------------------------------

\newcommand{\utikzimage}[3]
{
\begin{wrapfigure}{R}{#3\textwidth}
  \begin{center}
	 \resizebox{#3\textwidth}{!}{\input{../#1.tkz.tex}}
  \end{center}
\caption{#2}
\end{wrapfigure}
}

%-----------------------------------------------------
%														\link 
% Url et footnote du lien
%-----------------------------------------------------

% TT every character and hyphenate after it
\def\hyphenateAndTtWholeString #1{\xHyphenate#1$\wholeString\unskip}

\def\xHyphenate#1#2\wholeString {\if#1$%
    \else\transform{#1}%
    \takeTheRest#2\ofTheString\fi}

\def\takeTheRest#1\ofTheString\fi
{\fi \xHyphenate#1\wholeString}

\def\transform#1{\url{#1}\hskip 0pt plus 1pt}

% Define the \urlx command which works like \url, but with line brakes
\def\urlx #1{\href{#1}{\hyphenateAndTtWholeString{#1}}}
  

 %-----------------------------------------------------
%														\link 
%-----------------------------------------------------  
  
  \newcommand{\ulink}[2]
  {\href{#1}
  {#2~\raisebox{-0.2ex}{\faExternalLink}\footnote{\url{#1}}}\xspace}
  
  \newcommand{\wikipedia}[1]
  { \href{https://fr.wikipedia.org/wiki/#1}{\textcolor{grey}{#1}}\raisebox{0.5ex}{\tiny{\textcolor{uMainColor}{\faWikipediaW}}}}
  


 %-----------------------------------------------------
%														\utodo 
%-----------------------------------------------------

  \newcommand{\utodo}{\begin{center}{\noindent \textcolor{cyan}{\faGears~{En cours de rédaction, DRAFT non publiable} \faGears}}\end{center}}
  
    \newcommand{\utocomplete}{\begin{center}{\noindent \textcolor{blue}{\faClipboard~{à compléter, les éléments ne donnent qu'une vue trop réduite ou trop parcellaire du sujet}}}\end{center}}
  

 




%-------------------------------------------------------------
%               FR CYBERDEF SECOPS COURSE
%
%                                    EDX STYLE
%
%                              2020 eduf@ction
%-------------------------------------------------------------
\usepackage{enumitem}

\usepackage{parskip}
\setlength{\parindent}{0em}
\setlength{\parskip}{5pt} % 1ex plus 0.5ex minus 0.2ex}

\def\uCoverColor{}

%---------------------------------------------------------------
%	MARGINS
%---------------------------------------------------------------

\usepackage{geometry}

\geometry{
	paper=a4paper, 
%	top=1cm, 
%	bottom=1cm, 
%	left=1cm, 
%	right=1cm, 
%	headheight=20pt, 
%	footskip=1.4cm, 
%	headsep=10pt, 
%	showframe, 
	showcrop 
}

%---------------------------------------------------------------
%	FOOTNOTE (format)
%---------------------------------------------------------------

\makeatletter
\long\def\@makefntextFB#1{%
    \ifx\thefootnote\ftnISsymbol
        \@makefntextORI{#1}%
    \else
        \rule\z@\footnotesep
        \setbox\@tempboxa\hbox{\@thefnmark}%
            \ifdim\wd\@tempboxa>\z@
                \kern2em\llap{\@thefnmark.\kern0.5em}%
            \fi
        \hangindent2em\hangafter\@ne#1
    \fi}
\makeatother

%-----------------------------------------------------
%	COLOR
%-----------------------------------------------------
%\usepackage{xcolor} % Required for specifying colors by name
%\definecolor{uMainColor}{RGB}{160,0,0} 
\colorlet{edxcolorcover}{uCoverColor}
%\definecolor{edxcolorcover}{RGB}{160,0,0} 

%\definecolor{cnam}{RGB}{128, 0, 32} 
                 

  \definecolor{comments}{rgb}{0.7,0,0}    % rouge foncé
  \definecolor{link}{rgb}{0,0.4,0.6}      % ~RoyalBlue de dvips
  \definecolor{url}{rgb}{0.6,0,0}         % rouge-brun
  \definecolor{citation}{rgb}{0,0.5,0}    % vert foncé
  \definecolor{ULlinkcolor}{rgb}{0,0,0.3} % de ulthese.cls
  \definecolor{rouge}{rgb}{0.85,0,0.07}   % rouge bandeau identitaire
  \definecolor{or}{rgb}{1,0.8,0}          % or bandeau identitaire

%-----------------------------------------------------
%	FONTS
%-----------------------------------------------------

%%\usepackage{avant} % Use the Avantgarde font for headings
%\usepackage{times} % Use the Times font for headings
%\usepackage{mathptmx} % Use the Adobe Times Rgman as the default text font together with math symbols from the Sym­bol, Chancery and Com­puter Modern fonts

%%\usepackage{microtype} % Slightly tweak font spacing for aesthetics
%%\usepackage[utf8]{inputenc} % Required for including letters with accents
%\usepackage[T1]{fontenc} % Use 8-bit encoding that has 256 glyphs

% Using AVANT GARDE Family

%\renewcommand{\familydefault}{\sfdefault}

%\fontfamily{pag}\selectfont
%\renewcommand{\familydefault}{pag} % sfdefault pag lmss

%bch         Charter
%lmr         Latin Modern Roman
%lmss        Latin Modern Sans Serif
%lmssq       Latin Modern Sans Serif extended
%lmtt        Latin Modern Typewriter
%lmvtt       Latin Modern Typewriter proportional
%pag         Avant Garde
%pbk         Bookman
%pcr         Courier
%phv         Helvetica
%pnc         New Century Schoolbook
%ppl         Palatino
%ptm         Times
%put         Utopia
 
%\renewcommand{\rmdefault}{pag} % text 
%\renewcommand{\sfdefault}{pag} % titre


%\renewcommand{\rmdefault}{pag} % text 
%\renewcommand{\sfdefault}{phv} % titre


%-----------------------------------------------------
%	BIBLIOGRAPHY AND INDEX
%-----------------------------------------------------


\usepackage{calc} % For simpler calculation - used for spacing the index letter headings correctly

\usepackage{imakeidx}
%\makeindex % Tells LaTeX to create the files required for indexing

\makeindex[columns=2,intoc=true,options={-s \upath/indexstyle.ist}]

%-----------------------------------------------------
%	HEADERS AND FOOTERS
%-----------------------------------------------------


\usepackage{fancyhdr} % Required for header and footer configuration
\pagestyle{fancy} % Enable the custom headers and footers

\ifcsname chapter \endcsname   % verify if command name exists in the CLASS (Book, report, article ...)
% BEGIN IF
\renewcommand{\chaptermark}[1]{\markboth{\sffamily\normalsize\bfseries\chaptername\ \thechapter.\ #1}{}} % Styling for the current chapter in the header
  \else
%nothing
\fi
% END IF

\renewcommand{\sectionmark}[1]{\markright{\sffamily\normalsize\thesection\hspace{5pt}#1}{}} % Styling for the current section in the header

\fancyhf{} % Clear default headers and footers
\fancyhead[LE,RO]{\sffamily\normalsize\thepage} % Styling for the page number in the header
\fancyhead[LO]{\sffamily\normalsize\rightmark} % Print the nearest section name on the left side of odd pages
\fancyhead[RE]{\sffamily\normalsize\leftmark} % Print the current chapter name on the right side of even pages
\fancyfoot[RE,LO]{\sffamily\normalsize\uCoursetittle} % Uncomment to include a footer
\fancyfoot[LE,RO]{\hyperlink{toc}{\sffamily\normalsize\uMinilogo}} % Uncomment to include a footer
% \hyperlink{toc}
\renewcommand{\headrulewidth}{1.0pt} % Thickness of the rule under the header

\fancypagestyle{plain}{% Style for when a plain pagestyle is specified
	\fancyhead{}\renewcommand{\headrulewidth}{0pt}} 

% Removes the header from odd empty pages at the end of chapters
\makeatletter
\renewcommand{\cleardoublepage}{
\clearpage\ifodd\c@page\else
\hbox{}
\vspace*{\fill}
\thispagestyle{empty}
\newpage
\fi}


%----------------------------------------------------------------------------------------
%	MAIN TABLE OF CONTENTS
%----------------------------------------------------------------------------------------
%\usepackage{titletoc} % Required for manipulating the table of contents

\contentsmargin{0cm} % Removes the default margin

% Part text styling (this is mostly taken care of in the PART HEADINGS section of this file)
\titlecontents{part}
	[0cm] % Left indentation
	{\addvspace{20pt}\bfseries} % Spacing and font options for parts
	{}
	{}
	{}

% Chapter text styling
\titlecontents{chapter}
	[1.25cm] % Left indentation
	{\addvspace{12pt}\large\sffamily\bfseries} % Spacing and font options for chapters
	{\color{uMainColor!60}\contentslabel[\Large\thecontentslabel]{1.25cm}\color{uMainColor}} % Formatting of numbered sections of this type
	{\color{uMainColor}} % Formatting of numberless sections of this type
	{\color{uMainColor!60}\normalsize\;\titlerule*[.5pc]{.}\;\thecontentspage} % Formatting of the filler to the right of the heading and the page number

% Section text styling
\titlecontents{section}
	[1.25cm] % Left indentation
	{\addvspace{3pt}\sffamily\bfseries} % Spacing and font options for sections
	{\contentslabel[\thecontentslabel]{1.25cm}} % Formatting of numbered sections of this type
	{} % Formatting of numberless sections of this type
	{\hfill\color{black}\thecontentspage} % Formatting of the filler to the right of the heading and the page number

% Subsection text styling
\titlecontents{subsection}
	[1.25cm] % Left indentation
	{\addvspace{1pt}\sffamily\small} % Spacing and font options for subsections
	{\contentslabel[\thecontentslabel]{1.25cm}} % Formatting of numbered sections of this type
	{} % Formatting of numberless sections of this type
	{\ \titlerule*[.5pc]{.}\;\thecontentspage} % Formatting of the filler to the right of the heading and the page number

%\titlecontents{subsubsection}
%	[1.25cm] % Left indentation
%	{\addvspace{1pt}\sffamily\tiny} % Spacing and font options for subsections
%	{\contentslabel[\thecontentslabel]{1.25cm}} % Formatting of numbered sections of this type
%	{} % Formatting of numberless sections of this type
%	{\ \titlerule*[.5pc]{.}\;\thecontentspage} % Formatting of the filler to the right of the heading and the page number

\titlecontents{subsubsection}
[1.25cm]
  {\addvspace{1pt}\sffamily\tiny}
 {\contentslabel[\thecontentslabel]{1.25cm}}
  {}
  {}

% Figure text styling
\titlecontents{figure}
	[1.25cm] % Left indentation
	{\addvspace{1pt}\sffamily\small} % Spacing and font options for figures
	{\thecontentslabel\hspace*{1em}} % Formatting of numbered sections of this type
	{} % Formatting of numberless sections of this type
	{\ \titlerule*[.5pc]{.}\;\thecontentspage} % Formatting of the filler to the right of the heading and the page number

% Table text styling
\titlecontents{table}
	[1.25cm] % Left indentation
	{\addvspace{1pt}\sffamily\small} % Spacing and font options for tables
	{\thecontentslabel\hspace*{1em}} % Formatting of numbered sections of this type
	{} % Formatting of numberless sections of this type
	{\ \titlerule*[.5pc]{.}\;\thecontentspage} % Formatting of the filler to the right of the heading and the page number

%----------------------------------------------------------------------------------------
%	MINI TABLE OF CONTENTS IN PART HEADS
%----------------------------------------------------------------------------------------

% Chapter text styling
\titlecontents{lchapter}
	[0em] % Left indentation
	{\addvspace{15pt}\large\sffamily\bfseries} % Spacing and font options for chapters
	{\color{uMainColor}\contentslabel[\Large\thecontentslabel]{1.25cm}\color{uMainColor}} % Chapter number
	{}  
	{\color{uMainColor}\normalsize\sffamily\bfseries\;\titlerule*[.5pc]{.}\;\thecontentspage} % Page number

% Section text styling
\titlecontents{lsection}
	[0em] % Left indentation
	{\sffamily\small} % Spacing and font options for sections
	{\contentslabel[\thecontentslabel]{1.25cm}} % Section number
	{}
	{}

% Subsection text styling (note these aren't shown by default, display them by searchings this file for tgcdepth and reading the commented text)
%\titlecontents{lsubsection}
%	[.5em] % Left indentation
%	{\sffamily\footnotesize} % Spacing and font options for subsections
%	{\contentslabel[\thecontentslabel]{1.25cm}}
%	{}
%	{}


\usepackage[]{ccicons}


%----------------------------------------------------------------------------------------
%	THEOREM STYLES
%----------------------------------------------------------------------------------------

\usepackage{amsmath,amsfonts,amssymb,amsthm} % For math equations, theorems, symbols, etc

\newcommand{\intoo}[2]{\mathopen{]}#1\,;#2\mathclose{[}}
\newcommand{\ud}{\mathop{\mathrm{{}d}}\mathopen{}}
\newcommand{\intff}[2]{\mathopen{[}#1\,;#2\mathclose{]}}
\renewcommand{\qedsymbol}{$\blacksquare$}

\ifcsname chapter \endcsname   % verify if command name exists in the CLASS (Book, report, article ...)
	\newtheorem{notation}{Notation}[chapter]
 \else
	\newtheorem{notation}{Notation}[section]
\fi


% Boxed/framed environments
\newtheoremstyle{ocrenumbox}% Theorem style name
{0pt}% Space above
{0pt}% Space below
{\normalfont}% Body font
{}% Indent amount
{\small\bf\sffamily\color{uMainColor}}% Theorem head font
{\;}% Punctuation after theorem head
{0.25em}% Space after theorem head
{\small\sffamily\color{uMainColor}\thmname{#1}\nobreakspace\thmnumber{\@ifnotempty{#1}{}\@upn{#2}}% Theorem text (e.g. Theorem 2.1)
\thmnote{\nobreakspace\the\thm@notefont\sffamily\bfseries\color{black}---\nobreakspace#3.}} % Optional theorem note

\newtheoremstyle{ocrenumbox}% Theorem style name
{0pt}% Space above
{0pt}% Space below
{\normalfont}% Body font
{}% Indent amount
{\small\bf\sffamily\color{uMainColor}}% Theorem head font
{\;}% Punctuation after theorem head
{0.25em}% Space after theorem head
{\small\sffamily\color{uMainColor}\thmname{#1}\nobreakspace\thmnumber{\@ifnotempty{#1}{}\@upn{#2}}% Theorem text (e.g. Theorem 2.1)
\thmnote{\nobreakspace\the\thm@notefont\sffamily\bfseries\color{black}---\nobreakspace#3.}} % Optional theorem note

\newtheoremstyle{blacknumex}% Theorem style name
{5pt}% Space above
{5pt}% Space below
{\normalfont}% Body font
{} % Indent amount
{\small\bf\sffamily}% Theorem head font
{\;}% Punctuation after theorem head
{0.25em}% Space after theorem head
{\small\sffamily{\tiny\ensuremath{\blacksquare}}\nobreakspace\thmname{#1}\nobreakspace\thmnumber{\@ifnotempty{#1}{}\@upn{#2}}% Theorem text (e.g. Theorem 2.1)
\thmnote{\nobreakspace\the\thm@notefont\sffamily\bfseries---\nobreakspace#3.}}% Optional theorem note


\newtheoremstyle{blacknumboxS}% name of the style to be used
{5pt}% measure of space to leave above the theorem. E.g.: 3pt
{5pt}% measure of space to leave below the theorem. E.g.: 3pt
{\normalfont}% name of font to use in the body of the theorem
{}% measure of space to indent
{\small\bf\sffamily}% name of head font
{}% punctuation between head and body
{ }% space after theorem head; " " = normal interword space
{\color{black}\faEye~  \color{uMainColor}\small\sffamily\thmnote{#3 : }}

% Non-boxed/non-framed environments
\newtheoremstyle{ocrenum}% Theorem style name
{5pt}% Space above
{5pt}% Space below
{\normalfont}% Body font
{}% Indent amount
{\small\bf\sffamily\color{uMainColor}}% Theorem head font
{\;}% Punctuation after theorem head
{0.25em}% Space after theorem head
{\small\sffamily\color{uMainColor}\thmname{#1}\nobreakspace\thmnumber{\@ifnotempty{#1}{}\@upn{#2}}% Theorem text (e.g. Theorem 2.1)
\thmnote{\nobreakspace\the\thm@notefont\sffamily\bfseries\color{black}---\nobreakspace#3.}} % Optional theorem note
\makeatother

% Defines the theorem text style for each type of theorem to one of the three styles above



\ifcsname chapter \endcsname

\newcounter{dummy} 

\numberwithin{dummy}{section}
	\theoremstyle{ocrenumbox}
		\newtheorem{theoremeT}[dummy]{Proposition}
		\newtheorem{problem}{Problem}[chapter]
		\newtheorem{exerciseT}{Outillage}[chapter]
	\theoremstyle{blacknumex}
		\newtheorem{exampleT}{Example}[chapter]
	\theoremstyle{blacknumbox}
		\newtheorem{vocabulary}{Vocabulary}[chapter]
		\newtheorem{definitionT}{Definition}[section]
		\newtheorem{corollaryT}[dummy]{Concept}
	\theoremstyle{blacknumboxS}
		\newtheorem{remarqueST}[]{}
		\newtheorem{remarqueT}[dummy]{Remarque}
	\theoremstyle{ocrenum}
	\newtheorem{proposition}[dummy]{Proposition}
  
 \else
 
\newcounter{dummy} 
\numberwithin{dummy}{section}

	\theoremstyle{ocrenumbox}
		\newtheorem{theoremeT}[dummy]{Proposition}
		%\newtheorem{problem}{Problem}[section]
		\newtheorem{exerciseT}{Outillage}[section]
	\theoremstyle{blacknumex}
		\newtheorem{exampleT}{Example}[section]
	\theoremstyle{blacknumbox}
		\newtheorem{vocabulary}{Vocabulary}[section]
		\newtheorem{definitionT}{Definition}[section]
		\newtheorem{corollaryT}[dummy]{concept}
	\theoremstyle{blacknumboxS}
		\newtheorem{remarqueST}[]{}
		\newtheorem{remarqueT}[dummy]{Remarque}
	\theoremstyle{ocrenum}
		\newtheorem{proposition}[dummy]{Proposition}
\fi



%----------------------------------------------------------------------------------------
%	DEFINITION OF COLORED BOXES
%----------------------------------------------------------------------------------------

\RequirePackage[framemethod=default]{mdframed} % Required for creating the theorem, definition, exercise and corollary boxes

% Theorem box
\newmdenv[skipabove=7pt,
skipbelow=7pt,
backgroundcolor=black!5,
linecolor=uMainColor,
innerleftmargin=5pt,
innerrightmargin=5pt,
innertopmargin=5pt,
leftmargin=0cm,
rightmargin=0cm,
innerbottommargin=5pt]{tBox}

% Exercise box	  
\newmdenv[skipabove=7pt,
skipbelow=7pt,
rightline=false,
leftline=true,
topline=false,
bottomline=false,
backgroundcolor=uMainColor!10,
linecolor=uMainColor,
innerleftmargin=5pt,
innerrightmargin=5pt,
innertopmargin=5pt,
innerbottommargin=5pt,
leftmargin=0cm,
rightmargin=0cm,
linewidth=4pt]{eBox}	

% Definition box
\newmdenv[skipabove=7pt,
skipbelow=7pt,
rightline=false,
leftline=true,
topline=false,
bottomline=false,
linecolor=uMainColor,
innerleftmargin=5pt,
innerrightmargin=5pt,
innertopmargin=0pt,
leftmargin=0cm,
rightmargin=0cm,
linewidth=4pt,
innerbottommargin=0pt]{dBox}	

% Corollary box
\newmdenv[skipabove=7pt,
skipbelow=7pt,
rightline=false,
leftline=true,
topline=false,
bottomline=false,
linecolor=gray,
backgroundcolor=black!5,
innerleftmargin=5pt,
innerrightmargin=5pt,
innertopmargin=5pt,
leftmargin=0cm,
rightmargin=0cm,
linewidth=4pt,
innerbottommargin=5pt]{cBox}

% Creates an environment for each type of theorem and assigns it a theorem text style from the "Theorem Styles" section above and a colored box from above
%\newenvironment{theorem}{\begin{tBox}\begin{theoremeT}}{\end{theoremeT}\end{tBox}}
\newenvironment{exercise}{\begin{eBox}\begin{exerciseT}}{\hfill{\color{uMainColor}\tiny\ensuremath{\blacksquare}}\end{exerciseT}\end{eBox}}				  
%\newenvironment{definition}{\begin{dBox}\begin{definitionT}}{\end{definitionT}\end{dBox}}	
%\newenvironment{example}{\begin{exampleT}}{\hfill{\tiny\ensuremath{\blacksquare}}\end{exampleT}}		
%\newenvironment{corollary}{\begin{cBox}\begin{corollaryT}}{\end{corollaryT}\end{cBox}}	
\newenvironment{nota}{\begin{cBox}\begin{remarqueT}}{\end{remarqueT}\end{cBox}}	

%----------------------------------------------------------------------------------------
%	REMARK ENVIRONMENT
%----------------------------------------------------------------------------------------

\newenvironment{remark}{\par\vspace{10pt}\small % Vertical white space above the remark and smaller font size
\begin{list}{}{
\leftmargin=35pt % Indentation on the left
\rightmargin=25pt}\item\ignorespaces % Indentation on the right
\makebox[-2.5pt]{\begin{tikzpicture}[overlay]
\node[draw=uMainColor!60,line width=1pt,circle,fill=uMainColor!25,font=\sffamily\bfseries,inner sep=2pt,outer sep=0pt] at (-15pt,0pt){\textcolor{uMainColor}{i}};\end{tikzpicture}} % Orange R in a circle
\advance\baselineskip -1pt}{\end{list}\vskip5pt} % Tighter line spacing and white space after remark

%----------------------------------------------------------------------------------------
%	SECTION NUMBERING IN THE MARGIN
%----------------------------------------------------------------------------------------

\makeatletter
% \renewcommand{\@seccntformat}[1]{\textcolor{uMainColor}{\csname the#1\endcsname}\quad}   

%\makeatother

%----------------------------------------------------------------------------------------
%	PART HEADINGS
%----------------------------------------------------------------------------------------

% Numbered part in the table of contents
\newcommand{\@mypartnumtocformat}[2]{%
	\setlength\fboxsep{0pt}%
	\noindent\colorbox{uMainColor!20}{\strut\parbox[c][.7cm]{\ecart}{\color{uMainColor!70}\Large\sffamily\bfseries\centering#1}}\hskip\esp\colorbox{uMainColor!40}{\strut\parbox[c][.7cm]{\linewidth-\ecart-\esp}{\Large\sffamily\centering#2}}%
}

% Unnumbered part in the table of contents
\newcommand{\@myparttocformat}[1]{%
	\setlength\fboxsep{0pt}%
	\noindent\colorbox{uMainColor!40}{\strut\parbox[c][.7cm]{\linewidth}{\Large\sffamily\centering#1}}%
}

\newlength\esp
\setlength\esp{4pt}
\newlength\ecart
\setlength\ecart{1.2cm-\esp}
\newcommand{\thepartimage}{}%
\newcommand{\partimage}[1]{\renewcommand{\thepartimage}{#1}}%
\def\@part[#1]#2{%
\ifnum \c@secnumdepth >-2\relax%
\refstepcounter{part}%
\addcontentsline{toc}{part}{\texorpdfstring{\protect\@mypartnumtocformat{\thepart}{#1}}{\partname~\thepart\ ---\ #1}}
\else%
\addcontentsline{toc}{part}{\texorpdfstring{\protect\@myparttocformat{#1}}{#1}}%
\fi%
\startcontents%
\markboth{}{}%
{\thispagestyle{empty}%
\begin{tikzpicture}[remember picture,overlay]%
\node at (current page.north west){\begin{tikzpicture}[remember picture,overlay]%	
\fill[uMainColor!20](0cm,0cm) rectangle (\paperwidth,-\paperheight);
\node[anchor=north] at (4cm,-3.25cm){\color{uMainColor!40}\fontsize{220}{100}\sffamily\bfseries\thepart}; 
\node[anchor=south east] at (\paperwidth-1cm,-\paperheight+1cm){\parbox[t][][t]{8.5cm}{
\printcontents{l}{0}{\setcounter{tocdepth}{1}}% The depth to which the Part mini table of contents displays headings; 0 for chapters only, 1 for chapters and sections and 2 for chapters, sections and subsections
}};
\node[anchor=north east] at (\paperwidth-1.5cm,-3.25cm){\parbox[t][][t]{15cm}{\strut\raggedleft\color{white}\fontsize{30}{30}\sffamily\bfseries#2}};
\end{tikzpicture}};
\end{tikzpicture}}%
\@endpart}
\def\@spart#1{%
\startcontents%
\phantomsection
{\thispagestyle{empty}%
\begin{tikzpicture}[remember picture,overlay]%
\node at (current page.north west){\begin{tikzpicture}[remember picture,overlay]%	
\fill[uMainColor!20](0cm,0cm) rectangle (\paperwidth,-\paperheight);
\node[anchor=north east] at (\paperwidth-1.5cm,-3.25cm){\parbox[t][][t]{15cm}{\strut\raggedleft\color{white}\fontsize{30}{30}\sffamily\bfseries#1}};
\end{tikzpicture}};
\end{tikzpicture}}
\addcontentsline{toc}{part}{\texorpdfstring{%
\setlength\fboxsep{0pt}%
\noindent\protect\colorbox{uMainColor!40}{\strut\protect\parbox[c][.7cm]{\linewidth}{\Large\sffamily\protect\centering #1\quad\mbox{}}}}{#1}}%
\@endpart}
\def\@endpart{\vfil\newpage
\if@twoside
\if@openright
\null
\thispagestyle{empty}%
\newpage
\fi
\fi
\if@tempswa
\twocolumn
\fi}

%----------------------------------------------------------------------------------------
%	CHAPTER HEADINGS
%-------------------------------------------------%--------------------------------------

%
%
%% A switch to conditionally include a picture, implemented by Christian Hupfer
\newif\ifusechapterimage
\usechapterimagetrue
\newcommand{\thechapterimage}{}%
\newcommand{\chapterimage}[1]{\ifusechapterimage\renewcommand{\thechapterimage}{#1}\fi}%
\newcommand{\autodot}{.}

\def\@makechapterhead#1{%
{\parindent \z@ \raggedright \normalfont
\ifnum \c@secnumdepth >\m@ne
\if@mainmatter
\begin{tikzpicture}[remember picture,overlay]
	
\node at (current page.north west)
			{\begin{tikzpicture}[remember picture,overlay]
			\node[anchor=north west,inner sep=0pt] at (0,0) {\ifusechapterimage\includegraphics[width=\paperwidth]{\thechapterimage}\fi};
			\draw[anchor=west] (\Gm@lmargin,-9cm) node [line width=2pt,rounded corners=15pt,draw=uMainColor,fill=white,fill opacity=1,inner sep=15pt]{\strut\makebox[22cm]{}};
			
			\draw[anchor=west] (\Gm@lmargin+.3cm,-9cm) node {\fontsize{20}{30}\selectfont\sffamily\bfseries\color{uMainColor}~#1\strut};
			%{\Huge\sffamily\bfseries\color{uMainColor}\thechapter\autodot~#1\strut};
			\node at (current page.north east)  [xshift=-3.95cm, yshift=-4.05cm, text opacity=1]  {{\color{white}\centering\fontsize{200}{30}\selectfont \bfseries\sffamily\thechapter\strut} };
			\node at (current page.north east)  [xshift=-4cm, yshift=-4cm, text opacity=1]  {{\color{uMainColor}\centering\fontsize{200}{30}\selectfont \bfseries\sffamily\thechapter\strut} };
			
			\end{tikzpicture}};
	\end{tikzpicture}

\else

\begin{tikzpicture}[remember picture,overlay]
\node at (current page.north west)
{\begin{tikzpicture}[remember picture,overlay]
\node[anchor=north west,inner sep=0pt] at (0,0) {\ifusechapterimage\includegraphics[width=\paperwidth]{\thechapterimage}\fi};
\draw[anchor=west] (\Gm@lmargin,-9cm) node [line width=2pt,rounded corners=15pt,draw=uMainColor,fill=white,fill opacity=1,inner sep=15pt]{\strut\makebox[22cm]{}};

\draw[anchor=east] (\Gm@lmargin,-9cm) node [line width=2pt,rounded corners=15pt,draw=uMainColor,fill=white,fill opacity=1,inner sep=15pt]{\strut\makebox[22cm]{}};
\draw[anchor=west] (\Gm@lmargin+.3cm,-9cm) node {\huge\sffamily\bfseries\color{uMainColor}#1\strut};

\end{tikzpicture}};
 
\end{tikzpicture}
\fi\fi\par\vspace*{270\p@}}}

\makeatother
%-------------------------------------------

%\def\@makeschapterhead#1{%
%\begin{tikzpicture}[remember picture,overlay]
%\node at (current page.north west)
%{\begin{tikzpicture}[remember picture,overlay]
%\node[anchor=north west,inner sep=0pt] at (0,0) {\ifusechapterimage\includegraphics[width=\paperwidth]{\thechapterimage}\fi};
%\draw[anchor=west] (\Gm@lmargin,-9cm) node [line width=2pt,rounded corners=15pt,draw=uMainColor,fill=white, fill opacity=0.8,inner sep=15pt]{\strut\makebox[22cm]{}};
%\draw[anchor=west] (\Gm@lmargin+.3cm,-9cm) node {\huge\sffamily\bfseries\color{uMainColor}#1\strut};
%\end{tikzpicture}};
%\end{tikzpicture}
%\par\vspace*{270\p@}}




%----------------------------------------------------------------------------------------
%	LINKS
%----------------------------------------------------------------------------------------


\usepackage{bookmark}
\bookmarksetup{
open,
numbered,
addtohook={%
\ifnum\bookmarkget{level}=0 % chapter
\bookmarksetup{bold}%
\fi
\ifnum\bookmarkget{level}=-1 % part
\bookmarksetup{color=uMainColor,bold}%
\fi
}
}

%----------------------------------------------------------------------------------------
%	FANCY CHAPTER Header
%----------------------------------------------------------------------------------------





%\definecolor{uMainColor}{RGB}{255,128,0}   



%-------------------------------------------------
%	BIBLIOGRAPHY
%-------------------------------------------------

\addbibresource{../Latex/Sources/commons/bibliography.bib} % BibTeX bibliography file
\defbibheading{bibempty}{}

%-------------------------------------------------
%	COLORS & BORDERS
%-------------------------------------------------
\setlength{\fboxrule}{0.75pt} % Width of the border around the abstract
%\definecolor{uMainColor}{RGB}{255,128,0} % Color of the article title and section
\definecolor{color2}{RGB}{220,220,220} % Color of the boxes behind the abstract and headings

%-------------------------------------------------
% ITEMS DEFINITION
%-------------------------------------------------

\setlist [itemize,1]{label=\color{uMainColor}\faCaretRight }


\usepackage{environ}
\newwrite\audioexport



\begin{document}





  %-------------------------------------------------------
  % AudioNote Open File
  %-------------------------------------------------------
\mode<all>{\immediate\openout\audioexport=\jobname.txt}


\addtocontents{toc}{\protect\hypertarget{toc}{}}

%-------------------------------------------------
%	ARTICLE INFORMATION
%-------------------------------------------------

\Abstract {\uabstract%===========================
% COURS "INTRO CYBERDEF"
% Abstract général des articles
%===========================


Il fait partie du cours introductif aux fondamentaux de la sécurité des systèmes d'information vue sous deux prismes quelques fois opposés dans la littérature : la gouvernance et la gestion opérationnelle de la sécurité.
Le cours est constitué d'un ensemble de notes de synthèse indépendantes compilées en un document unique, mais édité par chapitre dans le cadre de ce cours.\\
C\edoc ne constitue pas à lui seul le référentiel du cours CYBERDEF101 (SEC101 du Cnam). Il compile des notes de cours mises à disposition de l'auditeur comme support pédagogique partiel à ce cours introductif à la cyberdéfense d'entreprise. }
\JournalInfo{\uJournalInfo} % Journal information
\Archive{Notes de cours éditées le  \DTMnow} % Additional notes (e.g. copyright, DOI, review/research article)
\PaperTitle{\utitle} % Article title
\Authors{\uauthor\textsuperscript{1,}\textsuperscript{2}*} % Authors
\affiliation{\textsuperscript{1}\textit{\uproa}} % Author affiliation
\affiliation{\textsuperscript{2}\textit{\uprob}} % Author affiliation
\affiliation{*\textbf{email}: \umaila\ -- \umailb} % Corresponding author
\Keywords{\ukeywords} % Keywords 
\newcommand{\keywordname}{Mots clefs} % Defines the keywords heading name

%-------------------------------------------------
%	MAKE TITLE
%-------------------------------------------------

\maketitle

\begin{tikzpicture}[remember picture,overlay]\node at (current page.north west)[xshift=3cm, yshift=-3cm, text opacity=1.0]{\includegraphics[width=0.1\paperwidth]{\utemplatepath/commons.inc/CommonsPictures/\utemplatetheme.theme/shield-20.pdf}}; 
\end{tikzpicture} 


\AddToShipoutPicture{%
\AtPageCenter{%
\begin{tikzpicture}[remember picture,overlay]
\node [rotate=60,scale=5,text opacity=0.1] at (current page.center) { \textit{\color{uAlertColor}\INFODistrib} };
\end{tikzpicture} 
%\begin{tikzpicture}[remember picture,overlay]
%\node [rotate=60,scale=10,text opacity=0.3]
%at (current page.center) {moi};
%\end{tikzpicture}
}%
}


%-------------------------------------------------	
\begin{center}
	\textcolor{uMainColor}{\textbf{\INFODistrib}} % Distribution information
\end{center}

\INFOGitHub % Link to download file and copyright
%-------------------------------------------------

%-------------------------------------------------
	\newpage
	\tableofcontents  % TOC
	\iftotalfigures
  			\listoffigures
	\fi
%-------------------------------------------------

%-------------------------------------------------
\newpage
\ubody % MAIN BODY defined external
%-------------------------------------------------

	\printbibliography
 \printglossary[type=\acronymtype]
%	\printglossary
%-------------------------------------------------	
\newpage
\Ucontribute % Contribution TEXT
%-------------------------------------------------

  %-------------------------------------------------------
  % AudioNote Open File
  %-------------------------------------------------------
  \mode<all>{\immediate\closeout\audioexport}

\end{document}
}
	\fi
\fi


 

%************************************************************
% Chargement des variables dédiées à l'article
%************************************************************
\lessonID{1}
\setauthor{EDU}

%========================================
\newcommand {\ukeywords}
%========================================
{%---------------------------------------------------------------------
Contexte, Menaces, Vulnérabilités
}%

%========================================
\newcommand {\utitle}
{%---------------------------------------------------------------------
Introduction générale à la Cybersécurité d'entreprise
}%---------------------------------------------------------------------

%========================================
\newcommand {\uabstract}
{%---------------------------------------------------------------------
C\edoc présente comment s'exprimer dans  le même langage que les acteurs du domaine, en ayant compris leurs valeurs, leurs codes, leur environnement et ses contraintes. 
Disposer des éléments de langages du domaine en comprenant les enjeux et le cadre de la Cybersécurité en entreprise.


}
%---------------------------------------------------------------------


%************************************************************
%  variable définissant  le corps de l'article
%************************************************************

%========================================
\newcommand {\ubody}
{%---------------------------------------------------------------------
%$BEGIN_UBODY

%\documentclass[a4paper, 11pt]{book}

\usepackage{graphicx}
\usepackage[french]{babel}
\usepackage{stringstrings}
\usepackage{luacode}
\title{''Titre du livre''}
\author{''Nom de l'auteur''}
\date{''date de fin de rédaction''} 


% Définition de la chaîne de caractères
\def\aline{path/to/your/file.pdf}
\def\mypicturesubpath{Inverse/}



\usepackage{luacode}

% Définition des fonctions Lua
\begin{luacode*}
    function extract_before_last_slash(input_string)
        local last_slash_index = 0
        local before_last_slash = ""
    
        -- Recherche du dernier slash
        for i = 1, #input_string do
            if string.sub(input_string, i, i) == "/" then
                last_slash_index = i
            end
        end
    
        -- Extraction de la sous-chaîne avant le dernier slash
        if last_slash_index > 0 then
            before_last_slash = string.sub(input_string, 1, last_slash_index - 1)
        end
    
        tex.print(before_last_slash) -- Affichage de la chaîne dans le document LaTeX
    end
    
function extract_after_last_slash(input_string)
        local last_slash_index = 0
        local after_last_slash = ""
    
        -- Recherche du dernier slash
        for i = 1, #input_string do
            if string.sub(input_string, i, i) == "/" then
                last_slash_index = i
            end
        end
    
        -- Extraction de la sous-chaîne après le dernier slash
        if last_slash_index > 0 then
            after_last_slash = string.sub(input_string, last_slash_index + 1)
        end
    
        tex.print(after_last_slash) -- Affichage de la chaîne dans le document LaTeX
    end
    
    
function remove_extension(input_string)
    local last_dot_index = input_string:find("%.[^%.]*$") -- Recherche du dernier point (pour l'extension)
    if last_dot_index then
        tex.print(input_string:sub(1, last_dot_index - 1)) -- Affiche la partie sans l'extension
    else
        tex.print(input_string) -- Si pas d'extension trouvée, affiche simplement la chaîne d'entrée
    end
end
    
    function extract_extension(input_string)
    local last_dot_index = input_string:find("%.[^%.]*$") -- Recherche du dernier point (pour l'extension)
    if last_dot_index then
        tex.print(input_string:sub(last_dot_index + 1):upper()) -- Affiche l'extension en majuscules
    else
        tex.print("") -- Si pas d'extension trouvée, affiche une chaîne vide
    end
end
\end{luacode*}
    
\newcommand{\ExtractBeforeLastSlash}[1]{\directlua{extract_before_last_slash("#1")}}
\newcommand{\ExtractAfterLastSlash}[1]{\directlua{extract_after_last_slash("#1")}}
\newcommand{\RemoveExtension}[1]{\directlua{remove_extension("#1")}}
\newcommand{\ExtractExtension}[1]{\directlua{extract_extension("#1")}}

\begin{document}

%PATH : \getBeginPart{\aline}

COMPLET : \aline

%FICHIER : \getLastPart{\aline} 

Sans extension  : \RemoveExtension{\aline}
L'extension : \ExtractExtension{\aline}

Recomposée : \ExtractBeforeLastSlash{\aline}/\ExtractAfterLastSlash{\aline}

\end{document}

%-------------------------------------------------------------
%               FR CYBERDEF SECOPS COURSE
%                                        SECOPS
%                                            Intro
%
%                           Introduction Cyberdefense
%                            % Chap-Intro-gen.tex
%
%                              2020 eduf@ction
%-------------------------------------------------------------


\section{Avant propos}

Chaque jour la presse se fait l'écho d'attaques et de piratages informatiques, de divulgations d'informations sensibles ou  de fragilités découvertes dans les produits et services numériques.
Derrière ces incidents, nous découvrons des menaces certaines fois complexes, des actions criminelles, étatiques ou activistes. Construire des systèmes sûrs, les protéger et les défendre, dans une société où accélérer la digitalisation est devenu un challenge quotidien pour les équipes spécialisées qui luttent contre ces menaces. 
La cybersécurité est un domaine de mythes et de légendes. Ses activités plongent au plus profond de notre histoire avec des luttes ancestrales entre le méchant et le gentil, le gendarme et le voleur, le corsaire et le pirate, en n'oubliant pas les luttes secrètes entre les espions et le contre-espionnage. Une thématique qui résonne, donc comme un domaine de romans, qui se traduit toutefois par une réalité souvent moins réjouissante pour les équipes chargées de la cybersécurité dans les entreprises. Les métiers de la cybersécurité sont nombreux, pour certains très techniques, d'autres plus fonctionnels, juridiques, ou managériaux. 

La cybersécurité est, en effet,  une discipline transverse et interdisciplinaire à plusieurs titres. Elle nécessite  :
\begin{itemize}
  \item de maîtriser les nombreuses techniques et technologies des systèmes d'information ainsi que leurs zones de fragilités;
  \item de maîtriser de nombreuses solutions de sécurité permettant de couvrir, en n'oubliant qu'elles aussi peuvent être fragiles \footnote{cf. Certification et Qualification de produits de sécurité et Critères communs };
%\autoref{chap:ISOCC}

  \item de faire coopérer des métiers et des cultures différentes ;
  \item de gérer l'entreprise dans des cadres de conformité souvent complexes et coûteux;
  \item d'intégrer ces démarches en tenant compte des cultures et des pratiques des nombreux métiers de l'entreprise.
\end{itemize}

\mode<all>{\picframe{../Latex/Sources/EDU/Pictures/img-4model}{Cybersécurité, un domaine hollistique}{0.7}{lbli4model}}

%\includegraphics{../Latex/Sources/EDU/Pictures/img-4model}

Les métiers de la cybersécurité concourent tous à une seule et même mission : \g{\textbf{assurer la continuité de la mission ou du service en préservant le patrimoine de l'entreprise contre toute menace dans l'environnement numérique}}.

Ce sont des métiers de passion, des métiers extrêmement techniques pour partie, fortement marqués par le fonctionnel pour d'autres. S'il est vrai qu'une grande partie des experts du domaine sont issus de formations en informatique ou en électronique, les domaines d'expertises s'élargissent et font naître de nouveaux chemins d'excellence. On trouvera en particulier, des métiers issus du domaine juridique comme celui du \uindex{Data Protection Officer}.    

%\section{les grands cyber-risques}

% les grands typologies de menaces liées au politique,, au vol, à la guerre économique" : Chantage, Images, Vole, Continuité d'activité ... manipulation, espionnage ...

% Les cyberrisques dependent de la cibles, et souvent la cible est liée à l'attaquant

%\section{les vecteurs cyber}

\section{Aborder la cybersécurité}

% ajouter le joke sur les corsaires et une hallebarde

La cybersécurité ou la sécurité du numérique\footnote{Historiquement d'autres termes sont utilisés comme  \ugls{aSSI} ou  Sécurité Informatique} peut être découverte par de nombreuses voies.

La plus courante est certainement pour les technophiles, l'aventure passionnante de découvrir ce domaine par la technique, et le hacking.  Longtemps abordé par le triptyque académique cryptologie,   sécurité  protocolaire des réseaux, et informatique fondamentale (compilation et théorie des langages, architecture système et bases de données), le domaine s'est vulgarisé avec une forme de gamification de l'apprentissage.

On y trouve en particulier :

\begin{itemize}
  \item Les challenges comme les \UKword{\ugls{aCTF}}, ou les \UKword{\ugls{aDTF}} qui permettent de mettre le pied dans les techniques et stratégies d'intrusion pour les pentesteurs et auditeurs techniques en herbe;
  \item Les bug-bounty qui permettent de se confronter à ses propres limites avec la recherche de failles dans les logiciels avec pour partie des rémunérations au niveau des difficultés;
\end{itemize}

Toutefois, un volet peu enseigné, qui ne mobilise pas spécialement les jeunes apprenants du domaine concerne la gouvernance de cette sécurité numérique de l'entreprise.

J'ai souhaité m'intégrer dans  une approche globale de la sécurité du numérique pas le biais de quelques processus, en particulier ceux de la sécurité opérationnelle. L'enjeu est de fournir une trame de connaissances pour déployer des actions de cybersécurité en entreprise.
Cette trame a pour intention de fournir des points d'accroche et des modèles de compréhension des différentes compétences, actions, et outils du large domaine de la cybersécurité.

Destiné à un public large, cet ouvrage tente d'offrir un niveau de lecture permettant à un expert technique de repositionner sa technicité dans un ensemble plus large, et  à un débutant de découvrir de nombreuses facettes du domaine avec quelques éclairages techniques.

La cybersécurité dans une entreprise est une co-activité d'hommes de l'art.  C'est aussi un domaine en perpétuelle évolution, soutenu et contraint par des lois, des règlements, des normes, des méthodologies,  des technologies spécialisées et en particulier des expertises.  Il nécessite pour être efficace d'être orchestré pour maintenir en condition de sécurité une organisation dont le périmètre peut être complexe face à des menaces elles aussi en perpétuelles évolutions.

Il y a de nombreuses manières d’aborder le pilotage de la cybersécurité au sein de l’entreprise, et nombreux ouvrages spécialisés en détaillent les concepts et les méthodologies. Nous avons toutefois délibérément choisi ici de confronter, si ce n'est de corréler, dans un seul support, trois domaines qui apparaissent souvent dans la littérature comme des domaines d'expertise différents : la gestion des risques, la gouvernance de la cybersécurité et la cybersécurité opérationnelle. 

Nous avons donc fait ce choix de structurer notre approche suivant le prisme de la cyberdéfense d'entreprise avec une analyse en trois axes majeurs qui résument les difficultés dont relève cette discipline holistique \ucite{sch13}. 


% begin PRZ===================
\begin{frame}{Les 3 axes de la cybersécurité d'entreprise}
	\begin{itemize}
 \item \tb{l'analyse des risques} informatiques sur les actifs les plus sensibles de l'entreprise avec les difficultés d'identifier ces actifs ainsi que leur sensibilité et les menaces qui pèsent sur ceux-ci ou sur l'environnement ;
\item la structuration d'une gouvernance efficace avec des \tb{politiques de sécurité} des systèmes d'information pour des architectures de sécurité de confiance, dans des systèmes d'information complexes, intégrant des services dans le cloud, des technologies obsolètes et des politiques de sécurité sédimentées ;
\item la construction et l’organisation d'une \textbf{sécurité opérationnelle} vue sous un angle d'anticipation et de veille, de détection, et enfin d'alerte et de réponse aux attaques, nécessitant une activité continue avec des ressources de plus en plus expertes et avec des outils plus \g{hyperspécialisés} et de plus en plus complexe à déployer.
\end{itemize}
\end{frame}
% end PRZ====================

%\rem{Politiques : le monde de la reflexion, Stratégies le monde l'action}
\subsection {Politiques versus stratégies}


\mode<all>{\picframe{../Latex/Sources/EDU/Pictures/img-cycle}{Processus Cyber d'entreprise}{0.7}{lblcycle}}


Nous avons structuré la dynamique des processus d'entreprise dans c\edoc suivant un schéma simple\ uref{processus}{lblcycle}.
De  \head{l’analyse de risque}, nous pouvons déduire et/ou modifier des politiques de sécurité adaptées au contexte de l'entreprise. 
Sur la base de l’existant, il est alors possible d’adapter ou simplement de mieux utiliser ou configurer les architectures techniques et organisationnelles pour définir un niveau de protection attendu.
Il y a malheureusement toujours un écart entre les mesures de sécurité souhaitées et la réalité des mesures déployées. Que ce soit à cause de défauts de configuration, de délais de mise en place plus longs que prévu, le système n’est que très rarement au niveau décrit dans les éléments de spécification ou les documents d'assurance sécurité.
Mesurer ce niveau, analyser les écarts et remédier relève d’un des grands thèmes de la gouvernance sécurité. Il reste à lui seul un consommateur à plus de 30\% des charges d'activité de cette gouvernance.

Après avoir défini des \head{politiques de sécurité} et mesurer leur déploiement dans l’environnement de l’entreprise, il n’en demeure pas moins que l’ennemi est toujours à ses portes et de plus en plus souvent, il arrive à pénétrer le périmètre de sécurité.
Non pas que les barrières et filtre de l'entreprise ne soient plus efficaces mais simplement parce que l'attaquant change plus souvent de stratégie que l'entreprise de politique. La sécurité se doit donc d'être plus dynamique.
L’entreprise doit faire face à des attaquants qui ne raisonnent pas sous forme de politiques d’attaque, mais en stratégie d'action. L’entreprise doit raisonner aussi de la même manière pour se défendre. C'est à ce titre que l'on parle de stratégie de Cyberdéfense.
C'est avec des stratégies de \g{cyberdéfense} que nous aborderons les moyens organisationnels et techniques à mettre en place.

Vous trouverez  dans ce document une terminologie qui peut être certaines fois éloignée des expressions classiques de la sécurité informatique. \\
J'ai choisi d'utiliser et mixer sans trop de complexes des termes et concepts issus du monde militaire (renseignement, tenir une position, infiltration …) et de nombreux autres issus de l'univers médical (infection, épidémie, comportement).\\ Ces incursions dans les analogies d'autres champs professionnels, bien que présents pour illustrer certains concepts, n'en demeurent pas moins justifiés par leurs usages de plus en plus répandus dans le monde de la cybersécurité. Par ailleurs, les termes sécurité et cybersécurité pourront être utilisés indifféremment dans le corps de ce document.

\subsection {Transformation numérique}

% Les causes de la complexité

La cybersécurité est devenue en quelques années un axe fondamental dans la prise en compte de ces nouveaux risques sociétaux qu'apporte l'informatique au cœur de chaque activité sociale, économique ou politique.

%Quelque que soit le niveau d'impact  --- personnel, entreprise, Etatique, 
Les  transformations digitales d'une grande partie des acteurs économiques apportent de nouveaux risques. Les modifications des conditions d'utilisation des technologies dans les crises comme celle du COVID-19 engendrent aussi des risques globaux réduisant les frontières entre les espaces professionnels et les espaces privés.

Depuis bien des années, le législateur s'en est saisi avec de nombreuses réglementations et lois permettant de protéger en particulier, le citoyen et l'Etat.

On notera en particulier dans cette évolution du cadre réglementaire,  la protection des données personnelles, mais aussi la protection des systèmes sensibles stratégiques \footnote{\ugls{aOIV}  et \ugls{aOSE} } en lien avec la protection de la nation avec la dynamique de Cyberdéfense soutenue par les différentes lois de programmation militaire depuis 2008, et récemment la directive \ulink{https://www.ssi.gouv.fr/actualite/adoption-de-la-directive-network-and-information-security-nis-lanssi-pilote-de-la-transposition-en-france/}{NIS 2} . L'entreprise se trouve quant à elle prise en sandwich entre les exigences de l'état et les désirs de liberté que défend le citoyen. Il faut aussi noter que le salarié est souvent un citoyen et son rôle dans la cybersécurité de l'entreprise peut soulever des problématiques complexes\footnote{Lanceurs d'alertes, comportements déviants des utilisateurs légitimes}.
 
Se sentir en sécurité dans un monde de transformation digitale c'est bien entendu disposer des moyens de se protéger et protéger son patrimoine, que ce dernier soit ou non informationnel, mais aussi de le défendre en continu. Il est de moins en moins accepté de le protéger, en érigeant des murs épais, solides et supposés infranchissables. L'entreprise a besoin de faire circuler rapidement les savoirs, de partager largement des informations entre les salariés, les clients, les citoyens, les fournisseurs...


Il est donc nécessaire de correctement définir les biens vitaux ou essentiels pour y mettre les meilleurs moyens pour les défendre. Par ailleurs comme toute activité protégée et défendue qui peut subir des dommages, il est important de structurer l'activité numérique d'une entreprise ou d'une organisation pour pouvoir fonctionner en mode dégradé, et revenir à la normale en moins de temps possible.


Entre une maîtrise des risques cyber et une capacité de se défendre et réagir, il est nécessaire de disposer déjà d'un bon niveau de protection adaptée aux enjeux du numérique. Il existe de nombreuses définitions de cette cybersécurité.

Pour ma part je vous propose de poser pour la suite de mon propos, une définition simple, qui fait consensus et résume en une pseudo équation la manière dont nous traiterons ce domaine dans c\ecours. \\
\begin{nota}[Une définition de la cybersécurité]
\begin{align}
Cybers\acute{e}curit\acute{e} \cong Cyberprotection\oplus Cyberd\acute{e}fense \oplus Cyberr\acute{e}silience
\end{align}
\end{nota}


% begin PRZ====================
\begin{frame}<presentation>{Une définition de la cybersécurité}
\begin{align}
{\large
Cybers\acute{e}curit\acute{e} \cong Cyberprotection\oplus Cyberd\acute{e}fense \oplus Cyberr\acute{e}silience
        }
\end{align}
\end{frame}
% end PRZ====================

% begin PRZ====================
\begin{frame}

La cybersécurité est l'enchaînement opéré, organisé, documenté, piloté, optimisé de trois environnements d'actions :
\begin{itemize}
 \item \tb{Protéger} l'environnement par les mesures et solutions technologies adaptées au niveau de risque que l'entreprise est prête à prendre; 
 \item \tb{Défendre} les actifs les plus sensibles de l'entreprise en surveillant et combattant la menace (y compris l'image de l'entreprise);
 \item assurer \tb{la continuité et la reprise d'activité} de l'entreprise face à tout incident rendant indisponible tout ou partie d'une fonction essentielle de celle-ci.
\end{itemize}

\end{frame}
% end PRZ====================

La Cybersécurité, est donc, avant-tout, le déploiement de mécanismes de protection des biens et des processus numériques sensibles. C’est avec cette première dynamique que l’entreprise déploie en premier lieu des solutions de sécurité. 

Toutefois, malgré ce niveau de protection et souvent les lourds investissements réalisés dans des composants de sécurité périmètrique, l’entreprise peut se faire surprendre avec des attaques contournant ces mesures. Face à ces attaques, l’entreprise découvre que la solidité de l’entreprise n’est pas directement liée aux investissements sur les systèmes de protections. Il lui faut anticiper les menaces, les détecter non seulement sur son périmètre mais aussi dans l’écosystème de l’environnement des menaces potentielles. Ces menaces exploitent des vulnérabilités qu’il convient de détecter en amont.

Malheureusement, malgré ces mesures de protection et de défense qui permet de réagir vite et efficacement, il arrive que des attaques informatiques arrivent à leurs fins. La capacité de l’entreprise à revenir à une situation normale, avec un contexte assaini est un critère dont un chef d’entreprise appréciera la valeur \textbf{qu’après} un incident.

% ajouter l’equation du risque dans le texte du RISQUE

Il fut une époque pas si lointaine, où l'analyste consacrait beaucoup de temps à l'évaluation de la probabilité d'une attaque. Aujourd'hui bien que ce paramètre continue quand même à être pris en compte, l'analyste positionne cette probabilité ou \uindex{vraisemblance}
 à 100\% car pour l'entreprise la nouvelle hypothèse n'est pas de savoir "si un attaque va avoir lieu mais plutôt quand elle aura lieu" \uref{figure}{lblriskassess}.

% I M A G E -----------> 


\mode<all>{\picframe{../Latex/Sources/EDU/Pictures/img-risk-assess}{le cyber-risque}{0.8}{lblriskassess}}

L'ensemble des experts du domaine est globalement en accord sur la posture que doivent prendre les entreprises et les organisations : \g{Le temps n'est plus de savoir si on sera attaqué ou pas, mais plutôt de savoir quand et comment on le sera}, qui concrètement se résume à la certitude que tout incident de sécurité peut se produire.

\mode<all>{\picframe{../Latex/Sources/EDU/Pictures/img-risk-prison}{La menace, une vision de l'attaquant}{0.8}{lblriskjail}}

 
Dans les modèles d'analyse de risque et d'évaluation de la cybersécurité, l'analyste se positionne aujourd'hui du point de vue de l'attaquant. Ce regard lui permet de mieux comprendre la menace comme équation duale du risque, mais vu de l'énergie dépensée par l'attaquant et du risque qu'il prend \uref{figure}{lblriskjail}. La menace sera d'autant plus grand que la valeur des actifs visés est importante à ses yeux et que les vulnérabilités sont nombreuses. Elle sera d'autant plus faible que les moyens à déployer sont importants et les risques pris plus élevés.

On notera que la valeur des actifs visés dépend de l'attaquant. Ils ne seront pas les mêmes si l'attaquant est un état, une groupe d'activiste, un groupe criminelle.
Quand aux vulnérabilités, la majorité des attaques utilise les failles perpétuelles du facteur humain.




\section{Sécurité du système d'information}

Le système d'information est au coeur de ce \g{monde digital}, et il est le lieu d'activités humaines très denses, permettant à des utilisateurs de réaliser leurs activités, professionnelles ou privées, à l'aide de processus informatiques et de services. 

Ces activités doivent de plus en plus faire face à tout un système d'agression orchestré par des attaquants non seulement humains mais aussi \g{automatiques}. On observe aujourd'hui une multitude de situations critiques, incertaines dont l'occurrence quasi quotidienne provient de phénomènes variés, humains (isolés, en réseau, …), physiques et/ou technologiques. Parmi ces difficultés qui profitent aux pirates informatiques il y a de nombreuses failles ou fragilités que nous découvrirons, provenant des systèmes du SI sans lesquelles ils ne pourraient exploiter leurs attaques. Ces phénomènes sont une menace pour les conditions de sécurité du système d'information.


\subsection{Gouvernance et conformité}

La conformité en matière de cybersécurité dans une entreprise fait référence au respect des normes, réglementations et exigences légales en matière de protection des données et de sécurité informatique. Les enjeux liés à la conformité dans la gouvernance de la cybersécurité d'une entreprise sont nombreux et cruciaux pour assurer la protection des informations sensibles et la réduction des risques.


\begin{itemize}
	\item Respect des réglementations : Les entreprises sont soumises à des réglementations de plus en plus strictes en matière de cybersécurité, telles que le Règlement général sur la protection des données (RGPD) en Europe ou les lois liées à la protection des infrastructure critiques sur la protection (Loi de programmation militaire en France), avec son cortèges de référentiels (CNIL, ANSSI, Défense ...). Le non-respect de ces réglementations peut entraîner des sanctions financières non néglieables, des poursuites et condamnations judiciaires pouvant induire des défauts de réputation de l'entreprise.
	\item Gestion des risques : La conformité permet aux entreprises de gérer efficacement les risques liés à la cybersécurité. En identifiant les menaces potentielles, en évaluant les vulnérabilités et en mettant en place des contrôles de sécurité adéquats, les entreprises peuvent réduire les risques d'attaques informatiques, de pertes de données et de fraude dans le numérique. Cette gestion des risques est aussi contrainte par les réglementations, obligeant les entreprises à documenter leur processus mais aussi la mise en œuvre de la gestion des risques dans les projets, produits et services de l'entreprise.
	\item Continuité des activités : Les incidents de cybersécurité peuvent paralyser les opérations d'une entreprise et entraîner des perturbations importantes. La conformité en matière de cybersécurité aide à garantir la continuité des activités en mettant en place des plans de reprise après sinistre, des sauvegardes régulières des données et des procédures de réponse aux incidents.
	\item Protection des données personnelles : La conformité en matière de cybersécurité vise à garantir la protection des données personnelles des clients, des employés et des partenaires commerciaux. Les entreprises doivent mettre en place des mesures de sécurité appropriées pour prévenir les violations de données, les fuites d'informations sensibles et les accès non autorisés.
	\item Confiance des clients : La conformité renforce la confiance des clients dans une entreprise. Les clients sont de plus en plus préoccupés par la confidentialité et la sécurité de leurs données personnelles. En démontrant leur conformité aux réglementations en matière de cybersécurité, les entreprises peuvent gagner la confiance de leurs clients et maintenir des relations commerciales solides.
	\item Avantage concurrentiel : La conformité en matière de cybersécurité peut également conférer un avantage concurrentiel aux entreprises. En étant en conformité avec les réglementations et les meilleures pratiques en matière de cybersécurité, les entreprises peuvent attirer de nouveaux clients, renforcer leur réputation et se différencier de leurs concurrents moins attentifs à ces enjeux.
	
\end{itemize}

Cette conformité en matière de cybersécurité est essentielle pour déployer et prouver l'utilisation des bonnes pratiques dans la protection des données, la gestion des risques et l'engagement sécurité vis à vis des clients (B2C et B2B). Gouverner la cybersécurité de l'entreprise c'est aussi investir dans des mesures de sécurité adéquates, la formation du personnel, des politiques et procédures solides, le tout encadrer par des mécanises d'audits réguliers pour garantir cette conformité aux réglementations dont le corpus est en constante évolution.

\subsection{SECOPS et lutte contre la malveillance}

\utodo

\subsection{Les fonctions SSI de gouvernance}

Au sein des grandes entreprises, il existe de nombreuses fonctions ou missions pour gouverner, piloter cette sécurité numérique.

%----begin FRAME-------
\mode<all>{\texframe
{SSI : Responsabilités}
{les métiers}
{
\begin{itemize}
 \item \head{Le gestionnaire de risque} ou \UKword{Risk Manager} qui porte l'animation de la gestion des risques dans les projets ou dans l'entreprise;
\item \head{Le responsable sureté / sécurité } généralement responsable de la sécurité physique ou sein de l'entreprise (vol, intrusion physique, contrôle d'accès). Il endosse le plus souvent la responsabilité des biens et des personnes;
\item \head{L'audit et le contrôle} : Au sein des grandes organisations, il peut exister un service \g{indépendant} dont la mission est d'auditer et de contrôler les activités des services;
\item \head{Les RSSI} : Responsables de la sécurité des Systèmes d'Information;
\item \head{Les DSSI} : Au sein des grandes entreprises, les RSSI globaux ne dépendent plus trop de DSI, et possèdent le rang de directeur;
\item \head{Le DPO} : la dernière responsabilité apparue dans l'environnement de la sécurité (En France successeur du CIL , Correspondant Informatique et Liberté) (\UKword{Data Protection Officer}).
\end{itemize}
}}%----end FRAME-------


Nous ne présentons rapidement ici que ceux qui seront utilisés directement dans ce document et qui sont fortement en liaison avec la sécurité des systèmes d'information.

\subsubsection{Les DSSI et RSSI}

Au sein de l'entreprise, il est important que quelqu'un porte la charge de suivre ces conditions de sécurité. C'est le rôle du RSSI (Responsable de la sécurité des systèmes d'information), ou DSSI (Directeur de la sécurité des systèmes d'information).
La mission de ce RSSI d'entreprise est de protéger son Système d'Information (SI), de le mettre dans une posture d'amélioration continue tant de son système de protection que de son système de défense. Le RSSI n'est pas seul pour assumer ces missions, à tous les niveaux de l'entreprise, s'organise des fonctions de sécurité tant au sein de la DSI (auquel est souvent rattaché le RSSI), qu'au sein d'autres activités de l'entreprise. 

\subsubsection{Le DPO}

A partir de mai 2018, une responsabilité plus juridique liée à la protection des données a été rendue plus visible avec la nécessité de disposer d'un \ugls{aDPO} en entreprise (Data Protection Officer) héritier en France  \textit{\ugls{aCIL}}.
Nous n'aborderons pas la fonction, les missions et la dynamique de responsabilité du \ugls{aDPO} ici. Il faut toutefois que cette fonction possède de nombreux recouvrements dans la chaine de gouvernance du risque \g{informatique} auprès des directions d'entreprise. Orienté vers la protection des données à caractère personnel, le tropisme de la fonction \ugls{aDPO} peut conduire certaines structures à oublier des pans importants des risques numériques comme :

\begin{itemize}
	\item la protection du patrimoine informationnel. (Espionage industriel).
	\item la protection des systèmes d'information contre les risques de ruptures de services (Continuité d'activité)
\end{itemize}

\subsubsection{L'officier de sécurité de défense}
Pour les entreprises traitant des informations classifiées de défense ou liées aux contraintes de la classification de l'état, il est indispensable de se doter d'une fonction \ugls{aOS} de défense. Son rôle est de s'assurer de la conformité à l'Instruction Générale Interministérielle IGI 1300 pour le \g{Confidentiel Défense} et l'instruction Interministérielle II901 pour le \g{Diffusion Restreinte}.

Nous resterons donc dans le cadre fonctionnel de la Cybersécurité dans son volet protection des Systèmes d'information et gestion des risques numériques.
Quand nous aborderons des sujets en forte adhérence avec les dynamiques de la \ugls{aGDPR} nous donnerons les liens et les indications adaptés pour les \ugls{aDPO}.
Par exemple, nous aborderons l'usage des données nominatives collectées et traitées dans les \ugls{aSIEM}, ou celles recueillies sur le DarkWeb etc ...

Consulter le site de la \ulink{www.cnil.fr}{CNIL} pour parfaire ses connaissances en matière de réglementation européenne sur la protection des données personnelles. 


\subsubsection{Responsabilités SSI}

Le maintien des conditions de sécurité du système d'information des grandes entreprises nécessite un \ugls{aRSSI} Central ou une fonction semblable rattachée à un niveau plus global de l'entreprise.
On découvre ainsi des  \ugls{aRSSI} rattachés à la direction des risques, la direction générale, au contrôle interne …\\
 Il n'y pas de rattachement bien défini. La couverture de responsabilité dépend grandement de la taille et de l'activité de l'entreprise, mais aussi de la maturité de celle-ci en matière de gestion de risque et de gouvernance. Il peut y avoir des RSSI par entité, par projet à l’intérieur d'une entreprise. Leur mandat est fixé en fonction des enjeux sécurité de ces entités ou ces projets. 

Le fin mot de l'histoire est le \g{R} de  \ugls{aRSSI}. Son domaine de responsabilité dépendra de son mandat pour assumer ce rôle de garant d'un environnement \g{possédant } des bonnes conditions de sécurité.
La gouvernance de la sécurité, est au coeur du métier du  \ugls{aRSSI}. Cette discipline que ce dernier pilote dans l’entreprise se nomme \ugls{aGRC}.

La notion de système d'information a profondément évoluée ces dernières années. Le périmètre des risques digitaux inclut maintenant des systèmes et services externes à l'entreprise. Beaucoup d'entre eux sous la forme de réseaux sociaux, de services \uindex{cloud}
 ouvrant par ailleurs le domaine de supervision à la téléphonie avec les smartphones et leurs applications professionnelles ou non.

Bien entendu en fonction de la taille de l'entreprise et de ses enjeux, on peut disposer au sein de l'entreprise de nombreuses personnes ayant une fonction de \ugls{aRSSI}. 

Le métier est riche et dispose d'un spectre de responsabilité et d'activité très large en terme de poste on y trouve par exemple :

% FRAME beamer PRZ ------------------------------------
\mode<all>{\texframe{Fonctions RSSI}{différents métiers} 
{
\begin{itemize}
\item \textbf{RSSI d'entreprise} : Responsable de la sécurité de sa structure.
\item  \textbf{RSSI d'un département, d'une organisation intermédiaire} : A l'image d'un RSSI d'entreprise, il assure toute les tâches de gouvernance, il applique et fait appliquer les directives et politique de sécurité aux équipes du département / division / structure intermédiaire, il déploie les actions décidées dans la chaîne fonctionnelle sécurité
\item   \textbf{RSSI d'un contrat, d'un projet contractualisé (Security Manager)} : Responsable de la sécurité du  déroulement d'un contrat. Souvent lié à un plan d'assurance sécurité, le RSSI contrat se doit d'assurer pour le client ou pour le fournisseur le suivi des exigences de sécurité du contrat.
\end{itemize}
}} % end FRAME.........................................................


\mode<all>{\texframe{Fonctions RSSI}{suite}
%. . . . . . . . . . . . . . . . . . . . . . . . . . . . . . . . . . . . . . . . . . . .
{\begin{itemize}
	\item \textbf{RSSI Projet} : La responsabilité sécurité couvre le projet, on parle souvent de \g{\uindex{security by design}
}. La responsabilité dans ce type de poste recouvre l'intégration de la sécurité dans le système, le suivi des indicateurs définis (contractuels, ou réglementaires), la remontée des indicateurs de suivi de sécurité à la MOA (Maitrise d'ouvrage), la prise de décision autour des choix de sécurité ....
	\item \textbf{RSSI Produit / Service} : Au delà de ce qui est fait pour un projet, le RSSI produit a en charge de gérer la sécurité opérationnel c'est à dire Maintenir la sécurité de son produit ou de son service.
 \item \textbf{RSOP} : Le responsable sécurité opérationnelle, est souvent un RSSI dépendant d'une DSI, il est généralement et dans beaucoup de d'entreprise de taille moyenne le RSSI technique. Il assure opérationnellement la mise en place technique des politiques de sécurité et maintien en condition de sécurité l'ensemble de l'environnement informatique. Il est aujourd'hui au coeur de la sécurité opérationnelle face aux attaques et aux crises cyber.
\end{itemize}
}} % end FRAME.........................................................


\subsection{Maintien en condition de sécurité}

Les conditions de sécurité représentent les propriétés fondamentales du SI, appelées : \ugls{aDICT} , qui favoriseront le fonctionnement optimisé du SI et éviteront l'avènement d'incidents de sécurité irréversibles ou même gênants pour son fonctionnement. D'un certain point de vue, les conditions de sécurité représentent le paramétrage du SI pour lequel le système fonctionne bien dans des conditions de sécurité \g{connues et approuvées}.

Ces fameux critères \textbf{\ugls{aDICT}} ou propriétés de sécurité des systèmes d'information visent les objectifs suivants :
%-----------------------------------------
\begin{nota}[DISPONIBILITE]
le système doit fonctionner sans faille (arrêt, ou dégradation) durant les plages d'utilisation prévues et garantir l'accès aux services et ressources définies et installées avec le temps de réponse attendu.
\end{nota}
%-----------------------------------------
\begin{nota}[INTEGRITE]
Les données doivent être celles que l'on attend, et ne doivent pas être altérées de façon fortuite, illicite ou malveillante. En clair, les éléments considérés doivent être exacts et complets.
\end{nota}
%-----------------------------------------
\begin{nota}[CONFIDENTIALITE] Seules les personnes autorisées peuvent avoir accès aux informations qui leur sont destinées. Tout accès indésirable doit être empêché.
\end{nota}
%-----------------------------------------
\begin{nota}[TRACABILITE] (ou preuve ) : garantie que les accès et tentatives d'accès aux éléments considérés sont tracés et que ces traces sont conservées et exploitables.
\end{nota}
%-----------------------------------------
D'autres aspects peuvent aussi être considérés comme des objectifs de la sécurité tels que :
%-----------------------------------------
\begin{nota}[AUTHENTICITE]
l'identification des utilisateurs est fondamentale pour gérer les accès aux espaces de travail pertinents et maintenir la confiance dans les relations d'échange. On voit aussi dans la littérature la terminologie \g{critères \ugls{aACID} (Authentification, Confidentialité, Intégrité, Disponibilité)}.
\end{nota}
%-----------------------------------------
\begin{nota}[NON-REPUDIATION]
La non-répudiation et l'imputation : aucun utilisateur ne doit pouvoir contester les opérations qu'il a réalisées dans le cadre de ses actions autorisées et aucun tiers ne doit pouvoir s'attribuer les actions d'un autre utilisateur.
\end{nota}

Dans un contexte d'activité économique dense et en perpétuel renouvellement, les conditions de sécurité sont aussi en perpétuelle évolution, c’est pourquoi nous parlons d'un cycle de vie vertueux au cours duquel les nouveaux paramètres tirent profit des expériences passées.  Ainsi, l’amélioration continue également appelée \g{ lean management} dans d’autres domaines (industrie, …) travaille elle sur le cycle de vie des conditions de sécurité souvent appelé \ugls{aPDCA}.\\
Ce cycle de vie doit néanmoins être maitrisé par le RSSI en place avec ses équipes, il faut co-produire ces conditions de sécurité, cette maîtrise est complexe, fortement dépendante du contexte de l'entreprise, c’est pourquoi elle doit être accompagnée d'une méthodologie rigoureuse et partagée qui constitue le savoir-faire de base du RSSI et de son équipe. Par ailleurs, parmi ces conditions, certaines sont universelles et d’autres propres à chaque entreprise. \\ Comme le montre le cycle de vie. \uref{diagramme}{lblcycledevie}, il est possible aussi d’utiliser un cycle de vie sécurité de type projet, qui se rapproche par ailleurs de la manière dont nous avons structuré c\edoc. 

\mode<all>{\picframe{../Latex/Sources/EDU/Pictures/img-cyclevie}{Cycle de vie sécurité dans les projets}{0.8}{lblcycledevie}}


Dans cette optique, ce cadre méthodologie a été défini par le sous-comité 27 de l'ISO, par l'ensemble de normes ISO 27x. Il s’agit également d'un ensemble de bonnes pratiques, qu'un RSSI peut suivre au travers de trois volets fondamentaux qui constituent les référentiels utilisés pour ce cours sur la cybersécurité. La norme 27001 est en particulier un cadre pour organiser la dynamique de la mise en condition de sécurité de l’entreprise et son maintien dans le temps. Cet environnement que le \g{RSSI} doit bâtir est le système de management de la sécurité (\g{SMSI}).\\
Notre dynamique méthodologique est soutenue dans  c\edoc, par trois cadres normatifs : 
 
% FRAME beamer PRZ ------------------------------------
\mode<all>{\texframe{Cadres normatifs}{3 modèles du cadre}
%. . . . . . . . . . . . . . . . . . . . . . . . . . . . . . . . . . . . . . . . . . . .
{
\begin{itemize}
\item Identifier ses cyber-risques sur la base de méthodologies que l’on retrouve dans l’environnement ISO/CEI 27001/27005 mais aussi sur la méthodologie \ugls{aEBIOS} de l’ANSSI (Méthode \ugls{aEBIOS} RM en particulier); 
\item Elaborer une politique de cybersécurité sur la base des cadres ISO/CEI  27001 et 27002, en n’oubliant pas les architectures de sécurité et la sécurité des architectures associées ; 
\item Détecter en amont des attaques et savoir réagir à ses cyber-incidents en se basant sur ISO 27035 et sur la continuité d’activité avec l’ISO 22301 et 27031.
\end{itemize}
}} % end FRAME.........................................................


\begin{nota}[Pourquoi des normes dans c\edoc?]
L’objectif de c\edoc, n’est pas de présenter en détail un cadre normatif, mais bien de les utiliser pour ce qu’elles sont : des langages communs permettant d’appréhender une terminologie, des méthodologies, des outils. L’ISO 27001 comporte un grand nombre de normes (plus de 50…) qu’il convient de connaître comme outils terminologiques et de référence. Leur maîtrise nécessite une spécialisation le plus souvent demandée pour des métiers de conseil ou d’implémentation pour une certification.
\end{nota}

Ces documents définissent un cadre méthodologique et normatif pour définir, créer, élaborer maintenir, améliorer les conditions ou les critères de sécurité pour le fonctionnement du système protégé et surveillé. 
Ils permettent aux acteurs de l'entreprise évoluant autour du métier RSSI un cadre méthodologique ainsi qu’un \g{how to} du maintien en conditions de sécurité. C’est en particulier au travers de ces trois axes que la mission de RSSI repose. 
Le nombre d’entreprises prêtes à accueillir des spécialistes de ce savoir-faire est en forte augmentation car les PME/PMI ont pris conscience que la sécurisation de l’entreprise est devenu primordiale pour \g{survivre} dans l’écosystème digital de nos sociétés modernes. Les contraintes légales issues de la Loi de Programmation Militaire (\ugls{gLPM}), de la Règlementation pour la Protection des Données Personnelles (RGPD), de la directive NIS nécessitent de disposer d’une vision globale et transverse tant technique, qu'organisationnelle ou humaine de la cybersécurité.

Nous tenterons donc dans le suite du cours, de vous donner des contextes d'usages de ces cadres normatifs indispensables pour aborder la cyberdéfense d’entreprise. 

\section{Enjeux légaux}

Beaucoup d'environnements normatifs sont issus de la pression des différents cadres législatifs sur le marché. Que ce soit avec la pression du grand public ou avec les enjeux stratégiques et économiques des pays, ces lois organisent profondément les modes de gouvernance de la sécurité en entreprise.

\subsection{Quelques cadres législatifs d'influence}
Parmi les grandes lois qui ont influencées  le monde de la sécurité des entreprises ces dix dernières années :

\begin{itemize}
 \item En France, le cyberdéfense est largement orientée par les différentes \g{Lois de programmation militaire} avec des directives nationale de sécurité par grands domaines d'infrastructures vitales.
 \item En Europe deux grandes directives ont donné plus de responsabilités aux entreprises dans l'engagement sécurité avec GDPR et NIS qui sont déclinés en droits français via la CNIL, et l'ANSSI. On notera par ailleurs la montée en puissance dans la confiance numérique avec le cadre de certification européen.
 \item Aux Etats Unis, le \UKword{Cloud security Act}, a bouleversé la vision des risques numériques des états avec les potentielles nuisances liées à l'extraterritorialité de lois américaines
 \item En Russie et en Chine, plusieurs lois autour de l'usage d'internet interpellent les entreprises et en particulier celles du numérique sur la protection des données de leurs clients ou utilisateurs de leurs services.
\end{itemize}


\subsection{Le cadre de certification européen}

Le règlement établit un \ulink{https://www.ssi.gouv.fr/entreprise/reglementation/cybersecurity-act-2/le-cadre-de-certification-europeen}{cadre européen de certification}
 de cybersécurité pour harmoniser à l’échelle européenne les méthodes d’évaluation et les différents niveaux d’assurance de la certification, au sein duquel l’ENISA trouve toute sa place. Les certificats délivrés bénéficieront d’une reconnaissance mutuelle au sein de l’Union européenne (UE). 

\subsection{Cyberdefense et loi de programmation militaire}
Pour ceux intéressés par les contraintes et cadres généraux de la cyberdéfense au sein des lois de programmation successives (2008, 2013, 2019 ...) il est conseiller d'aller voir sur le site de l'ANSSI. Les différentes LPM ont fait évoluer le cadre réglementaire pour assurer à la France une capacité de défendre la continuité de l'état et des infrastructures vitales du pays \ulink{https://www.ssi.gouv.fr/entreprise/protection-des-oiv/protection-des-oiv-en-france/}{(Cf. Opérateurs d'infrastructures vitales)}.

%-------------------------------------------------------------
%               FR CYBERDEF SECOPS COURSE
%                                        SECOPS
%                                            Intro
%
%                           Introduction Cyberdefense
%                            % Chap-Intro-Ref.tex
%
%                              2020 eduf@ction
%-------------------------------------------------------------



\section{Quelques organismes de référence}

Pour l'entreprise la cybersécurité est un domaine  de nombreux cadres normatifs et réglementaires soutenus bien souvent par contraintes legislatives propres à chaque pays.

Cette normalisation et ses réglementations sont riches mais certaines fois complexes.
Le plus simple pour s'enrichir de ces savoirs et surtout pour disposer des meilleurs informations à la sources autant \g{fréquenter} les sites internet institutionnels des organismes qui sont et continuent à être les points de  référence dans le domaine de cybersécurité.

De nombreux services étatiques et de normalisation possèdent des activités dites Cyber dans leur structures :

\begin{itemize}
    \item  Organismes français : AFNOR, Cert FR, CNIL, HADOPI, ANSSI, DGSE, DGSI, DGA/MI, Commandement de la cyberdéfense, C3N, OCLCTIC, BEFTI ...
  \item  Organismes internationaux : ISO, ETSI, CERT, Europo, lnterpol , ENISA , FIRST ...
  \item  Organismes étrangers : FBI, CIA, NSA, GCHQ, Unité 8200, Fapsi, The SANS institute, CISA ...
\end{itemize}

Je vous propose de donner quelques pointeurs par portée sur des organismes de référence du point de vue occidental.

\subsection{International et Etats-Unis}

Au niveau international, on ne peut éviter les Etats-Unis, un pays qui oeuvre fortement dans le domaine des standards.

\subsubsection{Le NIST}
Le National Institute of Standards and Technology, ou NIST est une agence du département du Commerce des États-Unis. Son but est de promouvoir l'économie en développant des technologies, la métrologie et des standards avec l'industrie. 

\begin{itemize}
  \item \ulink{https://csrc.nist.gov/}{NIST COMPUTER SECURITY RESOURCE CENTER}
  \item \ulink{https://www.nist.gov/itl/fips-general-information}{NIST INFORMATION TECHNOLOGY LABORATORY}
\end{itemize}

On notera en particulier les référentiels cryptographiques du NIST et ceux liées à la cyberdéfense en particulier avec le \UKword{CyberSecurity FrameWork}

\subsubsection{SEI - Université de Carnegie Mellon}

Le Software Engineering Institute (SEI) est un centre de recherche-développement financé par des fonds fédéraux et placé sous le parrainage du département de la Défense des États-Unis ; son fonctionnement incombe à Carnegie Mellon University. Le SEI travaille avec des organisations pour apporter des améliorations significatives à leurs capacités d’ingénierie logicielle en leur fournissant le leadership technique afin de faire progresser la pratique de l’ingénierie logicielle. Le CERT Division du SEI est l’entité qui fait autorité et cherche à améliorer la sécurité et la résilience des systèmes et réseaux en particulier dans le domaine du logiciel(\href{https://www.sei.cmu.edu/research-capabilities/cybersecurity/}{Carnegie Mellon University - Cybersecurity research}).

\subsubsection{l'ISO : International Organization for Standardization }


L’ISO est une Organisation Internationale participant à l’élaboration de Standards. En ce sens la conformité à une norme a l’avantage d’être reconnue internationalement.
 
Les normes de la famille ISO 27000 permettent d’organiser et structurer la démarche de la gestion de la sécurité des systèmes d’information, une grande famille de normes avec des  positionnement sur l'ensemble du spectre  de la sécurité des systèmes d'information : 

\begin{itemize}
  \item ISO 27001 décrit les processus permettant le management de la sécurité de l’information (SMSI);
  \item ISO 27002 présente un catalogue de bonnes pratiques de sécurité;
  \item ISO 27003 décrit les différentes phases initiales à accomplir afin d’aboutir à un système de Management tel que décrit dans la norme ISO 27001;
  \item ISO 27004 permet de définir les contrôles de fonctionnement du SMSI;
  \item ISO 27005 décrit les processus de la gestion des risques;
  \item ISO 27006 décrit les exigences relatives aux organismes qui auditent et certifient les SMSI des sociétés.
\end{itemize}

Nous aborderons dans le chapitre sur les politiques de sécurité, l'usage de ce cadre normatif dans la gouvernance globale de la cybersécurité au sein de l'entreprise.

%TODO : Pourquoi se faire certifier ISO27K, les points durs.

\subsection{Europe}

Au niveau européen, le règlement (CE) 460/2004 du Parlement européen et du Conseil du 10 mars 2004 à institué l'Agence européenne chargée de la sécurité des réseaux et de l'information  Agence européenne chargée de la sécurité des réseaux et de l'information \ulink{https://www.enisa.europa.eu}{ENISA}. Sont rôle est de : 

%BEGIN wikpedia
% TODO
\begin{itemize}
  \item Conseiller et assister la Commission et les États membres en matière de sécurité de l'information et les aider, en concertation avec le secteur, à faire face aux problèmes de sécurité matérielle et logicielle.
  \item Recueillir et analyser les données relatives aux incidents liés à la sécurité en Europe et aux risques émergents.
  \item Promouvoir des méthodes d'évaluation et de gestion des risques afin d'améliorer notre capacité de faire face aux menaces pesant sur la sécurité de l'information.
  \item Favoriser l'échange de bonnes pratiques en matière de sensibilisation et de coopération avec les différents acteurs du domaine de la sécurité de l'information, notamment en créant des partenariats entre le secteur public et le secteur privé avec des entreprises spécialisées.
  \item Suivre l'élaboration des normes pour les produits et services en matière de sécurité des réseaux et de l'information.
\end{itemize}
%END Wikipeida

\subsection{France}

En France, la Cybersécurité est pilotée par un organisme dépendant des services du 1er Ministre,  l'Agence National des Systèmes d'information (\uac{aANSSI}).
L'\uac{aANSSI} possède plusieurs rôles de fait. C'est un \g{régulateur} c'est à dire qu'elle définit des cadres réglementaires pour les entreprises mais c'est aussi une agence qui édicte des préconisations et des guides.

 \ulink{https://www.ssi.gouv.fr/agence/cybersécurité/ssi-en-france/}{Le site de l'agence} est riche en informations et guides sur la cybersécurité.

Dépendant aussi de l'état, la \ulink{https://www.cnil.fr/}{CNIL} (Commission National Informatique et Liberté) est une autorité dont la mission est de protéger le citoyen. Avec l'avènement du règlement de protection des données personnelles, la \uac{aCNIL} a vu son pouvoir étendu.  

Il faut aussi citer l'AFNOR \uac{aAFNOR}, qui relaie en France la normalisation internationale dont l'ISO au delà de ses actions de normalisation purement françaises.

\section{Quelques associations et groupements professionnels} 

A titre d'information, vous trouverez en France aussi quelques clubs et associations historiques de la sécurité de systèmes d'information qui offrent à leurs adhérents des lieux d'échanges très intéressants et publient régulièrement  :



\begin{itemize}

  \item\head{Observatoire de la Sécurité des Systèmes d'Information et des Réseaux} (Technique)
\ulink{https://www.ossir.org/}{OSSIR}, association plutôt technique, qui propose de nombreux échanges sur la cybersécurité et existant depuis les années 90.
  \item\head{Club de la sécurité de l'information Français} (Gouvernance)
\ulink{https://clusif.fr}{CLUSIF}, association qui propose de nombreux échanges sur la cybersécurité.

  \item\head{Club CyberEdu} (Education)
\ulink{https://www.cyberedu.fr}{CyberEdu}, issu des travaux sur la formation des enseignants en cybersécurité de l'ANSSI, l'association regroupe les écoles et les utilisateurs des travaux de CyberEdu.

  \item\head{Club HexaTrust} (Editeurs de produits et services de cybersécurité Français)
\ulink{https://www.hexatrust.com/le-club/}{HexaTrust}, regroupe les éditeurs et fournisseurs de services français en cybersécurité.

  \item \head{Club des Experts de la sécurité de l’Information et du Numérique}. (Club de RSSI)
le \ulink{https://www.cesin.fr}{CESIN} est une association regroupant les RSSI d'entreprises, l'adhésion à cette association nécessite un parrainage et vous devez être RSSI.
\end{itemize}



%%==============================%
% MODULE 1 - Section 1 - Bulle 1 - Intro et Concepts
%==============================%
% La réalité de la menace
%-----------------------------------------------------%
\section{Immersion}

\subsection{Quelques principes}											
%==============================%

% SLIDE 1 : Introduction générale
%%-----------------------------------------------------%
\begin{frame}
	\frametitle{Immersion}
	\framesubtitle{La cyber c'est quoi}
\begin{itemize}
\item Quelques définitions pour appréhender les enjeux
\item Un peu d'histoire pour ne pas penser que le sujet est récent
\item Les biens les plus sensibles sont les biens les plus convoités
\item Pourquoi des attaques par l'exemple
\item Qui nous attaques, qui sont-ils et comment sont-ils organisés
\end{itemize}

\audionote{Commençons pas une petite définition bien plus large que son application à l'informatique domaine qui va concentrer notre attention. La sécurité c'est trois choses, c'est se sentir en sécurité c'est oeuvrer en confiance, et pour cela il faut que l'environnement soit de confiance mais ce n'est généralement pas le suffisant, il faut donc se protéger contre les menaces connues, s'avoir, pouvoir se défendre pour faire face à l'imprévu d'une situation ou d'une attaque, et être suffisamment solide pour rester debout face à l'agression, on appelle cela la résilience.}

\end{frame}

\subsubsection{définitions}		

%-----------------------------------------------------%

% SLIDE 2 : definition Sécurité
%-----------------------------------------------------%
\begin{frame}[c]{}
	\frametitle{Sécurité \& co}
	\framesubtitle{Test}
	\centering \LARGE SECURITE = \pause PROTECTION \pause+ DEFENSE \pause+ RESILIENCE
	% ajouter les boxeurs
\end{frame}
%-----------------------------------------------------%

%-----------------------------------------------------%

% SLIDE 3 : definition Cyber
%-----------------------------------------------------%
\begin{frame}[c]{}
	\frametitle{Sécurité \& co}
	\framesubtitle{Test}
	\centering \Large \tcb{CYBER}SECURITE = \tcb{CYBER}PROTECTION + \tcb{CYBER}DEFENSE + \tcb{CYBER}RESILIENCE
	% ajouter les boxeurs
\end{frame}
\subsection{Immersion}	

% Sources interessantes
% Panorama de la menace
%\ulink{https://www.cert.ssi.gouv.fr/uploads/20220309_NP_WHITE_ANSSI_panorama-menace-ANSSI.pdf}{Panorama ANSSI menace 2021}
%\ulink{https://www.cert.ssi.gouv.fr/uploads/CERTFR-2023-CTI-001.pdf}{Panorama ANSSI menace 2022}

%$END_UBODY
}%---------------------------------------------------------------------


\newcommand {\ubodyquizz}
{%---------------------------------------------------------------------

}%---------------------------------------------------------------------


%************************************************************
% Chargement  du MODELE
%************************************************************

\umainload



